\documentclass[10pt, letterpaper]{article}

\usepackage[margin=0.7in, top=0.6in, bottom=0.6in]{geometry}

% rl-at-scale shared preamble
% Packages and custom environments used across all documents

% --- Packages ---
\usepackage{amsmath, amssymb}
\usepackage{enumitem}
\usepackage{titlesec}
\usepackage[hidelinks]{hyperref}
\usepackage{fontenc}
\usepackage{lmodern}
\usepackage{booktabs}
\usepackage{xcolor}
\usepackage{tcolorbox}
\usepackage{listings}
\usepackage{fancyvrb}

% --- Study Guide boxes ---

% BLUF box - the one thing to walk away knowing
\newtcolorbox{bluf}{
  colback=black!5, colframe=black!70, boxrule=0.8pt,
  left=4pt, right=4pt, top=3pt, bottom=3pt,
  fonttitle=\bfseries\small, title=The Key Idea
}

% "How to say it" box
\newtcolorbox{sayit}[1][]{
  colback=green!4, colframe=green!50!black, boxrule=0.5pt,
  left=4pt, right=4pt, top=2pt, bottom=2pt,
  fonttitle=\bfseries\small, title=How to say it: #1
}

% "If they probe deeper" box
\newtcolorbox{deeper}{
  colback=orange!5, colframe=orange!60, boxrule=0.5pt,
  left=4pt, right=4pt, top=2pt, bottom=2pt,
  fonttitle=\bfseries\small, title=If they probe deeper
}

% "Why it matters for Anthropic" box
\newtcolorbox{whyanthro}{
  colback=blue!4, colframe=blue!40, boxrule=0.5pt,
  left=4pt, right=4pt, top=2pt, bottom=2pt,
  fonttitle=\bfseries\small, title=Why this matters for Anthropic
}

% --- Workbook boxes ---

% Exercise box
\newtcolorbox{exercise}[1][]{
  colback=yellow!5, colframe=yellow!60!black, boxrule=0.6pt,
  left=5pt, right=5pt, top=3pt, bottom=3pt,
  fonttitle=\bfseries\small, title=Exercise: #1
}

% Think-first box (conceptual, no computer needed)
\newtcolorbox{thinkfirst}[1][]{
  colback=purple!5, colframe=purple!50, boxrule=0.5pt,
  left=5pt, right=5pt, top=3pt, bottom=3pt,
  fonttitle=\bfseries\small, title={Think First (no computer): #1}
}

% Portfolio project box
\newtcolorbox{portfolio}[1][]{
  colback=green!5, colframe=green!50!black, boxrule=0.6pt,
  left=5pt, right=5pt, top=3pt, bottom=3pt,
  fonttitle=\bfseries\small, title={Portfolio Project: #1}
}

% Why this matters box
\newtcolorbox{whythis}{
  colback=blue!4, colframe=blue!40, boxrule=0.5pt,
  left=5pt, right=5pt, top=3pt, bottom=3pt,
  fonttitle=\bfseries\small, title=Why This Matters at Frontier Scale
}

% Research taste box
\newtcolorbox{taste}{
  colback=gray!6, colframe=gray!60, boxrule=0.5pt,
  left=5pt, right=5pt, top=3pt, bottom=3pt,
  fonttitle=\bfseries\small, title=Research Taste Check
}


\pagestyle{plain}
\setlength{\parindent}{0pt}
\setlength{\parskip}{4pt}
\titleformat{\section}{\Large\bfseries\scshape}{}{0em}{}[\titlerule]
\titlespacing*{\section}{0pt}{12pt}{6pt}
\titleformat{\subsection}{\large\bfseries}{}{0em}{}
\titlespacing*{\subsection}{0pt}{8pt}{4pt}
\titleformat{\subsubsection}{\normalsize\bfseries\itshape}{}{0em}{}
\titlespacing*{\subsubsection}{0pt}{6pt}{3pt}
\setlist[itemize]{nosep, leftmargin=1.4em, label=\textbullet, topsep=2pt}
\setlist[enumerate]{nosep, leftmargin=1.8em, topsep=2pt}

\lstset{
  language=Python,
  basicstyle=\small\ttfamily,
  keywordstyle=\color{blue!70!black},
  commentstyle=\color{green!50!black},
  stringstyle=\color{red!60!black},
  breaklines=true,
  frame=single,
  framesep=4pt,
  xleftmargin=4pt,
  xrightmargin=4pt,
  aboveskip=6pt,
  belowskip=6pt,
}

\begin{document}

\begin{center}
  {\LARGE \textbf{RL \& Performance Engineering Workbook}}\\[4pt]
  {\large Intelligence per Watt: From Kernels to RLVR Pipelines}\\[6pt]
  {\small v2.0 --- \today}
\end{center}

\bigskip

\textbf{Governing metric.} Every chapter in this workbook is organized around a single number: \textit{intelligence per Watt} --- the ratio of model capability to the energy required to produce it. Maximizing that ratio requires compressing both sides simultaneously: fewer FLOPs per forward pass (hardware and kernel chapters) and higher quality per FLOP (RL algorithm chapters).

\textbf{Three-tier structure.}
\begin{itemize}
  \item \textbf{Tier~1 --- PE Foundations (Chapters~1--5):} Hardware benchmarking (AMD, NVIDIA, TPU), inference efficiency, and distributed training. These chapters optimize FLOPs/Watt.
  \item \textbf{Tier~2 --- RL Depth (Chapters~6--7):} Hypothesis-driven policy gradient study and a hands-on community OSS framework. These chapters optimize Quality/FLOP.
  \item \textbf{Tier~3 --- RL Systems (Chapters~8--9):} Pipeline verification, reward engineering, and scaling laws. These chapters supply the systems foundation that makes Tier~2 work at scale.
\end{itemize}

\textbf{Box types.}
\begin{itemize}
  \item \textbf{Purple (Think-First, no computer):} Conceptual foundations. Reason through the problem before touching code.
  \item \textbf{Yellow (Exercise):} Hands-on implementation with explicit deliverables.
  \item \textbf{Green (Portfolio):} Publishable projects with external verification criteria.
\end{itemize}

\textbf{Recommended priority ordering.}

\begin{table}[h]
\centering
\small
\begin{tabular}{@{}llll@{}}
\toprule
\textbf{Priority} & \textbf{Chapter} & \textbf{Rationale} & \textbf{Compute} \\
\midrule
1a & Ch1 (AMD MI300X)      & ``Derisk Infrastructure Decision''; chip performance is still unknown & \textasciitilde\$300/mo \\
1b & Ch6 (Policy Gradient) & \$0 parallel track; builds RL depth immediately           & \$0 \\
2  & Ch7 (TinyRL)          & Community artifact; OSS + leaderboard signal             & \textasciitilde\$100/mo \\
3  & Ch4 (RL Inference)    & \#1 RLVR bottleneck in production                        & \textasciitilde\$100/mo \\
4  & Ch2 or Ch3 (NVIDIA/TPU) & Current stack depth; covers GPU and TPU Kernel Eng     & \textasciitilde\$250/mo \\
5  & Ch5 (Distributed)     & Cluster efficiency at scale                               & \textasciitilde\$200/mo \\
6  & Ch8--9 (RL Systems)   & Systems background; read after core chapters              & \textasciitilde\$100/mo \\
\bottomrule
\end{tabular}
\end{table}

\newpage

\section{Benchmarking AMD MI300X for RLVR Training Workloads}

\smallskip\noindent\textit{Study Guide companion: Study Guide Section 8 (Research Proposals) and Section 12 (Hardware Deep Dive) cover the interview-ready version of this material.}
\smallskip

\begin{whythis}
Frontier labs currently run primarily on NVIDIA hardware. But RLVR training runs last \textit{weeks}, making every hardware decision a multi-billion-dollar call. AMD's MI300X offers 5.3~TB/s of HBM3 bandwidth (vs.\ H100's 3.35~TB/s) and 192~GB of on-chip memory (vs.\ 80~GB) --- advantages that map directly onto the memory-bandwidth-bound operations that dominate RLVR wall time. The engineering question is whether AMD's software ecosystem is mature enough for the hardware advantage to survive contact with real workloads. This chapter builds the evidence base. Simon Boehm did this for NVIDIA; nobody has done it rigorously for AMD RLVR workloads. That gap is wide open.
\end{whythis}

\subsection{Prerequisites and Hardware}

\textbf{Reading:} HIP Programming Guide (\texttt{rocm.docs.amd.com}); AMD CDNA3 Architecture Whitepaper (covers XCD topology, LDS, MFMA); \texttt{composable\_kernel} library documentation (AMD's equivalent of CUTLASS); \texttt{rocprof} and \texttt{omniperf} profiling guides.

\textbf{Hardware:} Rent MI300X on Crusoe or TensorWave at $\sim$\$1.70/hr. Budget $\sim$\$300/month for this chapter. All findings extend to the MI355X --- same CDNA architecture family, higher clocks.

\textbf{Setup:} ROCm 6.x stack, PyTorch with ROCm backend, \texttt{hipcc} compiler, \texttt{rocblas} and \texttt{hipblaslt} for BLAS baselines.

% -----------------------------------------------------------------------
\subsection{Part A: Conceptual Foundation}
% -----------------------------------------------------------------------

\begin{thinkfirst}{AMD vs.\ NVIDIA Architecture: What Changes About Your Mental Model?}
Before writing a single line of HIP, reason through the architectural differences. This shapes every tiling and scheduling decision.

\textbf{Execution width:} NVIDIA warps are 32 threads wide. AMD wavefronts are 64 threads wide. If you tile a GEMM with a 32-wide thread block strategy from a CUDA tutorial, you are leaving half the hardware idle on MI300X. How must your tile dimensions change? What happens to register pressure when you double the execution width?

\textbf{Local Data Share vs.\ shared memory:} MI300X LDS is 64~KB per compute unit; NVIDIA's shared memory is 48~KB (configurable up to 99~KB on H100 with \texttt{cudaFuncSetAttribute}). The MI300X has more on-chip scratch by default. For a double-buffered matmul tile that holds two BF16 sub-matrices of size $T \times T$: what is the largest $T$ you can hold in LDS before spilling? How does this compare to the H100 default?

\textbf{MFMA vs.\ Tensor Cores:} AMD MFMA (Matrix Fused Multiply-Add) instructions operate on 16$\times$16 or 32$\times$32 tiles. NVIDIA's \texttt{wmma}/MMA operates on 16$\times$16 tiles. Different instruction shapes mean different register-file layouts for accumulator tiles. Before writing Kernel~2, sketch out the data layout your wavefront needs in registers to feed a 32$\times$32 MFMA correctly.

\textbf{GCD topology:} MI300X has 8 Graphics Compute Dies (XCDs), each with its own LDS. Inter-XCD communication goes through the Infinity Fabric, not shared L1. A kernel that touches memory resident on a different XCD pays a higher penalty than expected. How would you detect this with \texttt{omniperf}? What workload patterns expose it?

\textbf{HBM3 bandwidth:} MI300X delivers 5.3~TB/s; H100 delivers 3.35~TB/s --- a 1.58$\times$ advantage. This advantage only materializes for memory-bandwidth-bound operations. Write down: which operations in an RLVR training loop do you expect to be memory-bandwidth-bound, and which compute-bound? Predict before you measure.
\end{thinkfirst}

\begin{thinkfirst}{Roofline Analysis for MI300X}
The roofline model places operations on a two-regime map: memory-bound (performance $\propto$ bandwidth) vs.\ compute-bound (performance $\propto$ FLOP/s peak). The crossover is the \textit{arithmetic intensity ridge point}.

\textbf{MI300X specs:} Peak BF16 compute $\approx 383$~TFLOP/s; peak memory bandwidth $= 5.3$~TB/s. Derive the ridge point:
$$I^* = \frac{383 \times 10^{12} \text{ FLOP/s}}{5.3 \times 10^{12} \text{ B/s}} \approx 72 \text{ FLOP/byte}$$

\textbf{Compare to H100:} Peak BF16 $\approx 989$~TFLOP/s (with sparsity; $\sim 495$ dense); bandwidth $= 3.35$~TB/s. H100 ridge point $\approx 148$~FLOP/byte (dense). The MI300X's ridge is lower --- it becomes compute-bound at \textit{lower} arithmetic intensity than H100. Counterintuitive result: for operations with intermediate arithmetic intensity ($72$--$148$~FLOP/byte), MI300X may deliver better absolute throughput even at lower peak FLOP/s.

\textbf{BF16 matmul arithmetic intensity:} An $N \times N$ GEMM ($C = AB$, $A,B,C$ in BF16) performs $2N^3$ FLOPs and reads/writes $3 \times 2N^2$ bytes. Arithmetic intensity $= \frac{2N^3}{6N^2} = N/3$. At what $N$ is MI300X compute-bound? ($N/3 > 72 \Rightarrow N > 216$.) For a 7B model, typical sequence batch matmuls have $N \gg 216$ --- firmly compute-bound during training.

\textbf{Autoregressive generation:} During decode, batch size is often 1--4 tokens. The weight matrix is $O(D^2)$ bytes; FLOPs are $O(D^2)$ per token but memory reads are also $O(D^2)$ --- arithmetic intensity $\approx 1$--$2$~FLOP/byte. Far below the ridge: \textit{pure memory-bandwidth bottleneck}. MI300X's 1.58$\times$ bandwidth advantage should give a near-linear generation speedup. Verify this in Part~B.

\textbf{Prediction log:} Before running a single benchmark, fill in this table for 5 RLVR operations: (operation, expected intensity, expected bound, expected AMD advantage). You will check this against measured results in Part~C.
\end{thinkfirst}

\begin{thinkfirst}{Why Generation Throughput Dominates RLVR Wall Time}
RLVR training iterates: generate rollouts with the current policy, score them with a verifier, compute policy gradient, update weights. Repeat for weeks.

\textbf{Amdahl's Law applied:} Suppose generation consumes 60\% of wall time and the training step (forward + backward + optimizer) consumes 40\%. If you optimize generation by 2$\times$:
$$\text{Speedup} = \frac{1}{0.40 + 0.60/2} = \frac{1}{0.70} \approx 1.43\times$$
If instead you optimize the matmul inside the training step by 2$\times$, and matmul is 50\% of that 40\% step (20\% of total wall time):
$$\text{Speedup} = \frac{1}{0.80 + 0.20/2} = \frac{1}{0.90} \approx 1.11\times$$

A 2$\times$ generation improvement beats a 2$\times$ matmul improvement by 29 percentage points of pipeline speedup. Over a 3-week run, the 1.43$\times$ speedup saves 8.5 days vs.\ 2.3 days.

\textbf{Why MI300X's bandwidth advantage is structural here:} Generation is memory-bandwidth-bound by the KV~cache and weight loading. No algorithm change fixes this --- the roofline ceiling is bandwidth. MI300X's 1.58$\times$ bandwidth advantage translates almost directly to 1.58$\times$ generation throughput. This is the hardware bet. Test whether the software stack delivers it.

\textbf{Dario's ``laminar flow'' principle:} Numerical stability during RL training matters. Unlike supervised learning, a gradient explosion mid-run wastes weeks of compute, not hours. Before optimizing for throughput, understand what precision regime keeps gradients stable. This is why Kernel~3 and Part~B Op~4 focus on FP8/BF16 mixed precision with gradient norm monitoring --- not just raw TFLOP/s.

\textbf{Write your answer before proceeding:} Which single operation should you optimize first on MI300X to maximize RLVR training throughput? Why? If your answer is not ``autoregressive generation,'' re-read this box.
\end{thinkfirst}

% -----------------------------------------------------------------------
\subsection{Part A: HIP Matmul Progression}
% -----------------------------------------------------------------------

\begin{exercise}{HIP Matmul: Naive to Near-rocBLAS in 4 Kernels (8--12 hours)}
This progression builds kernel credibility and MI300X proficiency. The pedagogy is in the profiling and diagnosis loop, not the kernel count. After each kernel: measure, profile with \texttt{omniperf}, form a hypothesis about the bottleneck, then write the next kernel to address exactly that bottleneck. Do not move on until you can articulate why each kernel is faster than the previous one.

\textbf{Kernel~0: rocBLAS Baseline.}
Call \texttt{rocblas\_sgemm} (or \texttt{hipblaslt} for BF16) for matrix sizes $N \in \{512, 1024, 2048, 4096, 8192\}$. Record TFLOP/s as your ceiling. Record bandwidth utilization. Compute achieved arithmetic intensity and verify it matches the roofline prediction from the Think First box. This number is your denominator for all subsequent kernels. If you cannot reach 80\% of theoretical peak with rocBLAS, something is wrong with your setup --- fix it before proceeding.

\textbf{Record:} TFLOP/s at each $N$, \% of BF16 theoretical peak (383 TFLOP/s), HBM bandwidth utilization.

\medskip
\textbf{Kernel~1: Naive HIP $\to$ LDS Tiling (Boehm stages 1--3 collapsed).}
Start with the simplest possible kernel: one thread computes one output element $C[i,j]$ by iterating over the $k$ dimension. Measure TFLOP/s --- expect $<5\%$ of rocBLAS. Then add memory coalescing (ensure threads in a wavefront access consecutive memory addresses) and LDS blocking in one progression. Tile size: start with $64 \times 64$; AMD-specific --- LDS is 64~KB, so a $64 \times 64$ BF16 tile is $64^2 \times 2 = 8$~KB; you can hold 8 such tiles. Double-buffer: load tile $k$ into LDS while computing tile $k-1$.

\textbf{AMD-specific pitfall:} Wavefronts are 64-wide. If your thread block is $32 \times 32 = 1024$ threads, you have 16 wavefronts. Ensure your LDS access pattern avoids bank conflicts for 64-wide accesses --- the bank structure is 32 banks, but the stride pattern differs from CUDA. Profile with \texttt{omniperf --dispatch-ids} to catch LDS bank conflicts directly.

\textbf{Record:} TFLOP/s, \% of rocBLAS, LDS bandwidth utilization, bank conflict rate.

\medskip
\textbf{Kernel~2: MFMA + Wavefront Tiling.}
This is the AMD-differentiating stage. Replace scalar FMA with MFMA instructions. Use \texttt{\_\_builtin\_amdgcn\_mfma\_f32\_32x32x8\_bf16} for 32$\times$32$\times$8 BF16 MFMA. Each MFMA instruction consumes registers from the full 64-wide wavefront --- you must lay out your register file so that each lane of the wavefront holds the correct $A$ and $B$ fragments. Tile dimensions: outer tile 128$\times$128 (computed by the thread block), inner tile 32$\times$32 (computed by one MFMA per wavefront). Load $A$ fragments (32$\times$8 BF16) and $B$ fragments (8$\times$32 BF16) from LDS into registers. Accumulate into 32$\times$32 FP32 accumulators.

\textbf{Profiling checkpoint:} Run \texttt{omniperf} and check VALU (vector ALU) utilization vs.\ MFMA utilization. If VALU is high and MFMA is low, your kernel is spending time on address computation, not matrix math. Reduce address arithmetic in the hot loop.

\textbf{Record:} TFLOP/s, \% of rocBLAS, VALU utilization, MFMA utilization, wavefront occupancy.

\medskip
\textbf{Kernel~3: FP8/BF16 Mixed Precision + Vectorized 128-bit Loads.}
RL training stability requires careful precision management. Use BF16 for weight storage and FP8 for activation matmul (where supported by ROCm), accumulating into FP32. This is the regime Dario's ``laminar flow'' principle applies to: monitor gradient norm variance across training steps. If norms spike, FP8 accumulation is the suspect --- fall back to BF16.

Replace scalar LDS loads with 128-bit vectorized loads (\texttt{float4} for 4$\times$FP32 or \texttt{short8} for 8$\times$BF16). This reduces the number of load instructions per LDS access and improves memory-level parallelism. AMD's memory subsystem prefers 128-bit aligned accesses --- verify alignment or profile for misalignment penalties.

\textbf{Target:} $>50\%$ of rocBLAS TFLOP/s by this kernel. If you fall short, profile with \texttt{omniperf} and identify the binding constraint before declaring done. Common culprits: register spilling (check SGPRs/VGPRs in \texttt{rocprof} output), LDS bank conflicts, insufficient wavefront occupancy.

\textbf{Record:} TFLOP/s, \% of rocBLAS, gradient norm variance (run 100 steps of a toy training loop with these matmuls), effective HBM bandwidth.

\medskip
\textbf{Diagnostic discipline:} After each kernel, write one paragraph: what was the bottleneck, how did you diagnose it, what did you change, what was the result. These paragraphs become the core of your blog post's technical narrative.

\textit{This establishes kernel credibility. The actual value for the portfolio is in Part~B.}
\end{exercise}

% -----------------------------------------------------------------------
\subsection{Part B: RLVR-Relevant Operations Benchmark}
% -----------------------------------------------------------------------

\begin{exercise}{RLVR Operations Benchmark on MI300X (15--20 hours)}
Part~A proved you can write HIP kernels. Part~B answers the question that matters: \textit{does MI300X's hardware advantage survive real RLVR workloads?} For each operation below: (1) state why it appears in RLVR training, (2) run the benchmark, (3) compare to H100 published numbers or your own H100 measurements if available.

\textbf{Benchmark harness:} Write a single Python script that runs all 7 operations, logs results to a JSON file, and prints a summary table. Use \texttt{torch.cuda.synchronize()} (which maps to \texttt{torch.hip.synchronize()} on ROCm) for accurate timing. Warm up 10 iterations; measure 100 iterations; report median and p95 latency.

\medskip
\textbf{Operation 1: Autoregressive Generation Throughput.} \textit{The RLVR bottleneck.}

RLVR spends 50--70\% of wall time generating rollouts. Each rollout is autoregressive: one token at a time, loading the full weight matrix and KV~cache at every step. This is memory-bandwidth-bound.

Load a 7B-parameter model in BF16 (\textasciitilde{}14~GB). Generate 512 tokens with batch sizes $\{1, 4, 8, 16, 32\}$. Measure tokens/second end-to-end (including KV~cache updates). MI300X's 5.3~TB/s should give $\approx 1.58\times$ vs.\ H100's 3.35~TB/s, since the bottleneck is weight loading bandwidth. Verify. If you observe less than $1.3\times$ speedup, profile KV~cache overhead separately --- the gap is likely in attention kernel efficiency or RCCL initialization overhead, not raw HBM bandwidth.

\textbf{Measure:} Tokens/second at each batch size. Roofline utilization (achieved HBM bandwidth / 5.3~TB/s). Ratio vs.\ H100 baseline.

\medskip
\textbf{Operation 2: Flash Attention Kernel.} \textit{The training compute kernel.}

Flash Attention fuses the attention computation to avoid materializing the full $N \times N$ attention matrix. On AMD, the reference implementation is in \texttt{composable\_kernel} (AMD's equivalent of CUTLASS). Alternatively, use \texttt{flash-attention} with the ROCm backend if available for your ROCm version.

Run for sequence lengths $\{512, 1024, 2048, 4096, 8192\}$, head dim 128, 32 heads, batch size 4. Compare to FlashAttention-2 on H100 (TFLOP/s numbers from Tri Dao's paper are public). Measure both forward and backward pass. If \texttt{composable\_kernel} FA underperforms by $>30\%$ vs.\ H100, identify whether the gap is in the MFMA utilization or the softmax numerics --- these have different fixes.

\textbf{Measure:} TFLOP/s (fwd), TFLOP/s (bwd), \% of theoretical MI300X peak, ratio vs.\ FlashAttention-2 H100.

\medskip
\textbf{Operation 3: All-Reduce (RCCL).} \textit{The distributed training bottleneck.}

Multi-GPU RLVR training synchronizes gradients via all-reduce after each backward pass. On NVIDIA, NCCL is battle-hardened. AMD uses RCCL (ROCm Collective Communication Library).

Using 4 MI300X GPUs (rent a 4-GPU node), run RCCL all-reduce at message sizes $\{1\text{MB}, 10\text{MB}, 100\text{MB}, 1\text{GB}, 10\text{GB}\}$. Plot bandwidth utilization (achieved GB/s / theoretical NVLink or Infinity Fabric bandwidth). Compare to published NCCL numbers at equivalent message sizes. Key question: does RCCL's bandwidth utilization fall off at small message sizes (the latency-dominated regime) --- and what is the crossover message size where it becomes bandwidth-efficient? For a 7B model, gradient tensors are $\sim$14~GB total --- is this in the efficient regime?

\textbf{Measure:} All-reduce bandwidth (GB/s) at each message size, \% of Infinity Fabric peak, ratio vs.\ NCCL H100.

\medskip
\textbf{Operation 4: Mixed-Precision Training Step.} \textit{Throughput + stability.}

A realistic training step: FP8 matmul for forward pass, BF16 for attention, FP32 accumulation for gradient updates. This matches the precision hierarchy used in production RLVR.

Construct a minimal training loop: 2-layer transformer, 7B-parameter size, batch size 4, sequence length 512. Run 100 forward-backward-optimizer steps. Measure: (a) TFLOP/s of the matmul-heavy forward pass, (b) gradient norm at each step --- flag if norm exceeds $10\times$ the initial norm (instability signal), (c) wall time per step.

This is numerically interesting because RL rewards are often sparse and bursty --- gradient norms spike around reward events. FP8 accumulation may be fine for supervised training but unstable for RL. Document what you find.

\textbf{Measure:} TFLOP/s, gradient norm variance, wall time/step, any instability events.

\medskip
\textbf{Operation 5: KV Cache Capacity.} \textit{The memory-capacity advantage.}

MI300X has 192~GB HBM3 vs.\ H100's 80~GB. For RLVR, long-context generation is increasingly important --- reasoning traces can be thousands of tokens. The question is at what context length MI300X's capacity advantage becomes the binding constraint.

For a 7B model in BF16: KV~cache per token per layer $= 2 \times 2 \times D_{\text{head}} \times H_{\text{heads}} = 2 \times 2 \times 128 \times 32 = 16,384$~bytes $= 16$~KB. For 32 layers: $512$~KB per token. At what context length does the KV~cache alone fill 80~GB? (Answer: $\approx 160K$ tokens.) At what length does it fill 192~GB? ($\approx 375K$ tokens.) Verify these numbers experimentally by running generation until OOM. Repeat for a 70B model.

\textbf{Measure:} Max context length before OOM for 7B and 70B models. Tokens/second as context length grows (bandwidth demand grows with KV~cache size). The point at which throughput degrades due to KV~cache bandwidth pressure.

\medskip
\textbf{Operation 6: MoE Memory Budget.} \textit{Dario's MoE bet.}

Frontier labs increasingly use MoE architectures. An 8-expert MoE (2 active per token) has $8\times$ the total parameter count but only $2\times$ the active FLOPs. The memory footprint for all experts must fit on-device.

Model: 8 experts, each a 7B MLP block (2-layer FFN, hidden dim 28,672). Expert weights: $8 \times 7B \times 2~\text{bytes} = 112$~GB for BF16. On H100 (80~GB): does not fit. On MI300X (192~GB): fits with 80~GB remaining for KV~cache, activations, and optimizer state. At what context length does the remaining 80~GB for KV~cache run out? (Use Op~5 calculation for 7B equivalent.) At what total model scale does MI300X's 192~GB advantage become the reason you pick AMD over NVIDIA?

\textbf{Measure:} Expert weight memory footprint. Remaining capacity for KV~cache. Max context length achievable with full MoE model loaded. Tokens/second for MoE decode vs.\ dense decode (expert routing overhead).

\medskip
\textbf{Operation 7: Decode-Path Attention (KV Cache Memory Movement).} \textit{Not matmul.}

Autoregressive decode is not a matmul problem --- it is a KV~cache memory movement problem. During decode step $t$, the model reads the full KV~cache (all $t$ previous tokens), computes attention scores, and writes one new KV pair. This is an irregular memory access pattern that MFMA instructions do not directly accelerate.

Use \texttt{composable\_kernel} decode attention kernels (or implement via \texttt{xformers} with ROCm backend). Profile with \texttt{omniperf}: what fraction of the decode step is (a) KV~cache reads, (b) attention score computation, (c) KV write. At sequence length 4096, what is the ratio of memory bandwidth consumed to FLOPs performed? Does this match the $\approx 1$~FLOP/byte prediction from the Think First roofline box?

\textbf{Measure:} KV~cache read bandwidth (GB/s), FLOPs/s for attention scoring, HBM bandwidth utilization as \% of 5.3~TB/s, ratio vs.\ H100 decode attention benchmark.
\end{exercise}

% -----------------------------------------------------------------------
\subsection{Part C: Intelligence-per-Watt Comparison}
% -----------------------------------------------------------------------

\begin{thinkfirst}{Predict Before You Measure: Intelligence-per-Watt}
Before running Part~C benchmarks, write down your predictions. Return to this box and grade yourself after.

MI300X TDP: 750~W. H100 SXM TDP: 700~W. The MI300X consumes 7\% more power. For AMD to win on intelligence-per-Watt, it must deliver $>7\%$ more useful throughput on a given operation.

\textbf{Predict for each Part~B operation:} Will MI300X deliver better FLOP/s-per-Watt (or tokens/s-per-Watt) than H100? Reason from first principles using the roofline analysis and architecture differences you studied.

\begin{enumerate}
  \item Autoregressive generation: AMD wins/loses? Why? (Hint: bandwidth is proportional to power consumption approximately linearly for HBM.)
  \item Flash Attention forward: AMD wins/loses? What determines the efficiency ratio?
  \item All-reduce (RCCL vs.\ NCCL): AMD wins/loses? What is the primary source of software overhead?
  \item Mixed-precision training step: AMD wins/loses? Does FP8 throughput on MFMA vs.\ Tensor Cores change the picture?
  \item KV~cache capacity (192~GB vs.\ 80~GB): This is not a throughput question --- it is a capability question. At what context length does MI300X unlock workloads that are simply impossible on H100?
\end{enumerate}

\textbf{Write your predictions now.} Do not read Part~C exercise results first.
\end{thinkfirst}

\begin{exercise}{Intelligence-per-Watt Analysis (2--3 hours)}
Using the benchmark results from Parts~A and~B, build the decision-useful analysis that justifies or refutes AMD investment.

\textbf{Step 1: Compute useful FLOP/s per Watt.} For each Part~B operation, record: (achieved throughput on MI300X) / 750~W and (H100 benchmark throughput) / 700~W. Use tokens/s for generation; TFLOP/s for matmul and attention; GB/s for all-reduce. Build a table with columns: Operation, MI300X metric, H100 metric, MI300X/W, H100/W, Ratio.

\textbf{Step 2: Gap analysis.} For operations where MI300X loses on per-Watt efficiency, categorize the gap source:
\begin{itemize}
  \item \textit{Software immaturity} --- kernel not yet optimized (e.g., RCCL, FA not tuned): gap closable with engineering effort.
  \item \textit{Compiler gap} --- \texttt{hipcc} generating worse code than \texttt{nvcc} for a class of patterns: harder to close, may require upstream ROCm contributions.
  \item \textit{Architectural mismatch} --- the operation genuinely suits NVIDIA's architecture better (e.g., Tensor Core scheduling, warp-level primitives): structural disadvantage, cannot be engineered away cheaply.
\end{itemize}
For each gap, estimate engineering effort to close: $<1$~person-week, 1--4~person-weeks, $>1$~person-month, or ``architectural, cannot close.''

\textbf{Step 3: Write the memo.} One page, structured as follows:
\begin{enumerate}
  \item \textbf{Recommendation} (one sentence): Should Anthropic invest in MI300X for RLVR training, and under what conditions?
  \item \textbf{Hardware case} (2--3 bullets): The specific operations where AMD wins on intelligence-per-Watt and why the win is structural.
  \item \textbf{Risks} (2--3 bullets): The specific gaps, their sources, and engineering cost to close.
  \item \textbf{Decision trigger} (1 paragraph): What would need to be true for the answer to change? E.g., ``If RCCL reaches NCCL parity at the 10~GB message size, the all-reduce gap closes and the hardware case strengthens by X\%.'' This is the section a senior RL engineer can act on.
\end{enumerate}

\textbf{Verification:} Your analysis should identify at least one operation where AMD wins on intelligence-per-Watt AND at least one where it loses. If your analysis shows AMD wins everywhere, you have not stress-tested the software gaps. If it shows AMD loses everywhere, you have not accounted for the bandwidth and memory-capacity advantages.
\end{exercise}

% -----------------------------------------------------------------------
\subsection{Portfolio Project}
% -----------------------------------------------------------------------

\begin{portfolio}{Evaluating AMD MI300X for Large-Scale RLVR Training}
\textbf{Deliverable:} Public blog post + open-source repository. Estimated total time: 30--40 hours including Parts A--C.

\textbf{Repository structure:}
\begin{itemize}
  \item \texttt{kernels/} --- HIP kernels from Part~A, each in a separate file with a build script.
  \item \texttt{benchmarks/} --- Python benchmark harness for all 7 Part~B operations, runnable on any MI300X with ROCm 6.x.
  \item \texttt{results/} --- Raw JSON output from your benchmark runs, plus a Jupyter notebook that generates all plots.
  \item \texttt{analysis/} --- The intelligence-per-Watt memo and gap analysis from Part~C.
  \item \texttt{README.md} --- Reproduction instructions: exact ROCm version, model checkpoint, cloud provider, instance type, and total cost to reproduce.
\end{itemize}

\textbf{Blog post structure:} Part~A kernels are the pedagogical hook that draws readers who want to learn HIP. Parts~B and~C are the decision-useful substance that RL infrastructure teams can act on. The narrative thread: ``I set out to understand whether AMD's hardware advantage survives real RLVR workloads. Here is what I found, operation by operation.''

\textbf{The monkey:} Does MI300X's 1.58$\times$ bandwidth advantage actually translate to faster RLVR training end-to-end, or does software ecosystem immaturity (RCCL, composable\_kernel, compiler maturity) erase the hardware advantage? The benchmark harness answers this empirically. The gap analysis answers what would need to change for the answer to flip. That combination --- empirical data plus actionable gap analysis --- is what makes the project decision-useful rather than merely educational.

\textbf{Verification criteria:}
\begin{enumerate}
  \item The analysis identifies $\geq 1$ operation where AMD wins on intelligence-per-Watt.
  \item The analysis identifies $\geq 1$ operation where AMD loses, with a categorized gap source.
  \item The gap analysis includes a concrete engineering effort estimate ($<$1 week / 1--4 weeks / $>$1 month / architectural) for each gap.
  \item The benchmark harness is reproducible by a third party with access to a single MI300X instance at the stated ROCm version.
  \item Total cloud cost to reproduce is documented and is under \$200.
\end{enumerate}

\textbf{Why this matters:} Simon Boehm's CUDA matmul post demonstrated kernel authorship on NVIDIA. This project does the same for AMD, but at a higher level of strategic relevance: it produces evidence that informs a hardware purchasing decision for any team evaluating AMD for RL workloads. A senior RL engineer reading your blog can directly use the gap analysis and cost estimates. That is the difference between a tutorial and a contribution.
\end{portfolio}

\newpage

\section{NVIDIA CUDA/Blackwell Kernel Optimization}

\smallskip\noindent\textit{Study Guide companion: Study Guide Section 12 (Hardware Deep Dive) covers the interview-ready version of this material.}
\smallskip

\begin{whythis}
Anthropic runs on NVIDIA hardware, currently transitioning to Blackwell (B200) via Azure. The PE/RE GPU role explicitly lists CUDA, Triton, CUTLASS, Flash Attention, and tensor core optimization as required skills. Simon Boehm's SGEMM\_CUDA project---13 kernels from naive to 93.7\% of cuBLAS on an A6000---is the canonical public demonstration of kernel authorship at this level. This chapter replicates and extends that progression to Blackwell-specific features: TMEM, native FP4/FP6, and thread block clusters up to size~16. The governing metric throughout is \textit{intelligence per Watt}: not raw TFLOP/s, but useful inference throughput per joule, which is the actual constraint in large-scale RLVR training where generation is the wall-time bottleneck.
\end{whythis}

\subsection{Prerequisites and Hardware}

\textbf{Reading:} Boehm's SGEMM\_CUDA repository (\texttt{github.com/siboehm/SGEMM\_CUDA}) and accompanying blog post ``How to Optimize a CUDA Matmul Kernel for cuBLAS-like Performance''; CUDA Programming Guide chapters on shared memory, warp execution, and tensor cores; CUTLASS documentation; Tri Dao's Flash Attention papers (v1 and v2); the Blackwell Architecture Whitepaper (GTC 2024); Nsight Compute user guide.

\textbf{Companion reference:} The SGEMM\_CUDA repo is at \texttt{SGEMM\_CUDA/} on your machine. It implements kernels 0--13 with profiling scaffolding. Study it as source material for understanding design decisions, not as code to copy.

\textbf{Hardware:} H100 SXM at $\sim$\$2/hr (Kernels 0--7, Lambda Labs or CoreWeave); B200 at $\sim$\$2.25/hr (Kernels 8--10, available via Azure or specialized providers). Budget $\sim$\$250/month. Total estimated spend for the full chapter: \$80--\$120 if you are disciplined about not leaving instances idle.

\textbf{Setup:} CUDA 12.4+, cuBLAS, Nsight Compute (\texttt{ncu}), PyTorch for Triton baseline, \texttt{nvcc} with \texttt{-arch=sm\_90} (H100) or \texttt{-arch=sm\_100} (B200).

% -----------------------------------------------------------------------
\subsection{Part A: Conceptual Foundation}
% -----------------------------------------------------------------------

\begin{thinkfirst}{Hopper vs.\ Blackwell Architecture: What Changes About Your Optimization Strategy?}
Before writing a single kernel, reason through the architectural differences. Not every Hopper technique ports cleanly to Blackwell, and Blackwell introduces memory and synchronization primitives that require new mental models.

\textbf{Hopper (H100) specifications:} 80~GB HBM3 at 3.35~TB/s; 132 SMs each with 256~KB L1 cache (configurable as up to 228~KB shared memory); Tensor Core generation 4 supporting FP8/BF16/FP16/TF32/FP32; thread block cluster size up to 8; no TMEM.

\textbf{Blackwell (B200) specifications:} 192~GB HBM3e at 8~TB/s; TMEM (Tensor Memory)---a new on-chip memory space separate from shared memory, reserved for advanced kernel programming; native FP4 and FP6 data types with hardware Tensor Core support; thread block clusters up to size~16; NVLink 5 for multi-GPU. The 8~TB/s bandwidth figure is a 2.4$\times$ improvement over H100.

\textbf{What changes:}
\begin{itemize}
  \item \textbf{Roofline ridge point:} At BF16, B200 peak compute is $\sim 4500$~TFLOP/s (dense); bandwidth is $8$~TB/s. Ridge: $4500/8 = 562$~FLOP/byte. H100 ridge is $\sim 148$~FLOP/byte (dense BF16). Operations that were memory-bound on H100 can become compute-bound on B200---including medium-batch inference matmuls.
  \item \textbf{TMEM:} The accumulator registers in Blackwell Tensor Cores can now reside in TMEM rather than the register file. This reduces register pressure and enables larger tile sizes without spilling. TMEM requires explicit allocation and has a different programming interface than shared memory.
  \item \textbf{FP4/FP6:} At FP4, each element is 0.5 bytes. A 4096$\times$4096 weight matrix in FP4 is 8~MB vs.\ 32~MB in FP16. The bandwidth savings compound: both reads and writes shrink, and the roofline ceiling (compute) rises because Tensor Cores execute more FP4 ops per cycle than FP16 ops.
  \item \textbf{Cluster size 16:} Hopper introduced thread block clusters (CGA) at size~8. Blackwell doubles this to 16. A cluster of 16 blocks shares distributed shared memory (DSMEM) via the NVLink fabric between SMs, enabling tiles up to $16 \times \text{block\_shared\_mem}$ without explicit global memory round-trips.
\end{itemize}

\textbf{What stays the same:} The memory hierarchy reasoning---registers $>$ shared memory (or TMEM) $>$ L2 $>$ HBM---is unchanged. Tiling to exploit locality, double-buffering to hide latency, coalescing global memory accesses, avoiding bank conflicts in shared memory: all of these apply on Blackwell with the same logic.

\textbf{Write before proceeding:} For a batched matmul at inference time (batch size 64, $M=N=K=4096$, FP16 weights): is this compute-bound or memory-bound on H100? On B200? Does your answer change at batch size 1? At FP4 precision? Compute arithmetic intensity for each case before checking against the roofline.
\end{thinkfirst}

\begin{thinkfirst}{When Triton Beats Hand-Written CUDA}
Triton is a Python DSL that compiles to efficient PTX. It auto-tiles, auto-vectorizes, and manages shared memory layout. For a standard FP16 matmul, Triton matches 90\%+ of hand-written CUDA with roughly 10$\times$ less code. This raises the question: why write CUDA at all?

\textbf{Where Triton wins:}
\begin{itemize}
  \item \textbf{Iteration speed:} A Triton matmul is $\sim$50 lines. An equivalent hand-written CUDA kernel with warptiling and vectorized loads is $\sim$500 lines. If the goal is to benchmark a new attention variant or test a quantization scheme, Triton lets you iterate in hours instead of days.
  \item \textbf{Auto-tuning:} Triton's \texttt{@triton.autotune} decorator searches the tile-size space automatically. Hand-written CUDA requires you to build your own sweep or use a tool like \texttt{ncu} to guide tuning.
  \item \textbf{Compiler improvements:} Triton's compiler is rapidly improving. Kernels that were 20\% below cuBLAS in 2022 are now within 5\% in 2025 for many shapes.
\end{itemize}

\textbf{Where hand-written CUDA wins:}
\begin{itemize}
  \item \textbf{Custom memory layouts:} Triton's abstraction assumes a block-based memory model. If you need a non-standard shared memory layout (e.g., a transposed tile for a specific warp schedule), Triton cannot express it.
  \item \textbf{Warp-level scheduling:} Hand-written CUDA lets you interleave MMA instructions with load instructions at warp granularity, hiding memory latency through instruction-level parallelism. Triton schedules at block granularity.
  \item \textbf{TMEM and FP4:} As of early 2026, Triton does not expose TMEM or native FP4 Tensor Core instructions. Accessing Blackwell-specific features requires CUDA or PTX.
  \item \textbf{Flash Attention:} The Flash Attention kernel relies on careful management of the online softmax trick and recomputation scheduling that Triton can express (FlashAttention-2 has a Triton implementation), but the highest-performance versions use hand-written CUDA to control warp scheduling precisely.
\end{itemize}

\textbf{The practical rule:} Use Triton for prototyping and for operations that are not on the critical path. Use hand-written CUDA for the 3--5 kernels that dominate RLVR wall time: autoregressive attention, KV-cache update, and the generation-side matmul. The chapter exercises build both skills; the Triton comparison exercise (Part~D) makes the performance gap concrete.

\textbf{Write before proceeding:} For RLVR generation at batch size 8 with a 7B model: which single kernel do you expect to contribute most to wall time? Is Triton likely sufficient for production use of that kernel, or do you need hand-written CUDA? Justify with roofline arithmetic.
\end{thinkfirst}

\begin{thinkfirst}{Arithmetic Intensity Shift: FP16 $\to$ FP8 $\to$ FP4}
As precision decreases, bytes per element halve (FP16: 2 bytes; FP8: 1 byte; FP4: 0.5 bytes). Simultaneously, Tensor Cores at lower precision can execute more operations per cycle: Blackwell FP4 Tensor Core throughput is roughly $4\times$ higher than FP16 (hardware processes 4$\times$ as many elements per instruction). This combination shifts the roofline ridge point dramatically.

\textbf{B200 specifications by precision (approximate):}
\begin{itemize}
  \item FP16/BF16: $\sim 4500$~TFLOP/s dense; bandwidth 8~TB/s; ridge $\approx 562$~FLOP/byte
  \item FP8: $\sim 9000$~TFLOP/s; bandwidth 8~TB/s; ridge $\approx 1125$~FLOP/byte
  \item FP4: $\sim 18000$~TFLOP/s; bandwidth 8~TB/s; ridge $\approx 2250$~FLOP/byte
\end{itemize}

\textbf{Matmul arithmetic intensity:} An $N \times N$ GEMM with FP4 inputs performs $2N^3$ FLOPs and reads $3 \times 0.5N^2$ bytes (FP4 A, B, C). Intensity $= 2N^3 / (1.5N^2) = (4/3)N$.

\textbf{Work out these cases} before proceeding:

\begin{center}
\begin{tabular}{lrrrr}
\hline
\textbf{Precision} & $N=1024$ & $N=2048$ & $N=4096$ & $N=8192$ \\
\hline
FP16 intensity (FLOP/byte) & 341 & 683 & 1365 & 2731 \\
FP16 bound on B200? & compute & compute & compute & compute \\
FP4 intensity (FLOP/byte) & 1365 & 2731 & 5461 & 10923 \\
FP4 bound on B200? & compute & compute & compute & compute \\
\hline
\end{tabular}
\end{center}

\textit{Note: at $N=1024$, FP4 intensity (1365) exceeds the FP4 ridge (2250)? Check your arithmetic. The ridge for FP4 is $\sim 2250$~FLOP/byte, so $N=1024$ gives intensity 1365~FLOP/byte, which is \textbf{below} the ridge---memory-bound at FP4. At $N=2048$, intensity 2731 exceeds 2250---compute-bound. Fill in the table above by working through each case, then verify: is there any $(N, \text{precision})$ pair in the table that is memory-bound? If so, what does that tell you about FP4 for small-batch inference on B200?}

\textbf{RLVR implication:} During generation (decode phase), batch size is small and sequence length grows. The effective ``N'' for the weight matmul is the batch size $\times$ head dimension, which at batch size 1 is far below 2048. At batch size 1 with FP4 weights, you are \textit{memory-bound} on B200 despite FP4's computational advantage---the bandwidth savings (4$\times$ fewer bytes) are the primary benefit, not the compute throughput. This means the B200's 8~TB/s bandwidth matters more than its FP4 TFLOP/s for single-stream generation.
\end{thinkfirst}

% -----------------------------------------------------------------------
\subsection{Part B: The Boehm Progression on H100}
% -----------------------------------------------------------------------

\begin{exercise}{CUDA Matmul: Boehm Kernels 0--7 (15--20 hours)}
This is the core credentialing exercise. The SGEMM\_CUDA repo at \texttt{SGEMM\_CUDA/} implements all kernels; your job is to understand \textit{why} each optimization works, not to copy the code. After each kernel: measure TFLOP/s, compute \% of cuBLAS, profile with Nsight Compute (\texttt{ncu}), identify the specific bottleneck that the next kernel addresses. Do not advance until you can write a one-sentence causal explanation: ``Kernel $k+1$ is faster than Kernel $k$ because \underline{\phantom{xxx}}.''

Compile with \texttt{nvcc -O3 -arch=sm\_90 --use\_fast\_math}. Benchmark at $N \in \{1024, 2048, 4096\}$ with 100 warm-up iterations and 1000 timed iterations. Report median TFLOP/s.

\medskip
\textbf{Kernel 0: cuBLAS Baseline.}
Call \texttt{cublasSgemm} (FP32) for your three matrix sizes. Record TFLOP/s as your ceiling. This is the number you are chasing. Note: cuBLAS uses internal heuristics to select its kernel; for large $N$ on H100, it uses a warptiled HMMA kernel. Your kernel~7 target is 70\%+ of this number (Boehm reached 93.7\% on an A6000; H100's higher-performance cuBLAS raises the bar).

\textbf{Record:} TFLOP/s at each $N$; bandwidth utilization from Nsight.

\medskip
\textbf{Kernel 1: Naive (1 thread per output element).}
One thread computes \texttt{C[row][col]} by looping over $k$. Global memory is accessed non-coalescingly: thread $(i,j)$ reads row $i$ of $A$ (stride-1 in $k$) and column $j$ of $B$ (stride-$N$ in $k$). The column access strides over $N$ elements between consecutive $k$ values---each access is a separate cache line.

\textbf{Profile:} Use \texttt{ncu --metrics l1tex\_\_t\_bytes\_pipe\_lsu\_mem\_global\_op\_ld.sum,sm\_\_cycles\_active.avg} to measure global load throughput. Expect $<5\%$ of cuBLAS.

\textbf{Predict before measuring:} Given HBM3 bandwidth at 3.35~TB/s and the arithmetic intensity of this naive kernel (approximately 1~FLOP/byte due to no data reuse), what is the theoretical TFLOP/s ceiling? Does your measurement match?

\medskip
\textbf{Kernel 2: Global Memory Coalescing.}
The fix: transpose the thread-to-output-element mapping so that threads in a warp access consecutive columns of $C$, which maps to consecutive rows of $B$ in memory. Concretely: thread $(t_x, t_y)$ in a $32 \times 32$ block computes $C[t_y + \text{blockIdx.y} \cdot 32][t_x + \text{blockIdx.x} \cdot 32]$. Now 32 threads in a warp access 32 consecutive elements of a $B$ row---a single 128-byte cache line transaction for the full warp.

\textbf{Profile:} Check \texttt{l1tex\_\_t\_sectors\_pipe\_lsu\_mem\_global\_op\_ld.sum} before and after. Each warp's 32 accesses should now generate 1 sector request instead of 32.

\textbf{Expected gain:} 2--5$\times$ over Kernel~1. Explain: coalescing turns 32 memory transactions per warp into 1.

\medskip
\textbf{Kernel 3: Shared Memory Blocking.}
Load \texttt{BLOCKSIZE}$\times$\texttt{BLOCKSIZE} tiles of $A$ and $B$ into shared memory before computing. Each element is loaded once from HBM into shared memory; then \texttt{BLOCKSIZE} multiply-accumulate operations reuse it from shared memory. Arithmetic intensity improves from $O(1)$ to $O(\text{BLOCKSIZE})$.

Key implementation detail: both $A$ and $B$ tiles are loaded cooperatively by all threads in the block. After loading, \texttt{\_\_syncthreads()} ensures all data is visible before computation begins; a second \texttt{\_\_syncthreads()} at the end of the tile loop ensures the tile can be safely overwritten on the next iteration. Missing either sync is a race condition.

Tile size: start with \texttt{BLOCKSIZE=32} (each tile is $32 \times 32 \times 4 = 4$~KB, well within the 48~KB shared memory default). Try \texttt{BLOCKSIZE=64} (16~KB per tile) and \texttt{BLOCKSIZE=128} (64~KB---requires increasing shared memory allocation with \texttt{cudaFuncSetAttribute}).

\textbf{Expected gain:} 5--10$\times$ over Kernel~2 for large $N$.

\medskip
\textbf{Kernel 4: 1D Blocktiling---Each Thread Computes TM Output Elements.}
Rather than one thread per output element, assign each thread a column of \texttt{TM} consecutive output elements. The thread registers accumulate \texttt{TM} partial sums (one per output row in its tile), reusing the loaded $B$ element across all \texttt{TM} accumulators. This increases arithmetic intensity per thread and reduces the number of threads needed (block occupancy changes accordingly).

Register caching: before the inner $k$ loop, load the single $B[k][col]$ value into a register. Then loop over the \texttt{TM} rows, loading $A[row+i][k]$ from shared memory and multiplying by the cached $B$ value. Without caching $B$, shared memory is read $TM \times$ per $k$ iteration; with caching, it is read once.

\textbf{Tune:} Try \texttt{TM} $\in \{4, 8, 16\}$. Profile register usage with \texttt{ncu --metrics sm\_\_maximum\_warps\_per\_active\_cycle\_pct}.

\medskip
\textbf{Kernel 5: 2D Blocktiling---Each Thread Computes TM$\times$TN Elements.}
Extend to 2D: each thread accumulates a \texttt{TM}$\times$\texttt{TN} sub-tile in registers. Cache a column of \texttt{TM} $A$ elements and a row of \texttt{TN} $B$ elements; compute the outer product to update the \texttt{TM}$\times$\texttt{TN} accumulator. This is the fundamental GEMM microkernel structure: load, outer product, accumulate.

Boehm's parameters: \texttt{BM=128, BN=128, BK=8, TM=8, TN=8}. Each thread holds $8 \times 8 = 64$ FP32 accumulators plus 8 $A$ registers plus 8 $B$ registers---79 registers total. H100 has 65536 registers per SM and 1024 threads/block maximum; at 79 registers/thread you can have $\lfloor 65536 / 79 \rfloor \approx 829$ threads before register-limited occupancy. Verify: \texttt{ncu --metrics launch\_\_registers\_per\_thread}.

\textbf{Expected gain:} This is the largest single jump in the progression---often 2--3$\times$ over Kernel~4.

\medskip
\textbf{Kernel 6: Vectorized Loads (float4 = 128-bit).}
Replace scalar \texttt{float} loads from global memory with \texttt{float4} loads (4 FP32 values = 128 bits). A single \texttt{float4} load issues one load instruction and fills 4 registers, vs.\ 4 separate instructions. This halves the number of issued load instructions and better utilizes the 128-byte cache line width.

Implementation: cast pointers to \texttt{float4*} for global-to-shared-memory loads. Alignment requirement: the starting address must be 16-byte aligned (guaranteed if your matrix is allocated with \texttt{cudaMalloc} with standard alignment).

\textbf{Profile:} \texttt{ncu --metrics l1tex\_\_t\_bytes\_pipe\_lsu\_mem\_global\_op\_ld.sum} before and after. The metric value in bytes should be unchanged; the metric in sectors should drop by $\sim 4\times$ if vectorization is working.

\medskip
\textbf{Kernel 7: Warptiling---Hierarchical Block $\to$ Warp $\to$ Thread.}
This is the structure that approaches cuBLAS. Add an explicit warp-level tile between the block tile and the thread tile. A block of 128$\times$128 output is divided into warp tiles of 64$\times$64; each warp of 32 threads computes its 64$\times$64 tile. Within the warp tile, each thread computes a 16$\times$16 sub-tile using WMMA or direct MMA instructions (on H100, use \texttt{mma.sync.aligned.m16n8k16} PTX instructions or the WMMA C++ API).

The benefit: warp-level tiling exposes the register file layout required for MMA instructions without explicit PTX. WMMA fragments hold $A$, $B$, and accumulator data in a warp-distributed layout---each lane holds a portion of the fragment. Loading into WMMA fragments from shared memory requires the \texttt{wmma::load\_matrix\_sync} operation.

\textbf{Boehm's result:} 93.7\% of cuBLAS on A6000. On H100, cuBLAS is more heavily tuned; targeting $>70\%$ of cuBLAS is a realistic bar for a hand-written kernel without CUTLASS internals.

\textbf{Final record:} For each kernel $k \in \{0, \ldots, 7\}$, fill in: TFLOP/s, \% of cuBLAS, top Nsight bottleneck metric, one-sentence causal explanation of the gain over the previous kernel. This table is your portfolio evidence.
\end{exercise}

% -----------------------------------------------------------------------
\subsection{Part C: Blackwell-Specific Optimizations}
% -----------------------------------------------------------------------

\begin{exercise}{Blackwell Kernels 8--10: TMEM, FP4, Clusters (6--8 hours)}
Rent B200 access for this section. Compile with \texttt{-arch=sm\_100}. These kernels are not in the Boehm repo---you are building them from scratch using the Blackwell architecture documentation and CUDA 12.6+ API references.

\medskip
\textbf{Kernel 8: TMEM Accumulator Storage.}
Blackwell introduces Tensor Memory (TMEM), a new on-chip memory space that can hold accumulator tiles for Tensor Core operations. Unlike shared memory (which is programmer-managed L1 cache visible to all threads in a block), TMEM is reserved for tensor core accumulators and requires explicit allocation via \texttt{cuda::ptx::tcgen05\_alloc} (or the equivalent intrinsic in CUDA 12.6).

\textbf{Implementation steps:}
\begin{enumerate}
  \item Start from Kernel~7 (warptiling with WMMA). The accumulator fragments are currently in registers.
  \item Allocate TMEM for the accumulator: \texttt{cuda::ptx::tcgen05\_alloc(\&tmem\_ptr, \textit{num\_cols})}. TMEM is addressed as a 2D structure; each ``column'' holds a portion of the accumulator for one MMA tile.
  \item Replace WMMA accumulator fragments with TMEM-resident accumulators using \texttt{mma.sp::ordered\_metadata} instructions or the CUDA 12.6 TMEM load/store APIs.
  \item After the $k$-loop, store the TMEM accumulator back to global memory via shared memory as a staging area.
\end{enumerate}

\textbf{Compare to Kernel~7:} Measure register usage (expect a reduction), occupancy (expect an increase), and TFLOP/s. If TMEM reduces register pressure enough to increase the number of active warps per SM, you should see a latency-hiding benefit. Profile with \texttt{ncu --metrics sm\_\_maximum\_warps\_per\_active\_cycle\_pct}.

\textbf{Write before measuring:} If TMEM reduces per-thread register count from 79 to 55 (hypothetical), how does SM occupancy change at the block size from Kernel~7? Does higher occupancy translate to higher TFLOP/s for this kernel? Why or why not? (Hint: if the kernel is compute-bound, more warps per SM only help by hiding inter-instruction latency, not by increasing peak FLOP/s.)

\medskip
\textbf{Kernel 9: FP4/FP6 Mixed Precision with FP32 Accumulation.}
Blackwell natively supports FP4 (E2M1 format: 1 sign bit, 2 exponent bits, 1 mantissa bit) and FP6 (E2M3 or E3M2) in hardware Tensor Cores. For RLVR inference, FP4 weights with FP32 accumulation may offer the best throughput-per-quality tradeoff: 4$\times$ fewer bytes loaded from HBM, hardware accumulation in FP32 to preserve gradient fidelity.

\textbf{Implementation:}
\begin{enumerate}
  \item Quantize a test weight matrix $W$ from FP16 to FP4 using symmetric per-channel quantization. Store as packed \texttt{uint8} (two FP4 values per byte).
  \item Write a kernel that unpacks FP4 pairs in shared memory and feeds them to Blackwell FP4 MMA instructions: \texttt{mma.sync.aligned.m16n8k32.row.col.f32.e4m3.e4m3.f32} (check the PTX ISA for Blackwell-specific variants).
  \item Accumulate into FP32 registers or TMEM (combine with Kernel~8 if time permits).
  \item Compare to FP16 Kernel~7 on the same matrix size: report TFLOP/s, memory bandwidth utilization, and---critically---the absolute difference in output values (FP4 vs.\ FP16 as ground truth). This quality metric is the signal Anthropic cares about for RLVR stability.
\end{enumerate}

\textbf{For RLVR generation specifically:} Run a microbenchmark that models the weight-load pattern during autoregressive decoding: batch size 1, weight matrix load from HBM, matmul, next-token prediction. Compare FP4 vs.\ FP16 in \textit{tokens per second}, not TFLOP/s. At batch size 1, you are memory-bandwidth-bound; the FP4 benefit is 4$\times$ fewer bytes loaded, which should approach a 4$\times$ throughput improvement modulo overhead.

\medskip
\textbf{Kernel 10: Thread Block Clusters (Size 16).}
Blackwell allows clusters of up to 16 thread blocks that can communicate via distributed shared memory (DSMEM) over the NVLink fabric. Each block in the cluster can directly read from another block's shared memory using \texttt{cuda::cluster::map\_shared\_rank} and the cluster barrier API.

\textbf{Why this matters for matmul:} A single block is limited to its 228~KB shared memory. A cluster of 16 blocks has access to $16 \times 228~\text{KB} = 3.6~\text{MB}$ of aggregated DSMEM. This allows a single cluster to hold a much larger tile of $A$ and $B$ in on-chip memory, dramatically increasing data reuse and reducing HBM traffic for large $N$.

\textbf{Implementation steps:}
\begin{enumerate}
  \item Declare a cluster of size 16 in the kernel launch configuration: \texttt{cudaLaunchConfig\_t cfg; cfg.gridDim = \{...\}; cfg.clusterDim = \{16, 1, 1\};}.
  \item Within the kernel, use \texttt{cluster.sync()} barriers instead of \texttt{\_\_syncthreads()} to synchronize across the cluster.
  \item Partition the output matrix so that each cluster computes a large tile; each block within the cluster computes a sub-tile and loads its portion of $A$ and $B$ into its own shared memory.
  \item Use \texttt{cuda::cluster::load} to read $A$ and $B$ fragments that reside in neighboring blocks' shared memory when needed.
\end{enumerate}

\textbf{Compare to Kernel~7 (non-clustered):} For $N = 8192$, report TFLOP/s and HBM read bytes (from Nsight). The cluster kernel should show lower HBM traffic for the same computation---the tile reuse is happening in DSMEM rather than HBM. If the DSMEM bandwidth (NVLink) is the new bottleneck, profile with the cluster-specific Nsight metrics.

\textbf{Key question:} At what matrix size does Kernel~10 outperform Kernel~7? For small $N$, cluster overhead (barrier synchronization, DSMEM addressing) may negate the tile-reuse benefit. Find the crossover empirically.
\end{exercise}

% -----------------------------------------------------------------------
\subsection{Part D: Triton Comparison}
% -----------------------------------------------------------------------

\begin{exercise}{Triton Matmul vs.\ Hand-Written CUDA (2--3 hours)}
Write the same FP16 matmul in Triton using \texttt{tl.dot}, \texttt{tl.load}, and \texttt{tl.store}. The canonical Triton matmul tutorial (docs.triton-lang.org) is the starting point; the exercise is in the profiling and gap analysis, not the implementation.

\begin{enumerate}
  \item Implement Triton matmul with \texttt{@triton.autotune} over \texttt{BLOCK\_M, BLOCK\_N, BLOCK\_K} $\in \{64, 128\}$ and \texttt{num\_warps} $\in \{4, 8\}$.
  \item Benchmark at $N \in \{512, 1024, 2048, 4096, 8192\}$ and record TFLOP/s alongside your Kernel~7 result.
  \item Profile both kernels with \texttt{ncu}. Compare: register usage, shared memory usage, warp occupancy, instruction mix (load vs.\ compute).
\end{enumerate}

\textbf{Specific questions to answer:}
\begin{itemize}
  \item At what matrix size does Triton match or exceed Kernel~7? (Expect Triton to be competitive at $N \geq 2048$ and potentially fall behind at $N = 512$ where the autotuner has fewer layout options.)
  \item Where does Triton fall short? Use Nsight to identify whether the gap is in: (a) suboptimal tile size selected by autotuner, (b) higher register usage due to compiler conservatism, (c) different shared memory layout causing more bank conflicts, or (d) worse warp-level instruction scheduling.
  \item For which RLVR operations would you choose Triton in production? Which would you hand-write? Justify with the profiling evidence from this exercise.
\end{itemize}

\textbf{Note:} If you are running on B200 and want to test Triton's FP4 support, check the Triton changelog for Blackwell backend status (as of early 2026, support is in active development). If available, compare Triton FP4 to your Kernel~9 result.
\end{exercise}

% -----------------------------------------------------------------------
\subsection{Part E: Flash Attention}
% -----------------------------------------------------------------------

\begin{exercise}{Flash Attention: Implementation and Analysis (3--4 hours)}
Flash Attention is the most consequential single-kernel optimization in recent ML systems history. For RLVR training, attention over long-context rollouts (sequences of 8K--32K tokens) is a significant wall-time component. Understanding Flash Attention is a prerequisite for any serious discussion of inference optimization at frontier scale.

\textbf{The key insight---memory-IO analysis:}
\begin{itemize}
  \item \textbf{Standard attention:} Compute $S = QK^\top / \sqrt{d}$ ($O(N^2 d)$ FLOPs, reads $Q, K \in \mathbb{R}^{N \times d}$, writes $S \in \mathbb{R}^{N \times N}$ to HBM---$O(N^2)$ bytes); apply softmax to $S$ (reads $N^2$, writes $N^2$); compute $O = S V$ ($O(N^2 d)$ FLOPs, reads $S$ and $V$, writes $O$). Total HBM reads/writes: $\Theta(N^2 + Nd)$---dominated by the $N \times N$ attention matrix at large $N$.
  \item \textbf{Flash Attention:} Process $Q$, $K$, $V$ in blocks. For each block of $Q$, iterate over blocks of $K$ and $V$, maintaining a running softmax normalization (the ``online softmax trick''). The $N \times N$ attention matrix is \textit{never materialized in HBM}---it exists transiently in on-chip shared memory. Total HBM reads/writes: $\Theta(Nd)$---linear in sequence length.
\end{itemize}

\textbf{Online softmax trick:} Standard softmax over a row requires two passes: first find the maximum (for numerical stability), then compute $\exp(x_i - \max) / \sum \exp(x_j - \max)$. The online algorithm maintains a running $(m, \ell)$ where $m$ is the current maximum and $\ell$ is the current sum, updating both as new elements arrive. This enables single-pass softmax over a sequence that does not fit in shared memory.

\textbf{Exercise options (choose based on time):}
\begin{enumerate}
  \item \textbf{Study option (2 hours):} Read the Flash Attention v2 paper (Dao 2023) and the Triton Flash Attention implementation in the Triton repo (\texttt{python/triton/ops/flash\_attention.py}). Trace through the forward pass, identifying: (a) which tensors live in HBM vs.\ shared memory at each step, (b) where the online softmax state is maintained, (c) how the backward pass handles recomputation instead of storing the $N \times N$ matrix.
  \item \textbf{Implementation option (4 hours):} Implement Flash Attention forward pass in Triton. Start from the Triton tutorial (\texttt{triton-lang.org/main/getting-started/tutorials/06-fused-attention.html}). Verify correctness against the PyTorch reference implementation. Benchmark: plot throughput (tokens/sec) vs.\ sequence length $N \in \{512, 1024, 2048, 4096, 8192\}$. At what $N$ does Flash Attention outperform standard PyTorch attention? (Expect crossover around $N = 512$--$1024$ on H100.)
\end{enumerate}

\textbf{Blackwell relevance:} Flash Attention benefits compound on B200's higher HBM bandwidth. The $\Theta(Nd)$ HBM traffic for Flash Attention scales directly with bandwidth; the $\Theta(N^2)$ traffic for standard attention does too, but the ratio of avoided bytes is $N/d$---at $N=8192$ and $d=128$, Flash Attention avoids $64\times$ more HBM traffic. On B200's 8~TB/s vs.\ H100's 3.35~TB/s, this means Flash Attention's absolute time is lower, but the \textit{relative} benefit of Flash Attention vs.\ standard attention is the same. The bandwidth upgrade does not change the algorithmic advantage; it scales it.
\end{exercise}

% -----------------------------------------------------------------------
\subsection{Part F: Blackwell Inference Analysis}
% -----------------------------------------------------------------------

\begin{exercise}{Amdahl's Law for RLVR on B200 (2--3 hours)}
B200 marketing claims 5$\times$ inference throughput over H100. Apply Amdahl's law rigorously to determine the realistic system speedup for RLVR training.

\textbf{Setup:} RLVR training spends approximately 60\% of wall time on generation (autoregressive decoding) and 40\% on the training step (forward pass, backward pass, optimizer update).

\textbf{Naive calculation:} If generation is 5$\times$ faster on B200:
$$\text{Speedup}_{\text{naive}} = \frac{1}{0.40 + 0.60/5} = \frac{1}{0.52} \approx 1.92\times$$

But generation itself is not a monolith. Break it into components and apply B200's improvements to each:

\begin{center}
\begin{tabular}{lrrl}
\hline
\textbf{Component} & \textbf{Fraction of gen.} & \textbf{B200 speedup} & \textbf{Reason} \\
\hline
Attention (KV-cache reads) & 40\% & $\sim 2.4\times$ & Bandwidth: 8/3.35 \\
Weight matmul (decode) & 30\% & $\sim 2.4\times$ & Bandwidth-bound \\
Sampling / argmax & 10\% & $\sim 1.5\times$ & Partially memory-bound \\
KV cache management & 20\% & $\sim 2.4\times$ & Bandwidth-bound \\
\hline
\end{tabular}
\end{center}

\textbf{Work out:}
\begin{enumerate}
  \item Compute the effective speedup for the generation phase as the harmonic mean of component speedups weighted by their fractions. (Note: the arithmetic weighted average is the correct formula here, not the harmonic mean---each component's speedup reduces its fraction of total time; use the Amdahl formula component by component.)
  \item Plug the resulting generation speedup into the top-level Amdahl formula. What is the realistic system speedup?
  \item Now consider FP4 quantization: the weight matmul component (30\% of generation) may see $4\times$ speedup from FP4 due to 4$\times$ fewer bytes loaded. How does this change the calculation?
  \item If B200's 5$\times$ marketing claim assumes all generation components improve equally, what component fraction and speedup would produce exactly 5$\times$? Is this realistic given the memory-bandwidth analysis?
\end{enumerate}

\textbf{Write a 200-word memo} summarizing your analysis. Claim: ``The B200 delivers X$\times$ realistic speedup for RLVR training, not the marketed 5$\times$, because [specific components] are [bandwidth/compute]-bound and improve only [fraction]$\times$.'' This memo---backed by arithmetic from your kernel measurements in Parts B--C---is a concrete artifact demonstrating systems reasoning at the level Anthropic hires for.
\end{exercise}

% -----------------------------------------------------------------------
\subsection{Portfolio Project}
% -----------------------------------------------------------------------

\begin{portfolio}{CUDA/Blackwell Kernel Optimization Guide}
\textbf{Output options:}
\begin{itemize}
  \item \textbf{Primary:} Blog post titled ``What Blackwell-specific optimizations (TMEM, FP4, clusters) actually matter for RLVR inference---and which are marketing?'' Published on your personal site or Substack. Target length: 3000--4000 words with profiling screenshots and the table from Part B's exercise.
  \item \textbf{Alternative} (if B200 cloud is not available): Submit the Boehm progression (Kernels 0--7) to the MLSys 2026 FlashInfer contest or equivalent open benchmark. The H100 results alone are credible portfolio evidence; the Blackwell analysis can be done analytically using published specs.
\end{itemize}

\textbf{The monkey:} The blog's core argument is the distinction between hardware improvements that deliver real RLVR throughput and those that are marketing. The kernel implementations are the pedestal---they establish that you can actually write these kernels and measure them. The monkey is the analysis: which of TMEM, FP4, and clusters gives a measurable RLVR speedup, and what fraction of the marketed 5$\times$ improvement is real at the system level?

\textbf{Verification criteria:}
\begin{itemize}
  \item At least 70\% of cuBLAS on the warptiling kernel (Kernel~7) at $N=4096$ on H100.
  \item Blackwell-specific kernels (8--10) show measurable improvement over Hopper equivalents for at least one operation---document the specific metric (TFLOP/s, memory transactions, or tokens/sec for generation).
  \item Amdahl's law analysis (Part F) backs the blog's central claim with numbers.
\end{itemize}

\textbf{Hardware budget:} H100 at $\sim$\$2/hr for Kernels 0--7 (estimate 10--15 hours of actual GPU time, including debugging: $\sim$\$20--30); B200 at $\sim$\$2.25/hr for Kernels 8--10 (estimate 6--10 hours: $\sim$\$14--23). Total: $\sim$\$35--55 in GPU time for the full chapter. Budget $\sim$\$250/month allows multiple iteration cycles.

\textbf{Why this matters:} Simon Boehm's SGEMM\_CUDA post is the canonical example of kernel credibility. This project extends it to Blackwell at a higher level of strategic relevance. A senior engineer reading your blog can directly use the Amdahl analysis to inform hardware purchasing decisions and RLVR infrastructure planning. That is the difference between a tutorial and a contribution.
\end{portfolio}

\newpage

\section{TPU/Pallas Kernel Optimization}

\smallskip\noindent\textit{Study Guide companion: Study Guide Section 12 (Hardware Deep Dive) covers the interview-ready version of this material.}
\smallskip

The RLVR training pipeline is ``held together by duct tape.'' Every percentage point of hardware utilization matters at scale. Simon Boehm's blog post walking through CUDA matmul optimization from 1.3\% to 93.7\% of cuBLAS remains the reference for GPU kernel pedagogy. No equivalent guide exists for TPUs and Pallas. This chapter follows the same methodology: start naive, profile, fix one thing, measure, repeat.

\subsection{Conceptual Foundation}

\begin{thinkfirst}{GPU vs.\ TPU mental models}
The biggest conceptual shift: on GPUs you think about \textbf{threads and memory access patterns}. On TPUs you think about \textbf{pipeline efficiency and keeping the systolic array fed}.

\begin{enumerate}
  \item \textbf{The GPU model:} Thousands of lightweight threads execute the same instruction on different data (SIMT). The bottleneck is usually memory bandwidth --- getting data to the threads fast enough. Key optimizations: coalescing memory accesses, using shared memory, avoiding bank conflicts.

  \item \textbf{The TPU model:} A systolic array (128$\times$128 or 256$\times$256 multiply-accumulate units) flows data through in a deterministic pattern. There's no warp scheduling randomness. The bottleneck is keeping data flowing through the array without bubbles. Key optimizations: tiling to match MXU dimensions, pipelining HBM-to-VMEM transfers, maximizing arithmetic intensity.

  \item \textbf{Why this matters for RLVR:} During RLVR training, the pipeline alternates between generation (autoregressive, memory-bandwidth-bound) and training (matmul-heavy, compute-bound). RLVR runs are particularly long because verification-based rewards require generating complete solutions and checking them end-to-end --- the generation phase is therefore even longer per update than in simpler RL setups. The optimization strategy is different for each phase. On a TPU, the generation phase underutilizes the MXU.

  Now answer: if RLVR training spends 60\% of time on generation and 40\% on gradient updates, which phase should you optimize first? What's the theoretical maximum speedup from optimizing just that phase? (Hint: Amdahl's law.)
\end{enumerate}
\end{thinkfirst}

\begin{thinkfirst}{Memory hierarchy reasoning}
TPU memory hierarchy: HBM (off-chip, large, slow) $\to$ VMEM (on-chip SRAM, small, fast) $\to$ VREGs (registers).

\begin{enumerate}
  \item A TPU v5e has 16GB HBM at 820 GB/s bandwidth and $\sim$30 MB VMEM. You want to multiply two 4096$\times$4096 bf16 matrices. Each matrix is 32 MB --- neither fits in VMEM. How do you tile this computation so each tile fits in VMEM while maximizing MXU utilization?

  \item The MXU does 128$\times$128 matmuls per cycle. If your tile is 64$\times$64, what fraction of the MXU is utilized? What about 256$\times$256? Work out the utilization for each and explain the tradeoff between tile size and MXU efficiency.

  \item Pallas uses ``BlockSpec'' to define tiles and ``Buffered(buffer\_count=2)'' to double-buffer HBM$\to$VMEM transfers. Draw the pipeline on paper: while tile $N$ is being computed on the MXU, tile $N+1$ is being loaded into the second VMEM buffer via DMA. What determines whether you're compute-bound or memory-bound in this pipelined setup?
\end{enumerate}
\end{thinkfirst}

\begin{thinkfirst}{Arithmetic intensity and the roofline model}
Arithmetic intensity = FLOPs / bytes transferred. For a matrix multiply: $2mnk$ FLOPs, $(mn + mk + nk) \times \text{bytes\_per\_element}$ transferred. For square matrices: intensity $\approx n/3$ (for bf16).

\begin{enumerate}
  \item A TPU v5e has $\sim$200 TFLOP/s compute and 820 GB/s memory bandwidth. At what matrix size does a bf16 matmul become compute-bound rather than memory-bound? Use the roofline model: compute the crossover arithmetic intensity, then solve for $n$.

  \item RLVR training involves many small matmuls (e.g., reward model forward passes on individual completions). Are these compute-bound or memory-bound? What does this mean for optimization strategy?

  \item The autoregressive generation phase produces one token at a time. Each step involves a matmul of shape $[1, d] \times [d, V]$ where $V$ is vocabulary size. Calculate the arithmetic intensity. Is this compute-bound or memory-bound? (This is why generation is the bottleneck --- it's always memory-bandwidth-bound.)
\end{enumerate}
\end{thinkfirst}

\subsection{Hands-On Exercises}

\begin{exercise}{Pallas matmul: naive to optimized (10--15 hours)}
Requires: TPU access (Google Cloud TPU Research program or Colab TPU).

Write a Pallas matmul kernel and optimize it step by step, measuring TFLOP/s at each stage:

\textbf{Kernel 0: JAX baseline.} \texttt{jnp.dot(A, B)}. This is XLA's auto-tuned kernel. Record as your ceiling.

\textbf{Kernel 1: Naive Pallas.} Single block, no tiling. Process the entire matrix in one kernel invocation.
\begin{lstlisting}
def naive_matmul_kernel(x_ref, y_ref, z_ref):
    z_ref[...] = x_ref[...] @ y_ref[...]

result = pl.pallas_call(
    naive_matmul_kernel,
    out_shape=jax.ShapeDtypeStruct((m, n), jnp.float32),
    grid=(1,),
    in_specs=[pl.BlockSpec((m, k), lambda i: (0, 0)),
              pl.BlockSpec((k, n), lambda i: (0, 0))],
    out_specs=pl.BlockSpec((m, n), lambda i: (0, 0)),
)(x, y)
\end{lstlisting}

\textbf{What to measure:} TFLOP/s. This will be terrible because the matrices probably don't fit in VMEM.

\textbf{Kernel 2: Block tiling.} Add a grid and proper BlockSpecs. Tile along the reduction dimension.

\textbf{Kernel 3: FP32 accumulator.} Use bf16 inputs but accumulate in FP32 (use \texttt{scratch\_shapes} for the accumulator in VMEM).

\textbf{Kernel 4: Pipeline the DMA.} Add \texttt{Buffered(buffer\_count=2)} to overlap DMA with compute.

\textbf{Kernel 5: Block size sweep.} Try $bm, bn, bk \in \{64, 128, 256, 512\}$. Plot TFLOP/s vs.\ block size. Why does 128 usually win? (It matches the MXU dimension.)

\textbf{Kernel 6: Megacore.} Set \texttt{dimension\_semantics=("parallel", "parallel", "arbitrary")} to use both TensorCores. Expected: $\sim$2$\times$ on supported chips.

\textbf{At each step:} Record TFLOP/s and \% of JAX baseline. Each optimization should produce a measurable improvement --- if it doesn't, diagnose why before moving on. By Kernel 4--5, you should be at $>50\%$ of the JAX baseline. If you're below 20\% after Kernel 3, something fundamental is wrong with your tiling.
\end{exercise}

\begin{exercise}{Extended Pallas Kernels (6--8 hours)}
Continue directly from Kernel 6. These four kernels push into production-quality territory: fused ops, mixed precision, chip-specific tuning, and distributed execution.

\textbf{Kernel 7: Fused transpose + activation.} Write a single Pallas kernel that applies a matrix transpose followed by an element-wise activation (GELU or SiLU) without writing the intermediate transposed result back to HBM.

\begin{itemize}
  \item Load a tile from HBM, transpose it in VMEM registers, apply the activation, write the fused result.
  \item \textbf{Measure:} TFLOP/s improvement vs.\ two separate kernels (transpose then activation). What fraction of the speedup comes from eliminating the intermediate HBM write?
  \item \textbf{Think:} GELU requires $\tanh$ --- how does the TPU handle transcendental functions? Is it approximated? Does SiLU (which avoids $\tanh$) outperform GELU on TPUs specifically?
\end{itemize}

\textbf{Kernel 8: Custom precision paths.} bf16 input, fp32 accumulate, bf16 output.

\begin{itemize}
  \item Use \texttt{scratch\_shapes} to allocate an fp32 accumulator in VMEM. Accept bf16 inputs via \texttt{in\_specs}, cast to fp32 on load, accumulate, cast back to bf16 on store.
  \item \textbf{Measure:} Compare gradient norm variance across 100 training steps vs.\ pure bf16 accumulation. Numerical stability improvement should be visible if your learning rate is high or your model is deep.
  \item \textbf{Think:} Pure bf16 accumulation introduces quantization error on every add. At what model size or learning rate does this become practically significant? Cross-check against JAX's default XLA behavior --- does XLA already do this internally?
\end{itemize}

\textbf{Kernel 9: TPUv6e-specific optimization (if accessible).} TPUv6e has a larger MXU and more VMEM than v5e. Re-tune block sizes for v6e architecture.

\begin{itemize}
  \item Re-run your Kernel 5 block size sweep on v6e. Does the optimal $(bm, bn, bk)$ change?
  \item If v6e is not accessible, use the TPU Research Cloud application to request access (this is worth doing regardless --- document the application process as part of your writeup).
  \item \textbf{Think:} If the optimal block size scales with MXU dimensions, what block sizes would you predict for v6e before running the sweep? Make the prediction first, then measure.
\end{itemize}

\textbf{Kernel 10: Distributed matmul with RDMA.} Promote the ring all-reduce exercise below into a full kernel step: implement a distributed matmul where each device computes a shard and communicates partial results.

\begin{itemize}
  \item Partition matrix $A$ along rows across $N$ devices. Each device computes its shard of $C = A_{\text{shard}} \cdot B$. Use ring RDMA (see the distributed exercise below) to accumulate partial sums.
  \item \textbf{Measure:} Effective TFLOP/s per device as $N$ scales from 1 to 8. Does it scale linearly? Where does communication overhead dominate?
  \item \textbf{Think:} This is a simplified version of tensor parallelism. What would need to change to make this production-viable in a real RLVR training loop?
\end{itemize}
\end{exercise}

\begin{exercise}{Profile the bottleneck in RLVR training (4--6 hours)}
Profile an actual RLVR training step (use TRL's \texttt{GRPOTrainer} or your nanoRLVR from Chapter 8).

\begin{enumerate}
  \item Use \texttt{jax.profiler} or PyTorch profiler to trace one full GRPO update step.
  \item Break down time: what fraction is generation (sampling completions)? What fraction is forward pass (reward model or verifier)? What fraction is backward pass (policy gradient)?
  \item Identify the actual bottleneck. Is it what you expected? Note: because RLVR requires complete solution generation for verification, generation typically dominates even more than in RLHF-style training.
  \item For the bottleneck operation: is it compute-bound or memory-bound? Calculate the arithmetic intensity to verify.
  \item Propose \textit{one} optimization for the bottleneck. Implement it if you can. Measure the speedup.
  \item \textbf{Write a 1-paragraph memo:} You have one week of engineering time. Based on your profiling, which optimization would you prioritize for the RLVR training pipeline, and why? What's the expected speedup vs.\ engineering cost? This is the judgment call a senior RE makes daily.
\end{enumerate}

\end{exercise}

The point of this exercise: \textbf{measure before optimizing}. If you assumed the bottleneck before profiling and got it wrong, that's the lesson. Diagnosing the real bottleneck rather than the assumed one is exactly what ``attacking the problem directly'' looks like.

\begin{exercise}{Distributed Pallas: ring all-reduce (8--12 hours)}
Requires: multi-chip TPU access (TPU v4-8 or larger pod slice).

Implement a ring all-reduce in Pallas using RDMA primitives:

\begin{enumerate}
  \item Each device holds a shard of the gradient tensor.
  \item In $N$ steps (where $N$ = number of devices), each device sends its shard to the right neighbor and accumulates the incoming shard.
  \item After $N$ steps, every device has the complete sum.
\end{enumerate}

\textbf{Key Pallas primitives:}
\begin{lstlisting}
copy = pltpu.make_async_remote_copy(
    src_ref=local_buf,
    dst_ref=remote_buf,
    send_sem=send_sem,
    recv_sem=recv_sem,
    device_id=right_neighbor
)
copy.start()    # Non-blocking
# ... compute on local data ...
copy.wait_recv()  # Block until received
\end{lstlisting}

\textbf{Questions to answer:}
\begin{enumerate}[label=\alph*)]
  \item Derive the theoretical bandwidth utilization of ring all-reduce as a function of $N$ (number of devices). Does it scale well?
  \item Measure actual vs.\ theoretical bandwidth. Where's the gap?
  \item Can you overlap the all-reduce communication with local computation? (Hint: start the DMA, compute on local data, then wait for completion.)
\end{enumerate}
\end{exercise}

\begin{thinkfirst}{TPUv7 Ironwood: Inference-First Chip for Training?}
TPUv7 Ironwood is Google's latest TPU, designed as an inference-first chip: 192GB HBM3E, optimized for large-scale serving workloads. Anthropic has access to 1M Ironwoods and runs RLVR training on TPUs. This creates a tension worth reasoning through carefully.

\begin{enumerate}
  \item \textbf{The roofline question.} Ironwood has 192GB HBM3E --- the same capacity as MI300X. But memory \textit{bandwidth} and peak \textit{compute} may be balanced differently from v5e. Recall the roofline crossover: ridge point = peak compute / memory bandwidth. If Ironwood has higher bandwidth but the same or lower peak compute than v5e, what happens to the ridge point? Does this make it easier or harder to stay in the compute-bound regime for training matmuls?

  \item \textbf{RLVR's dual personality.} RLVR training interleaves two very different workloads: generation (autoregressive, memory-bandwidth-bound) and gradient updates (dense matmul, compute-bound). If Ironwood is optimized for the inference-like generation phase, does it actually \textit{help} RLVR more than a training-optimized chip would? Sketch the utilization profile across a single RLVR step on Ironwood vs.\ v5e.

  \item \textbf{Argue both sides.} The inference-first design may feature higher memory bandwidth but lower peak compute. For RLVR where generation dominates wall time:
  \begin{itemize}
    \item \textbf{Argument FOR Ironwood being the right chip:} Generation is memory-bandwidth-bound; higher bandwidth directly reduces wall time in the dominant phase. RLVR's bottleneck is where Ironwood is strongest.
    \item \textbf{Argument AGAINST:} The gradient update phase, even if shorter, still gates convergence. If Ironwood is significantly weaker on dense matmuls, the training phase takes longer per batch, limiting how fast you can improve the policy. A training-optimized chip that keeps the ratio of generation-to-update time better balanced might converge in fewer wall-clock hours overall.
  \end{itemize}
  Which argument do you find more convincing? What empirical measurement would settle the question?

  \item \textbf{The 1M chip question.} Anthropic reportedly has 1M Ironwoods. At that scale, even a 5\% improvement in RLVR throughput has enormous leverage. What single kernel optimization from this chapter --- if ported to Ironwood --- would have the highest expected impact? (You may not have Ironwood access; reason from first principles about the architecture.)
\end{enumerate}
\end{thinkfirst}

\begin{portfolio}{TPU Optimization Guide: ``The Pallas Companion''}
Write the Simon Boehm blog post but for TPUs. An iterative, pedagogical guide to writing high-performance Pallas kernels, with profiling at every step.

\textbf{Structure (10 kernels):}
\begin{enumerate}
  \item JAX baseline (XLA auto-tuned)
  \item Naive Pallas (single block)
  \item Block tiling with grid
  \item FP32 accumulator in VMEM scratch
  \item Double-buffered pipelining
  \item Block size autotuning
  \item Megacore parallelism
  \item Fused operations (transpose, activation)
  \item Custom precision (bf16 input, fp32 accumulate)
  \item Distributed matmul with RDMA
\end{enumerate}

\textbf{What makes this portfolio-quality:}
\begin{itemize}
  \item Performance table at each step (\% of peak, TFLOP/s)
  \item Roofline model analysis showing compute vs.\ memory bound transitions
  \item Diagrams of memory hierarchy, pipeline stages, and tiling patterns
  \item Honest documentation of what didn't work and why
  \item Comparison of TPU vs.\ GPU optimization strategies
\end{itemize}

\textbf{The monkey:} ``Can you teach the optimization journey clearly enough that a reader with CUDA experience but no TPU experience reaches $>$50\% of peak on their first try?'' Writing the kernels is the pedestal. Writing the \textit{explanation} that makes the reader understand \textit{why} each optimization works is the monkey.

\textbf{Why this fills a gap:} Boehm's CUDA guide is ``the reference'' for GPU optimization. No equivalent exists for TPUs. The official Pallas docs are reference material, not pedagogy.

\textbf{Hardware note on Ironwood:} The results in this guide are developed on TPU v5e (and v6e if accessible via TRC). For Ironwood-specific results, apply to the Google TPU Research Cloud program and request v7 access when available. The findings on v5e extrapolate \textit{architecturally} but not \textit{quantitatively} to Ironwood --- the optimal block sizes, roofline ridge points, and bandwidth utilization numbers will differ due to the inference-first memory architecture. Documenting the delta between v5e and Ironwood predictions vs.\ measurements is itself a valuable contribution.

\textbf{Verification:} Find one person with CUDA experience and no TPU experience. Have them follow your guide. If they reach 50\% of peak TFLOP/s, you've succeeded.
\end{portfolio}

\newpage

\section{RL Inference \& Serving Efficiency}

\smallskip\noindent\textit{Study Guide companion: Study Guide Section 9 (Questions That Show Depth --- Hardware \& Generation Bottleneck) covers the interview-ready version of this material.}
\smallskip

\begin{whythis}
Dario Amodei put it plainly: ``This isn't a research problem. This is an engineering and inference problem.'' RLVR training runs last weeks, and 50--70\% of that wall time is spent on autoregressive token generation --- not on gradient computation, not on optimizer updates, but on the forward pass that produces rollouts. A 2$\times$ speedup in matmul kernels might recover a few percent of overall runtime. A 2$\times$ speedup in generation throughput cuts your training time nearly in half.

The asymmetry matters for compounding reasons. Every gradient step requires a fresh batch of rollouts. If rollout generation takes 5 minutes and the gradient step takes 1 minute, you run 6-minute cycles for weeks. Halve generation time and you run 3.5-minute cycles --- not 2$\times$ faster overall, but 1.7$\times$ faster, which over a 3-week run is 6 days of recovered wall time per run. Across dozens of experiments, that is months.

This chapter is hardware-agnostic. The techniques here --- paged attention, speculative decoding, continuous batching, MoE serving, fused sampling kernels, mixed-precision stability --- apply on H100s, MI300Xs, TPUs, and whatever comes next. Chapters 1--3 benchmark these on specific hardware. Here we build the conceptual and quantitative toolkit for understanding \textit{why} each technique exists and when it helps.
\end{whythis}

% -----------------------------------------------------------------------
\subsection{KV Cache Management}
% -----------------------------------------------------------------------

Every autoregressive generation step recomputes attention over all previous tokens unless the key-value projections are cached. KV caching is the dominant technique for making generation efficient: each new token attends to cached KV tensors from previous positions rather than recomputing from scratch. The cost is memory: for a model with $L$ layers, $H$ attention heads, head dimension $d$, batch size $B$, and sequence length $S$, the KV cache occupies:

$$\text{KV memory} = 2 \times L \times H \times d \times S \times B \times \text{bytes\_per\_element}$$

The factor of 2 is for keys and values. At FP16 (2 bytes), a 70B model with 80 layers, 64 heads, head dimension 128, batch 32, context 4096 tokens occupies roughly 54 GB --- most of an H100's 80 GB just for the cache.

\textbf{Paged attention.} The naive approach allocates a contiguous memory block per sequence at the maximum possible length. On variable-length sequences (the RLVR case, where completions range from 100 to 2000+ tokens), this wastes HBM proportionally to the gap between actual and maximum length. vLLM's paged attention instead manages KV cache as fixed-size \textit{blocks} (analogous to OS virtual memory pages), allocated on demand and mapped to non-contiguous physical memory. Fragmentation drops from $O(S_\text{max})$ to $O(\text{block size})$ per sequence. For RLVR with high length variance, this directly increases the number of concurrent rollouts that fit in memory.

\textbf{KV cache quantization.} FP8 or INT8 quantization of the KV cache reduces memory by 2--4$\times$ at the cost of quantization noise in attention scores. For supervised training, this noise is generally tolerable. For RLVR, the generation quality directly determines the reward signal, which directly determines the gradient. A question worth taking seriously: is KV quantization noise tolerable for RL? The answer is empirical and workload-dependent. Short-horizon tasks with binary rewards (correct/incorrect) may be insensitive to small attention noise. Long chain-of-thought generation, where a subtle early error compounds over 2000 tokens, may see meaningful reward degradation. Measure before deploying.

\textbf{Eviction and compression for long-context rollouts.} When the KV cache exceeds available memory, three strategies exist:
\begin{itemize}
  \item \textbf{Eviction:} Drop KV entries for positions that are least recently used or received low attention weight. Simple and fast but permanently loses context. Appropriate when the task does not require long-range dependencies.
  \item \textbf{Compression:} Retain ``attention sink'' tokens (the first few tokens, which accumulate disproportionate attention weight) plus a sliding recent window. StreamingLLM demonstrated that this maintains generation quality on long documents far beyond training context length.
  \item \textbf{Offloading:} Spill KV cache to CPU DRAM (typically 512--2048 GB on a modern server). Bandwidth-limited: PCIe 5.0 delivers around 64 GB/s, versus HBM3's 3.35 TB/s. Offloading is viable only when the CPU-side data is accessed infrequently --- for example, prompt KV that was computed once and is shared across all $G$ completions in a GRPO group.
\end{itemize}

\textbf{MoE memory interaction.} Dario Amodei confirmed Anthropic is training MoE models. The memory budget per GPU has three major competing consumers: (1) dense model weights (embedding, attention projections, layer norms), (2) expert weights (each expert is a separate FFN; with $E$ experts and $k$ active per token, only $k$ are used per forward pass but all $E$ must reside somewhere), and (3) the KV cache. On a sharded MoE deployment, expert weights may be sharded across GPUs, so each GPU holds $E / \text{num\_gpus}$ experts. But the KV cache is local --- each GPU caches the keys and values for the sequences it is processing. The three-way tradeoff is unavoidable: larger experts, longer context, or larger batch --- not all three.

\begin{thinkfirst}{MoE Memory Budget (no computer)}
You have a 70B MoE model with 8 experts (2 active per token) on 8 GPUs with 80 GB each. Draw the memory budget per GPU: dense model weights, expert weights (sharded across GPUs), KV cache, activations, and optimizer states.

Work through the numbers. Assume the 70B parameter count is split roughly 30B in dense components (embeddings, attention) and 40B in expert FFN weights across all 8 experts. Each expert holds 5B parameters; sharded across 8 GPUs, each GPU holds $5\text{B} / 8 \approx 0.625\text{B}$ expert parameters --- but it also needs all 8 experts' weights for routing, so expert weights per GPU are the full $8 \times 0.625\text{B} = 5\text{B}$ if not sharded, or $0.625\text{B}$ per GPU if expert-parallelism is used. Clarify which model parallelism strategy you are assuming.

At what context length does the GPU OOM? Show the arithmetic for FP16 weights and FP16 KV cache at batch size 16, varying context from 1K to 16K tokens.

Now redo the same analysis for an MI300X with 192 GB HBM. The extra 112 GB per GPU --- where does it go? Is it better invested in more context (longer rollouts without eviction), larger batch size (more parallel completions, higher GPU utilization), or in fitting a larger model without tensor parallelism overhead?

There is no single right answer. The point is to make explicit which constraint is binding and how extra memory relaxes it.
\end{thinkfirst}

% -----------------------------------------------------------------------
\subsection{Speculative Decoding for RLVR Rollouts}
% -----------------------------------------------------------------------

Speculative decoding is an exact-or-approximate method for accelerating autoregressive generation using a small \textit{draft model} that proposes $K$ tokens in parallel, followed by a single forward pass of the large \textit{target model} that verifies all $K$ proposals simultaneously. The key insight: a forward pass through the large model at sequence length $K+1$ costs roughly the same as a forward pass at sequence length 1 (for large models, the dominant cost is parameter loading from HBM, not flops) --- so if even one draft token is accepted, the effective throughput increases.

\textbf{Acceptance rate and breakeven.} Let $\alpha \in [0,1]$ be the per-token acceptance rate (the probability that the draft model's proposed token matches what the target would have sampled). The expected number of tokens produced per target forward pass is:

$$\mathbb{E}[\text{tokens per pass}] = K\alpha + 1$$

The $+1$ accounts for the one token the target always produces (either accepting the draft's final token or generating a correction). Standard autoregressive produces exactly 1 token per target forward pass. Speculative decoding is beneficial when:

$$K\alpha + 1 > 1 \implies \alpha > 0$$

But this ignores the cost of draft generation. Let the draft model cost $c_d$ per token and the target model cost $c_t$ per token (with $c_d \ll c_t$). The cost per $K$ draft tokens is $K c_d$. The cost of one target verification pass is $c_t$. Total cost to produce $K\alpha + 1$ tokens: $K c_d + c_t$. Standard autoregressive cost for $K\alpha + 1$ tokens: $(K\alpha + 1) c_t$. Speculative decoding wins when:

$$K c_d + c_t < (K\alpha + 1) c_t \implies K c_d < K\alpha \cdot c_t \implies \frac{c_d}{c_t} < \alpha$$

The breakeven acceptance rate is $\alpha^* = c_d / c_t$. If the draft model is 10$\times$ cheaper than the target, speculative decoding wins for $\alpha > 0.1$.

\textbf{The RLVR problem: a moving target.} In RLVR, the target model IS the policy being trained. Its distribution shifts every gradient update. The draft model, typically distilled from the target at initialization, approximates target step 0. At step 1000, the policy may have diverged substantially --- especially in reasoning models, where the chain-of-thought distribution changes as the model learns to solve harder problems. Acceptance rate $\alpha$ degrades monotonically with policy divergence (in general). At some training step $t^*$, $\alpha$ falls below the breakeven threshold and speculative decoding stops helping --- or actively hurts throughput due to the overhead of draft generation and verification.

\textbf{Re-distillation frequency.} The practical question is: how often must you re-distill the draft model to maintain $\alpha > \alpha^*$? This depends on the KL divergence between draft and target distributions, which in turn depends on the reward signal magnitude and learning rate. There is no closed-form answer; it must be measured experimentally. The exercise in this chapter builds that measurement.

\textbf{An elegant approach.} The previous RLVR checkpoint is itself an approximation of the current policy --- and re-distillation is free, since you already have the checkpoint. Using checkpoint $t-k$ as the draft model for checkpoint $t$ gives speculative decoding whose acceptance rate depends on how much the policy shifts in $k$ steps. This is the self-speculative decoding regime for RLVR. The open research question: at what checkpoint gap $k$ does acceptance rate fall below breakeven? Can you adapt $k$ dynamically based on measured KL divergence?

\begin{thinkfirst}{Speculative Decoding Breakeven Math}
Work through the following carefully before continuing.

\textbf{Part 1: Basic breakeven.} Draft model has acceptance rate $\alpha$ per token. Each draft step proposes $K$ tokens. The expected number of tokens produced per target forward pass is $K\alpha + 1$ (derive this from first principles using the geometric distribution of how many consecutive tokens are accepted before the first rejection). At what $\alpha$ and $K$ does speculative decoding beat standard autoregressive, assuming the draft model is free? Now add the draft cost: if the draft model is $1/r$ the cost of the target, what is the minimum $\alpha$ as a function of $K$ and $r$?

\textbf{Part 2: Distributional shift.} The target model is being RLVR-trained and its distribution shifts every update. The draft model was distilled from the target at step 0. Consider two extreme cases:

(a) The target learns to prefer a completely different vocabulary for chain-of-thought (e.g., it switches from English reasoning to symbolic reasoning). What happens to $\alpha$? What is the minimum $\alpha$ in the worst case?

(b) The target shifts gradually --- each update moves the distribution by KL divergence $\delta$. If $\alpha$ degrades linearly with cumulative KL divergence (a simplifying assumption), at what training step does speculative decoding stop helping?

\textbf{Part 3: Re-distillation strategy.} The draft model was distilled from the target at step 0 and costs $C_\text{distill}$ to re-distill. Speculative decoding provides a throughput gain of $g(\alpha)$ tokens/second. Re-distillation is worthwhile when the throughput gain integrated over the next $\Delta t$ steps exceeds $C_\text{distill}$. Derive the optimal re-distillation interval as a function of $\alpha$'s decay rate, $g(\cdot)$, and $C_\text{distill}$.
\end{thinkfirst}

% -----------------------------------------------------------------------
\subsection{Continuous Batching Under RLVR}
% -----------------------------------------------------------------------

RLVR rollout generation has an extreme length distribution. A single GRPO training step generates $G$ completions per prompt (typically $G = 8$--$32$) for each of $B$ prompts in the batch, producing $G \times B$ concurrent sequences. For reasoning-capable models solving math or coding problems, sequence lengths follow a heavy-tailed distribution: many completions finish quickly (100--300 tokens for easy problems), while hard problems produce extended chains of thought (1000--4000 tokens). The ratio of $\ell_\text{max} / \ell_\text{mean}$ can exceed 10:1.

\textbf{Static batching waste.} The naive strategy pads all sequences to the maximum length in the batch and runs them to completion simultaneously. For a batch where 90\% of sequences are 200 tokens and 10\% are 2000 tokens, static batching computes 10$\times$ as many tokens as needed for the short sequences --- and 90\% of GPU cycles in the tail are wasted on padding. At GRPO's scale, this is not a minor inefficiency; it can dominate GPU utilization metrics.

\textbf{Continuous/dynamic batching.} The fix is to insert new sequences into the running batch as old ones finish, rather than waiting for the entire batch to complete. This is the core technique in vLLM, TGI, and similar serving systems. The GPU is kept busy with real work rather than padding operations. Throughput increases proportionally to how much time was previously wasted.

\textbf{The RLVR synchronization barrier.} Continuous batching across independent requests is straightforward. RLVR adds a constraint: all $G$ completions for a single prompt must complete before advantage computation can proceed. The advantage $\hat{A}_i$ for completion $i$ is computed relative to the mean reward across all $G$ completions for that prompt:

$$\hat{A}_i = \frac{r_i - \mu_G}{\sigma_G}$$

where $\mu_G$ and $\sigma_G$ are the mean and standard deviation of rewards within the group. You cannot compute $\mu_G$ until all $G$ completions are finished. This creates a \textit{prompt group} abstraction: a group of $G$ sequences that must be synchronized before the gradient step.

\textbf{Batching strategy.} The recommended approach is to batch \textit{across} prompt groups while maintaining per-group synchronization. Concretely:
\begin{enumerate}
  \item Assign each incoming prompt a group ID. Launch all $G$ completions for that prompt simultaneously.
  \item Maintain a per-group completion counter. When a completion finishes, increment the counter.
  \item When a group's counter reaches $G$, mark the group as ready for advantage computation.
  \item Free the memory for all $G$ sequences in that group and immediately schedule a new prompt group.
\end{enumerate}

This gives continuous batching across groups (memory is recycled promptly) while preserving the intra-group synchronization needed for GRPO. The key insight: you do not need to wait for all groups to finish --- only for each individual group. Long-running groups can coexist with short groups in the same batch.

\begin{thinkfirst}{GRPO Batching Scenario}
GRPO with $G = 16$ completions per prompt, batch of 32 prompts. That is 512 concurrent generations. RLVR models produce extended chain-of-thought (100--2000 tokens). Some finish in 100 tokens, some in 2000.

\textbf{Scenario 1: Static batching.} All 512 sequences are padded to 2000 tokens. Compute the number of wasted token-forward-passes if the true length distribution is uniform on $[100, 2000]$. Express this as a fraction of total compute.

\textbf{Scenario 2: Dynamic batching without group awareness.} New sequences are inserted as old ones finish, ignoring prompt group membership. A group with 15/16 completions done but one slow outlier (2000 tokens) is stranded --- advantage computation cannot proceed. How do you handle this? What are the tradeoffs between waiting for the outlier versus dropping it?

\textbf{Scenario 3: Group-aware batching.} Implement the group abstraction described above. When does the synchronization barrier dominate? At $G = 4$ versus $G = 64$, how does the distribution of group completion times change (assume i.i.d. completion lengths within a group)? At what $G$ does the expected wait for the slowest completion in a group become the binding constraint on throughput?

Design a batching strategy that handles the RLVR synchronization constraint without significant padding waste. Identify the scheduling decisions: which groups to run concurrently, how to handle stragglers, how to bound memory usage.
\end{thinkfirst}

% -----------------------------------------------------------------------
\subsection{MoE Serving Patterns}
% -----------------------------------------------------------------------

Mixture-of-Experts models replace the dense FFN in each transformer layer with $E$ parallel expert FFNs, routing each token to the top-$k$ experts via a learned gating function. For generation, this introduces communication overhead that does not exist in dense models.

\textbf{All-to-all per token.} At each transformer layer during generation, the routing decision must be made and tokens dispatched to the appropriate expert. In a distributed MoE setting where experts are sharded across GPUs (expert parallelism), each token's computation requires an all-to-all communication: send token representations from their current GPU to the GPU holding the selected expert, compute the expert FFN, and send results back. For a single token step, this is two all-to-all operations per MoE layer. For a model with $L/2$ MoE layers (alternating dense and MoE, a common design), that is $L$ all-to-all operations per autoregressive step. At generation batch sizes of hundreds of sequences, this communication can dominate step latency.

\textbf{Expert caching.} In practice, token routing is not uniform. A given model tends to send similar tokens to similar experts --- the routing function is deterministic given the hidden state. For a fixed prompt distribution, a small number of ``hot'' experts handle a disproportionate share of tokens. Expert caching exploits this: keep the hot experts in faster memory (HBM or L2 cache), evict cold experts. At inference time with a static prompt distribution, expert caching can reduce effective all-to-all volume by 2--5$\times$. Under RLVR, the prompt distribution shifts as the policy learns, so hot experts shift over time --- caching must adapt.

\textbf{Load balancing under RLVR.} MoE models use an auxiliary load-balancing loss to prevent expert collapse (where all tokens route to a single expert, wasting $E-1$ experts). Under RLVR training, the routing patterns change as the policy update shifts the hidden state distribution. This creates dynamic load imbalance: at some training steps, the load-balancing loss has not yet corrected for the new routing distribution, and some GPUs holding cold experts are idle while others are saturated. Monitoring per-expert load variance during RLVR training is a useful diagnostic for whether expert capacity is a throughput bottleneck.

% -----------------------------------------------------------------------
\subsection{Sampling Kernels}
% -----------------------------------------------------------------------

At each generation step, after the forward pass computes logits over the vocabulary, sampling requires: (1) divide logits by temperature, (2) compute softmax, (3) apply top-$k$ or top-$p$ filtering, (4) sample from the resulting distribution. These steps are often implemented as separate GPU kernel launches, each requiring a round-trip through HBM. For large vocabulary sizes (32K--256K tokens), this overhead is measurable.

\textbf{Fused sampling kernel.} A single fused kernel that performs temperature scaling, softmax, top-$k$ filtering, and multinomial sampling in one pass eliminates intermediate HBM reads and writes. The speedup depends on vocabulary size: for small vocabularies (32K), the softmax computation is fast and HBM bandwidth is not the bottleneck. For large vocabularies (256K, as in modern multilingual models), the softmax over the full vocabulary requires reading and writing 512 KB per sequence per step (at FP16), and kernel fusion provides 10--30\% latency reduction on this step. As a fraction of total generation time for a 7B model, sampling is typically 5--15\% --- not the dominant cost, but not negligible across weeks of training.

\textbf{Temperature scheduling during RLVR.} The sampling temperature controls exploration versus exploitation during rollout generation. A natural schedule is high temperature early in training (when the policy is uncertain and exploration is valuable) ramping down to low temperature as the policy converges (to reduce variance in the reward signal). This is analogous to $\epsilon$-annealing in DQN. Implementing temperature as a dynamic parameter in the sampling kernel (rather than a compile-time constant) is a systems detail that matters for RLVR workflows where temperature changes every few hundred gradient steps.

\textbf{Rejection sampling and best-of-$N$.} A common alternative to temperature-based exploration is best-of-$N$: generate $N$ completions, score them all with the reward model, and keep the highest-scoring one for training (or for final output at inference time). The naive implementation generates all $N$ completions before scoring any. An early-exit optimization scores completions as they finish and terminates generation for remaining sequences once a sufficiently high-scoring sample is found --- reducing average generation length at the cost of some suboptimality (the best completion might have been among the ones terminated early). For RLVR training where you need the \textit{distribution} of completions (not just the best), best-of-$N$ with early exit is less applicable, but it matters for evaluation and deployment.

% -----------------------------------------------------------------------
\subsection{Numerical Stability in Mixed Precision (Dario's Factors 6--7)}
% -----------------------------------------------------------------------

Dario Amodei's ``Big Blob'' framework for frontier model training identifies numerical stability --- ``laminar flow of gradients'' --- as a first-class engineering concern, not an afterthought. Factors 6 and 7 in his framework emphasize that precision choices affect training stability in ways that interact with the optimizer, the architecture, and the loss signal.

\textbf{The precision stack.} Modern RLVR training operates at multiple precision levels simultaneously:
\begin{itemize}
  \item \textbf{Weights:} BF16 or FP8, depending on hardware support and stability requirements. FP8 weights halve memory versus BF16 at the cost of a smaller dynamic range.
  \item \textbf{Activations:} BF16 for most operations. FP32 for numerically sensitive operations (softmax, layer norm, loss computation).
  \item \textbf{KV cache:} As discussed above, FP8 or INT8 reduces memory. The precision here affects generation quality, not the backward pass.
  \item \textbf{Gradient accumulation:} FP32 is standard. Accumulating gradients in BF16 introduces rounding errors that compound over micro-batches. At large gradient accumulation counts (common in RLVR where the effective batch size must be large for stable advantage estimates), BF16 gradient accumulation can cause silent training instabilities.
  \item \textbf{Optimizer states:} FP32 for Adam first and second moments. This is the largest memory cost in the optimizer --- twice the model size in FP32. Mixed-precision optimizers (e.g., bf16 master weights with fp32 moments) reduce this.
\end{itemize}

\textbf{RL gradients are different.} In supervised learning, gradient noise from quantization is small relative to the signal --- the loss function is well-conditioned. Policy gradient estimates are inherently high-variance: the gradient is $\nabla_\theta \log \pi_\theta(a|s) \cdot \hat{A}$, and $\hat{A}$ can be large and variable. Quantization noise in FP8 activations adds an additional stochastic perturbation on top of this already-noisy estimator. The concern is not that FP8 always breaks RL training --- it may be fine in practice. The concern is that the failure mode is silent: training appears to proceed, rewards increase, but at a slower rate or with higher variance than BF16 would achieve. Without a BF16 baseline to compare against, you would not know. The recommendation is to run short ablations comparing FP8 and BF16 activation precision early in a new RLVR setup, before committing to a weeks-long run.

\textbf{Softmax overflow in long-context generation.} At context lengths of 4000+ tokens, attention logits can overflow FP16 (max value $\approx 65504$). The standard fix is to compute attention in FP32 (or use Flash Attention's online softmax, which avoids materializing the full logit matrix). Overflow in the attention softmax causes the generation distribution to collapse to the argmax --- equivalent to zero-temperature sampling --- which destroys exploration in RLVR. This is a correctness issue, not just a precision tradeoff.

% -----------------------------------------------------------------------
\subsection{Exercises}
% -----------------------------------------------------------------------

\begin{exercise}{KV Cache Budget Calculator (2--3 hours)}
Build a Python tool that computes memory budgets for transformer (including MoE) models.

\textbf{Inputs:}
\begin{itemize}
  \item Model config: number of parameters, number of layers, hidden dimension, number of heads, head dimension
  \item MoE config: number of experts, experts per token (active experts)
  \item Precision config: dtype for weights (FP16/BF16/FP8), KV cache dtype (FP16/FP8/INT8), activation dtype
  \item Hardware: GPU memory in GB, number of GPUs (for tensor/expert parallelism)
\end{itemize}

\textbf{Outputs:}
\begin{itemize}
  \item Memory breakdown: dense weights, expert weights (per GPU given expert parallelism), KV cache at given batch/context, activations, optimizer states
  \item Maximum context length given a fixed batch size (solve for $S$ such that total memory $\leq$ GPU memory)
  \item Maximum batch size given a fixed context length
  \item Utilization \% (what fraction of GPU memory is used by each component)
\end{itemize}

\textbf{Required runs:} H100 SXM5 (80 GB), MI300X (192 GB), TPU v5e (16 GB HBM per chip). For the TPU, use model parallelism across 4 chips.

\textbf{Deliverable:} A table comparing all three hardware targets for a 7B dense model and a 70B MoE model (8 experts, 2 active), at batch size 16 and context 4096, with FP16 weights and FP16 KV. Annotate which component is the binding constraint on each platform.
\end{exercise}

\begin{exercise}{Speculative Decoding Prototype (4--6 hours)}
Implement speculative decoding for GPT-2 small (117M, draft) + GPT-2 medium (345M, target) using HuggingFace. GPT-2 small is approximately $1/3$ the compute cost of GPT-2 medium --- use this as a proxy for the draft/target cost ratio.

\textbf{Part 1: Baseline measurement.}
\begin{enumerate}
  \item Implement standard autoregressive generation for GPT-2 medium. Measure tokens/second.
  \item Implement speculative decoding with draft length $K \in \{1, 2, 4, 8, 16\}$. For each $K$, measure: (a) empirical acceptance rate $\alpha$, (b) tokens/second, (c) speedup ratio vs autoregressive.
  \item Plot speedup vs $K$. Identify the $K$ that maximizes speedup. Does it match the theoretical prediction from Part 1 of the Think First box?
\end{enumerate}

\textbf{Part 2: Distributional shift under RLVR.}
\begin{enumerate}
  \item RLVR-train GPT-2 medium for 200 PPO steps on a simple reward (e.g., reward for generating text that contains a specific word or matches a pattern). Keep GPT-2 small fixed as the draft model.
  \item Every 20 training steps, measure the acceptance rate $\alpha$ between the current GPT-2 medium and the fixed GPT-2 small draft.
  \item Plot $\alpha$ vs training step. Find $t^*$: the step at which $\alpha$ falls below the breakeven threshold. Compute the breakeven threshold from the cost ratio.
  \item Repeat with the draft model updated every 50 steps (re-distillation via copying weights + fine-tuning). Does periodic update maintain $\alpha$ above breakeven?
\end{enumerate}

\textbf{Deliverable:} The degradation curve $\alpha(t)$ is the core artifact. Write a paragraph interpreting it: at what training phase does speculative decoding help most, and what does this imply for production RLVR systems?
\end{exercise}

\begin{exercise}{Variable-Length Batching Simulator (3--4 hours)}
No GPU required --- pure Python simulation. This exercise builds intuition for how batching strategy interacts with the GRPO synchronization constraint.

\textbf{Setup.} Simulate 1000 RLVR rollouts. Length distribution: log-normal with $\mu = 6$, $\sigma = 1$ (mean $\approx 403$ tokens, with heavy right tail extending to 2000+ tokens). Group size $G = 16$. Total prompts: 1000~/$G$ = 62 groups, each with 16 completions.

\textbf{Simulate three strategies:}
\begin{enumerate}
  \item \textbf{Static batching:} Pad all sequences in a batch to the maximum length. Batch size fixed at 64 sequences. Measure: total token-steps computed, wasted token-steps (padding), effective GPU utilization (real tokens / total tokens), wall time (proportional to total token-steps), peak memory (batch size $\times$ max length).
  \item \textbf{Dynamic batching per prompt group:} Each group of $G = 16$ completions is processed together, but sequences within the group are not padded to global max --- only to the group's max. New groups start when all 16 completions in the current group finish. Measure same metrics.
  \item \textbf{Fully continuous batching across groups:} Sequences are inserted and removed individually as they finish. The only constraint: advantage computation for a group is triggered when the group's $G$-th completion finishes. Measure same metrics.
\end{enumerate}

\textbf{Analysis:}
\begin{itemize}
  \item Which strategy wins on each metric? Is there a single winner, or is there a Pareto frontier?
  \item At what group size $G \in \{4, 8, 16, 32, 64\}$ does the synchronization barrier dominate throughput? Plot expected group completion time (= max of $G$ i.i.d. samples from the log-normal) vs $G$.
  \item What is the 95th percentile completion time for a group at $G = 16$? What does this imply for tail latency in production?
\end{itemize}

\textbf{Deliverable:} A table of utilization \%, wall time, and memory peak for all three strategies at $G \in \{8, 16, 32\}$. A one-paragraph capacity planning recommendation for an RL team choosing a batching strategy.
\end{exercise}

\begin{exercise}{Sampling Kernel Fusion (2--3 hours)}
Implement and benchmark fused vs.\ unfused sampling.

\textbf{Part 1: Implementation.}
\begin{enumerate}
  \item \textbf{Unfused baseline:} In PyTorch, implement the sampling pipeline as separate operations: \texttt{logits / temperature} $\to$ \texttt{torch.softmax} $\to$ top-$k$ masking $\to$ \texttt{torch.multinomial}. Time each operation separately.
  \item \textbf{Fused with \texttt{torch.compile}:} Wrap the entire pipeline in \texttt{torch.compile()} with \texttt{mode="reduce-overhead"}. Measure end-to-end latency.
  \item \textbf{Custom kernel (stretch):} Implement a Triton kernel that performs all four steps in a single pass. Compare to \texttt{torch.compile} performance.
\end{enumerate}

\textbf{Part 2: Scaling analysis.} Measure latency at vocab sizes $V \in \{32768, 50257, 65536, 131072, 256000\}$, batch sizes $\{1, 16, 64, 256\}$. For each $(V, \text{batch size})$ pair, compute:
\begin{itemize}
  \item Absolute latency (ms) for fused and unfused
  \item Speedup ratio = unfused / fused
  \item Theoretical roofline: is this operation memory-bandwidth-bound or compute-bound?
\end{itemize}

\textbf{Part 3: Generation time fraction.} For a 7B model generating at batch size 64, estimate what fraction of total per-step wall time is spent on sampling. Use your latency measurements to compute this. At what model size does sampling become negligible ($< 1\%$ of step time)?

\textbf{Deliverable:} A heatmap of speedup ratio across $(V, \text{batch size})$. Annotation of the crossover point where fusion provides $> 10\%$ speedup. A calculation of the sampling fraction for representative model sizes.
\end{exercise}

\begin{portfolio}{RL Inference Efficiency Guide}
Write a technical blog post with real measurements titled ``RL Inference Efficiency: KV Cache, Speculative Decoding, and Batching for GRPO-Scale Training.''

\textbf{Required sections:}
\begin{enumerate}
  \item \textbf{KV Cache Management:} Present your memory budget calculator results. Show the per-platform breakdown table. Explain which component is the binding constraint on each hardware target and what practitioners should do about it.
  \item \textbf{Speculative Decoding Under Distributional Shift:} Present the $\alpha(t)$ degradation curve. This is the monkey --- the measurement that requires building something and interpreting a result. Describe: (a) the speedup at training step 0, (b) the crossover point $t^*$ where speculative decoding stops helping, (c) the practical implication for production RLVR systems (when to use spec dec, when to abandon it, whether periodic re-distillation is worth the cost).
  \item \textbf{Batching Strategies for GRPO Rollouts:} Present your simulator results. Give the actionable capacity planning recommendation: for a team running GRPO with $G = 16$ on a cluster of H100s, what batching strategy should they use, and why?
\end{enumerate}

\textbf{The monkey.} ``How does speculative decoding's acceptance rate degrade when the target model is being RLVR-trained, and what is the practical implication for training throughput?'' Building the prototype (Exercise 2) is the pedestal. The degradation curve and the concrete recommendation (``use spec dec for the first $N$ gradient steps, then switch to standard autoregressive, re-distill every $M$ steps if you want to extend the benefit'') is the monkey. The blog post is not complete without a number for $N$ and a recommendation for $M$ grounded in your measurements.

\textbf{Verification checklist:}
\begin{enumerate}
  \item Speculative decoding shows measurable throughput speedup at training step 0 (at least $1.3\times$ for $K = 4$, $\alpha \geq 0.5$ on natural language).
  \item The degradation curve $\alpha(t)$ has a concrete crossover point $t^*$ that you can identify from your measurements.
  \item The batching simulator produces a utilization table that an RL team could use for capacity planning (numbers are in the right ballpark for real RLVR workloads).
  \item The blog post is written for a technical audience (ML engineers) and contains enough implementation detail that a reader could reproduce your measurements.
\end{enumerate}
\end{portfolio}

\newpage

\section{Distributed Training \& Communication}

\smallskip\noindent\textit{Study Guide companion: Study Guide Section 12 (Hardware Deep Dive) covers the interview-ready version of this material.}
\smallskip

\begin{whythis}
A single GPU at 90\% utilization means nothing if 1000 GPUs spend 40\% of their time on communication.
At frontier scale---millions of accelerators across TPU, Trainium, and GPU clusters---the relevant metric
is not per-device FLOP/s but cluster-level intelligence-per-Watt.
Every wasted cycle waiting for a gradient synchronization or an expert routing token is a wasted Joule
multiplied by a million.
Distributed systems knowledge is therefore not a peripheral skill for a research engineer at this scale;
it is the primary lever for translating compute budget into capability gains.
RLVR training compounds this: unlike pretraining, it interleaves inference and training, produces
variable-length outputs that create load imbalance, and---for MoE architectures---relies on all-to-all
communication patterns that scale far worse than the all-reduce used in dense pretraining.
Understanding where the bubbles are, and which ones are worth eliminating, is the job.
\end{whythis}

% ----------------------------------------------------------------
\subsection*{Conceptual Foundation}
% ----------------------------------------------------------------

\begin{thinkfirst}{Parallelism Strategies: Memory Math for a 70B Model}
Before reading further, work through the following from first principles.
You have a 70B-parameter model (assume all weights in \texttt{bfloat16}) and 8 GPUs each with 80 GB HBM.

\medskip
\textbf{Data Parallel (DDP).}
Every GPU holds a full copy of the model.
70B parameters $\times$ 2 bytes/param = 140 GB per GPU.
Each GPU also holds gradients (140 GB) and, with Adam, two optimizer state tensors (280 GB).
Total per GPU: $\sim$560 GB.
\emph{Does it fit?} No --- 560 GB $\gg$ 80 GB.
DDP alone is ruled out for a 70B model on 80 GB GPUs.

\medskip
\textbf{Tensor Parallelism (TP).}
Split each weight matrix along one dimension across $T$ GPUs.
An attention projection of shape $[d, d_{\text{head}} \cdot h]$ is split column-wise; each GPU holds
$1/T$ of the columns.
With $T = 8$, parameter memory per GPU: $140/8 = 17.5$ GB --- fits comfortably.
Cost: each forward pass requires an \texttt{all-reduce} after every tensor-parallel layer
(to combine partial sums).
This is a \emph{synchronous, latency-sensitive} communication on the critical path.

\medskip
\textbf{Pipeline Parallelism (PP).}
Assign contiguous layer groups to different GPUs (stages).
With 8 stages and 80 layers, each GPU holds 10 layers $\approx 17.5$ GB --- fits.
Cost: \emph{pipeline bubbles}.
With $M$ micro-batches and $P$ stages, the bubble fraction is $(P-1)/(M+P-1)$.
For $P=8$ and $M=8$: bubble fraction $= 7/15 \approx 47\%$ --- nearly half the step is idle.
Increasing $M$ reduces the bubble but increases memory for activations.

\medskip
\textbf{Fully Sharded Data Parallel (FSDP).}
Shard \emph{all} tensors (parameters, gradients, optimizer states) across all $N$ GPUs.
Memory per GPU: $\text{total} / N$.
For 70B + Adam: $\sim$560 GB / 8 = 70 GB --- barely fits with careful activation checkpointing.
Protocol: (1) \texttt{all-gather} the full parameter tensor for each layer before its forward pass,
(2) compute forward, (3) compute backward using the gathered weights,
(4) \texttt{reduce-scatter} gradients so each GPU receives its shard's gradient,
(5) update local shard.
The gather and reduce-scatter double the communication volume relative to DDP's all-reduce, but
overlap with computation when implemented carefully.

\medskip
\textbf{Hybrid strategies.}
Production systems combine all three.
A typical recipe: TP within a node (NVLink), PP across nodes (InfiniBand), FSDP across all replicas.
The reason TP stays intra-node is that tensor-parallel all-reduces are on the critical path of every
forward pass --- they must complete before the next layer starts.
NVLink bandwidth ($\sim$900 GB/s bidirectional on H100 NVLink4) can absorb this; InfiniBand cannot.

\medskip
\textbf{Design question:} For the 70B model on 8$\times$80 GB GPUs, sketch a strategy that fits in
memory and minimizes compute bubbles.
What does your communication graph look like?
Where is the bottleneck?
\end{thinkfirst}

\begin{thinkfirst}{Why RLVR Distribution Is Harder Than Pretraining}
RLVR training is not a single workload --- it is a pipeline of at least three distinct
compute phases that must be orchestrated across the same (or overlapping) hardware.
Before reading the analysis, ask yourself: which phase is memory-bandwidth-bound and which is
compute-bound?
Why does this distinction matter for how you schedule work across GPUs?

\medskip
\textbf{Phase 1 --- Rollout generation (inference).}
The policy generates $G$ completions per prompt (GRPO uses $G \sim 8$--$16$).
Generation is autoregressive: each token step requires loading the full KV cache plus the model
weights, performing a small matrix-vector multiply, and sampling.
This is \emph{memory-bandwidth-bound}, not compute-bound.
A single A100 has $\sim$2 TB/s HBM bandwidth; with a 7B model at bfloat16, loading weights takes
$\sim$7 ns per byte but the flop count per token is tiny relative to that bandwidth.
GPU utilization during generation: typically 30--60\%, even with batching.

\medskip
\textbf{Phase 2 --- Reward evaluation.}
Each completion must be scored by a reward model (or verified by a verifier in RLVR math settings).
This is a separate forward pass.
If the reward model differs in architecture from the policy, it may need its own GPU partition or
must time-share.
Scheduling these evaluations without stalling the generation pipeline requires careful batching.

\medskip
\textbf{Phase 3 --- Policy update (training).}
Once all $G$ completions for a batch of prompts are collected, advantages are computed and a gradient
step is taken.
This is \emph{compute-bound}: large matrix multiplications in the backward pass.
GPU utilization should be high here --- but only if the batch is large enough and
fully padded, which leads to the next problem.

\medskip
\textbf{Variable-length outputs and load imbalance.}
Completions vary in length.
A prompt that terminates early leaves its GPU idle while others finish.
The training step cannot begin until \emph{all} $G$ completions per prompt are done
(needed for advantage normalization in GRPO).
This creates a synchronization barrier at the end of every generation round.
The slowest completion in the batch determines when training can start --- the classic straggler problem.

\medskip
\textbf{MoE expert parallelism under RLVR.}
Dense models use all-reduce: every device sends/receives the same aggregated gradient.
MoE models route each token to $k$ of $E$ experts; experts are sharded across GPUs.
This requires \texttt{all-to-all} communication: device $i$ sends a \emph{different} chunk to
every other device.
During pretraining, routing distributions are approximately stationary.
During RLVR, the policy changes at every step, and routing patterns shift dynamically ---
some experts become more popular, others are underutilized.
This creates \emph{dynamic} load imbalance on top of the static straggler problem.

\medskip
\textbf{Design question:} Sketch a timeline diagram for one RLVR step on 4 GPUs using GRPO with
$G=8$ and a 7B MoE policy.
Label each phase and mark the synchronization barriers.
Where are the bubbles?
Which ones can be overlapped and which cannot?
\end{thinkfirst}

\begin{thinkfirst}{Communication Patterns: All-Reduce vs.\ All-to-All}
Before reading the derivations, estimate: for 8 GPUs exchanging 1 GB of gradient data,
how long does a ring all-reduce take if each GPU-to-GPU link has 100 GB/s bandwidth?
Then read below and check your estimate.

\medskip
\textbf{Ring all-reduce.}
Arrange $N$ devices in a ring.
Each device starts with a tensor of size $S$ bytes.
The ring all-reduce proceeds in two phases, each with $N-1$ steps.

\emph{Reduce-scatter phase:} Each device sends a chunk of size $S/N$ to the next device in the ring,
receives a chunk from the previous device, and adds it to its local buffer.
After $N-1$ steps, each device holds one fully reduced chunk of the global sum.

\emph{All-gather phase:} Each device sends its reduced chunk to the next device in the ring.
After $N-1$ steps, every device holds all $N$ reduced chunks, i.e., the full sum.

Total data sent per device: $2 \cdot (N-1)/N \cdot S$ bytes.
As $N \to \infty$, this approaches $2S$ bytes --- essentially twice the tensor size, regardless of $N$.
Bandwidth utilization: $(N-1)/N$.
For $N=8$, utilization $= 7/8 = 87.5\%$.
Latency: $2(N-1) \cdot \alpha + 2(N-1)/N \cdot S / \beta$,
where $\alpha$ is per-hop latency and $\beta$ is per-link bandwidth.
Ring all-reduce is bandwidth-optimal but has $O(N)$ latency --- poor for large $N$ at small message sizes.

\medskip
\textbf{Tree (recursive halving/doubling) all-reduce.}
Devices are paired, then groups of 4, then 8, etc.
Number of steps: $\log_2 N$.
Latency: $O(\log N \cdot \alpha)$ --- much better for small messages.
But each step doubles the message size being transmitted.
Bandwidth utilization is lower because early steps transmit small messages inefficiently.
Tree all-reduce wins when $\alpha$ dominates (small tensors, high-latency links).
Ring all-reduce wins when $\beta$ dominates (large tensors, high-bandwidth links like NVLink).

\medskip
\textbf{All-to-all (MoE routing).}
Every device sends a \emph{different} chunk to every other device.
Total data sent per device: $(N-1) \cdot c$ bytes, where $c$ is the chunk size per destination.
Total cluster communication: $N \cdot (N-1) \cdot c$ bytes --- $O(N^2)$.
Compare to ring all-reduce: $O(N \cdot S)$ bytes cluster-total, which for uniform sharding is
$O(S)$ per device regardless of $N$.
All-to-all scales catastrophically with $N$ for fixed per-device capacity.
For MoE with $E$ experts sharded across $N$ GPUs, each forward pass requires two all-to-all calls
(dispatch tokens to experts, receive results back).
Under RLVR, routing is non-uniform and changes every step --- some all-to-all messages are large,
some are tiny, and the imbalance stresses the network fabric in unpredictable ways.

\medskip
\textbf{Hardware implications.}
NVIDIA NVLink provides all-to-all bandwidth within a node that is roughly symmetric.
AMD Infinity Fabric (used with RCCL, the ROCm equivalent of NCCL) has a different topology:
bandwidth between two specific GPUs depends on which Infinity Fabric links connect them,
which may be asymmetric.
This matters for all-to-all: if expert 3 is always sending to GPU 7 via a slower link,
that link becomes a bottleneck every step.
RCCL's all-to-all implementation may need topology-aware routing that NCCL does not,
because NVLink's switched topology is more uniform.
This is a concrete hardware-software co-design question relevant to AMD evaluation (Chapter 1).

\medskip
\textbf{Design question:} For a 16-GPU cluster running MoE with 64 experts, compute the all-to-all
communication volume per step if each token's expert routing vector is 1 KB and you process 4096
tokens per step.
Compare this to the gradient all-reduce volume for a 7B dense model.
At what batch size does the all-to-all dominate?
\end{thinkfirst}

% ----------------------------------------------------------------
\subsection*{Exercises}
% ----------------------------------------------------------------

\begin{exercise}{Ring All-Reduce from Scratch (3--4 hours)}
Implement ring all-reduce using raw Python multiprocessing on a single machine.
Do not use \texttt{torch.distributed.all\_reduce} for the implementation itself --- use it only as a
reference for correctness checking.

\medskip
\textbf{Setup.}
Spawn 4 processes using \texttt{torch.multiprocessing} or Python's \texttt{multiprocessing}.
Each process represents one ``device'' and holds a local tensor of size $M$ bytes.
Use shared memory or Unix sockets for inter-process communication.

\medskip
\textbf{Step 1 --- Reduce-scatter.}
Divide each tensor into 4 chunks.
Implement the ring: process $i$ sends chunk $i$ to process $(i+1) \bmod 4$,
receives from process $(i-1) \bmod 4$, and accumulates.
Repeat for $N-1 = 3$ steps, rotating which chunk is sent.
After this phase, process $i$ holds the globally summed version of chunk $i$.

\medskip
\textbf{Step 2 --- All-gather.}
Process $i$ sends its reduced chunk to process $(i+1) \bmod 4$.
After $N-1$ steps, every process holds the full summed tensor.
Verify correctness against \texttt{torch.distributed.all\_reduce}.

\medskip
\textbf{Step 3 --- Benchmarking.}
Measure achieved bandwidth vs.\ theoretical for message sizes:
1 KB, 10 KB, 100 KB, 1 MB, 10 MB, 100 MB, 1 GB.
Plot bandwidth (GB/s) vs.\ message size (log scale).
You should see two distinct regimes:
\begin{itemize}
  \item \emph{Latency-dominated} (small messages): bandwidth is far below peak because
        setup and synchronization overhead dominate.
  \item \emph{Bandwidth-dominated} (large messages): throughput approaches the
        theoretical limit of your inter-process transport.
\end{itemize}
Identify the crossover point in bytes.
Compute theoretical bandwidth for your transport and compare.

\medskip
\textbf{Step 4 --- NCCL comparison.}
Repeat the bandwidth plot using \texttt{torch.distributed.all\_reduce} with NCCL backend on actual
GPUs (or \texttt{gloo} backend if only CPU is available).
At what message size does NCCL reach $>$80\% of the theoretical NVLink bandwidth?
How does NCCL's crossover point compare to your implementation's?

\medskip
\textbf{Deliverables.} Bandwidth-vs-message-size plot with both curves.
Annotated crossover point.
One paragraph explaining the shape of each curve in terms of $\alpha$--$\beta$ cost model.
\end{exercise}

\begin{exercise}{NCCL Profiling on Multi-GPU Instance (2--3 hours)}
Rent a multi-GPU instance (2$\times$H100 $\approx$ \$5/hr, or 4$\times$A100 $\approx$ \$6/hr).
Budget approximately \$15--20 for this exercise.

\medskip
\textbf{Benchmark setup.}
Use \texttt{torch.distributed} with NCCL backend.
Write a script that launches all-reduce for message sizes:
1 MB, 10 MB, 100 MB, 1 GB.
For each size: run 10 warmup iterations, then 50 timed iterations.
Record mean and standard deviation of latency.
Compute achieved bandwidth as $2(N-1)/N \cdot S / \text{latency}$.

\medskip
\textbf{Profiling tools.}
Enable \texttt{NCCL\_DEBUG=INFO} and redirect stderr to a log file.
Parse the log to find: which NCCL algorithm is selected for each message size (ring, tree, or
collnet), and what the reported bus bandwidth is.
Optionally, run under \texttt{nsys profile} to get a GPU timeline.
In the timeline, identify: NCCL kernel launches, data copy operations, and synchronization points.

\medskip
\textbf{Overhead decomposition.}
For the 1 MB case, what fraction of the all-reduce time is:
\begin{enumerate}
  \item Kernel launch and setup overhead?
  \item Actual data transfer?
  \item Synchronization/barrier?
\end{enumerate}
How does this fraction change at 1 GB?
At what message size does NCCL switch algorithms?
Does the switch correspond to the crossover you found in the previous exercise?

\medskip
\textbf{Scaling.}
If your instance has 4 GPUs, run all-reduce on 2, 3, and 4 GPUs for a fixed 100 MB tensor.
Does bandwidth scale as $(N-1)/N$ predicts?
If not, why not?

\medskip
\textbf{Deliverables.} Bandwidth-vs-message-size table with algorithm annotations.
Overhead decomposition bar chart for 1 MB and 1 GB.
Written answer: at what message size does NCCL reach $>$80\% of peak NVLink bandwidth?
\end{exercise}

\begin{exercise}{Toy FSDP: Manual Sharding of GPT-2 Small (4--6 hours)}
Implement FSDP manually for GPT-2 small (117M parameters) on 2 GPUs.
Do not use PyTorch's built-in \texttt{FullyShardedDataParallel} for the core implementation;
use it only for comparison.

\medskip
\textbf{Step 1 --- Shard parameters.}
For each parameter tensor $\theta \in \mathbb{R}^{d}$, GPU $i$ stores
$\theta_i = \theta[i \cdot d/2 : (i+1) \cdot d/2]$.
Do this for all parameters: weights, biases, embedding tables.
Verify that memory per GPU is approximately half of full-model memory.

\medskip
\textbf{Step 2 --- All-gather before forward.}
Before each layer's forward pass, issue \texttt{dist.all\_gather} to reconstruct the full parameter
tensor on both GPUs.
Perform the forward computation.
Immediately free the gathered tensor to reclaim memory (critical: if you keep all gathered tensors,
you have replicated the full model and defeated the purpose).

\medskip
\textbf{Step 3 --- Backward and reduce-scatter.}
PyTorch's autograd will compute gradients with respect to the gathered tensor.
After the backward pass, issue \texttt{dist.reduce\_scatter} to distribute gradient shards:
GPU $i$ accumulates the gradient for its parameter shard.
Each GPU now has the gradient for its shard only.

\medskip
\textbf{Step 4 --- Local optimizer update.}
Each GPU runs its optimizer (Adam) on its local parameter shard using its local gradient shard.
Optimizer state (momentum, variance) is also sharded and never needs to be gathered.
This is the memory saving: optimizer state is $2\times$ parameter memory for Adam,
and it is always sharded.

\medskip
\textbf{Step 5 --- Measurement.}
Compare peak memory per GPU for: (a) full replication (DDP), (b) your manual FSDP,
(c) PyTorch FSDP.
Measure throughput (samples/sec) for a fixed batch.
Where is your implementation slower than PyTorch FSDP?
Profile to find: is the bottleneck in the all-gather, the forward pass, or the reduce-scatter?

\medskip
\textbf{Deliverables.}
Memory comparison table.
Throughput comparison.
Profiler output annotated with where time is spent relative to PyTorch FSDP.
One paragraph: what does PyTorch FSDP do that your implementation does not, based on the profiler evidence?
\end{exercise}

\begin{exercise}{RLVR Distribution: Profiling Generation Bubbles (4--6 hours)}
Take a minimal RLVR training loop (nanoRLVR from Chapter 8, or build a stripped-down GRPO loop
around a 1B--3B model) and instrument it for distribution profiling on 2 GPUs.

\medskip
\textbf{Setup.}
Split rollout generation across 2 GPUs (each GPU generates $G/2$ completions per prompt).
Both GPUs participate in the training step.
Use a dataset of math problems (e.g., GSM8K) with verifier-based rewards so reward evaluation is
deterministic and fast.

\medskip
\textbf{Instrumentation.}
Wrap each phase with \texttt{torch.profiler} or manual \texttt{time.perf\_counter} checkpoints:
\begin{itemize}
  \item $t_{\text{gen\_start}}$, $t_{\text{gen\_end}}$: generation phase per GPU
  \item $t_{\text{reward\_start}}$, $t_{\text{reward\_end}}$: reward evaluation
  \item $t_{\text{sync}}$: time waiting at the synchronization barrier (straggler cost)
  \item $t_{\text{train\_start}}$, $t_{\text{train\_end}}$: backward pass
  \item $t_{\text{comm}}$: gradient communication time
\end{itemize}

\medskip
\textbf{Bubble identification.}
For a batch of 16 prompts with $G=8$:
compute the fraction of total step time spent in each phase.
The useful compute fraction is $t_{\text{train}} / t_{\text{step}}$.
Everything else is overhead: generation (unavoidable but schedulable), reward (schedulable),
sync wait (the bubble), and gradient communication (reducible with overlap).
What fraction of the step is the synchronization wait (straggler bubble)?

\medskip
\textbf{Optimization.}
Propose and implement one optimization targeting the largest bubble you found.
Options:
\begin{itemize}
  \item \emph{Async generation:} start the next round of generation while the training backward pass
        runs on completed rollouts.
  \item \emph{Staggered scheduling:} assign shorter prompts to one GPU and longer prompts to the
        other to balance generation time.
  \item \emph{Early cutoff:} truncate the slowest completions once a timeout is reached
        (analyze the reward distribution impact).
\end{itemize}

\medskip
\textbf{Measurement.}
Before and after your optimization: report step time, useful compute fraction,
and synchronization wait fraction.
Target: reduce the synchronization bubble by $>$30\%.

\medskip
\textbf{Deliverables.}
Flame graph (torch profiler output) showing step phases and communication stalls.
Before/after utilization bar chart.
One paragraph: which bubble was largest, which optimization you chose, and why.
\end{exercise}

% ----------------------------------------------------------------
\subsection*{Intelligence-per-Watt at Cluster Scale}
% ----------------------------------------------------------------

For $N$ GPUs, define the useful compute fraction $\eta(N)$ as the fraction of aggregate FLOP/s
that produces gradients contributing to model improvement.
The rest is overhead: communication, synchronization barriers, and idle bubbles.

For \textbf{pretraining with all-reduce} (DDP or FSDP), $\eta(N)$ declines slowly with $N$
because ring all-reduce volume is $O(S)$ per device regardless of $N$.
The bottleneck is inter-node bandwidth; $\eta$ is typically $>$70\% even at thousands of GPUs
when network is not the bottleneck.

For \textbf{RLVR with GRPO synchronization}, $\eta(N)$ declines faster due to the straggler
barrier at the end of every generation round.
As $N$ grows, the probability that at least one GPU has a slow completion increases,
and the expected straggler delay grows as $O(\log N)$ for i.i.d.\ exponential completion times.
At large $N$, generation bubbles can consume $>$40\% of step time.

For \textbf{MoE with all-to-all routing}, $\eta(N)$ declines most steeply.
All-to-all volume scales as $O(N)$ per device (each device sends to all others), so communication
time grows linearly with $N$ while compute per device stays fixed.
There exists a critical $N^*$ beyond which adding more GPUs reduces $\eta$ faster than it adds
parallel throughput, making additional GPUs net-negative in terms of intelligence-per-Watt.

This is the scaling ceiling.
Identifying $N^*$ for each workload --- and building the measurement infrastructure to detect when
you are approaching it --- is a core research engineering responsibility at frontier scale.

% ----------------------------------------------------------------
\subsection*{Portfolio Project}
% ----------------------------------------------------------------

\begin{portfolio}{RLVR Training Distribution Profiler}
\textbf{The question.}
What fraction of RLVR training compute is wasted on distribution overhead, and which specific
overhead is most cost-effective to eliminate?

\medskip
\textbf{The monkey.}
Quantify the largest utilization gap in an RLVR training loop
(generation bubble, communication stall, GRPO synchronization wait, or MoE routing imbalance)
and implement one optimization that reduces it by a measurable amount.

\medskip
\textbf{The pedestal.}
A profiling harness that instruments a real RLVR training loop (minimum 1B-parameter policy,
real reward signal) and produces per-step breakdowns of:
generation time, reward evaluation time, synchronization wait time, gradient communication time,
and backward pass time.
The harness must produce flame graphs and utilization charts automatically.

\medskip
\textbf{Implementation plan.}
\begin{enumerate}
  \item Run the RLVR loop uninstrumented for 100 steps to establish baseline step time and throughput.
  \item Add \texttt{torch.profiler} instrumentation. Run 20 steps and export the profiler trace.
        Open the trace in Chrome \texttt{chrome://tracing} or TensorBoard profiler.
        Identify the top-3 time sinks per step.
  \item For the largest bubble ($>$10\% of step time), design one targeted optimization.
        Write a one-paragraph justification before implementing.
  \item Implement the optimization. Run 20 instrumented steps.
        Compare: step time, useful compute fraction ($\eta$), and peak GPU memory.
  \item Write the blog post.
        Required figures: (a) pre-optimization flame graph with bubble annotated,
        (b) post-optimization flame graph, (c) per-phase utilization bar chart before/after,
        (d) step time distribution histogram showing variance reduction from straggler mitigation.
\end{enumerate}

\medskip
\textbf{Verification criteria.}
\begin{enumerate}
  \item Flame graph shows at least one identifiable bubble $>$10\% of step time.
  \item Proposed optimization reduces that bubble by $>$30\% in measured wall time.
  \item Blog post is written for an ML engineer audience with sufficient implementation detail
        to reproduce the measurement setup and the optimization.
\end{enumerate}

\medskip
\textbf{Hardware budget.}
2$\times$H100 at $\approx$\$5/hr or 4$\times$A100 at $\approx$\$6/hr.
Profiling runs are short; budget $\approx$\$50 for this project within the \$200/month allocation.
\end{portfolio}

\newpage

\section{Policy Gradient --- Hypothesis-Driven Depth}

\smallskip\noindent\textit{Study Guide companion: Study Guide Section 2 (RL for LLMs) covers the conversational, interview-ready version of this material.}
\smallskip

\begin{whythis}
There is a deep mathematical insight about what policy gradient is \textit{truly} doing --- an unlock that changes how you think about algorithm design for RLVR. This chapter does not guess the answer. Guessing is a trap: you pick your favorite theory, confirm it, and stop looking.

Instead, this chapter builds the more durable skill: \textit{generating hypotheses, deriving their predictions, and designing experiments that could falsify them}. Six candidate ``unlocks'' are developed here, each with a concrete prediction and a minimal experiment. Some will be wrong. The point is the discipline.

The right posture is not ``I know the secret.'' It is: ``I have been attacking this problem systematically. Here are the hypotheses I have generated, the predictions I have derived, and the experiments I have designed to test them. Here is where the evidence points --- and here is what would change my mind.'' That is the research taste Anthropic is hiring for.
\end{whythis}

% -----------------------------------------------------------------------
\subsection{Foundation: Policy Gradient from First Principles}
% -----------------------------------------------------------------------

\begin{thinkfirst}{Policy Gradient Theorem --- Two Derivations}
Work both derivations on paper before reading further. The goal is to understand not just the formula but which assumptions each derivation requires --- because those assumptions determine where each estimator breaks down.

\textbf{Derivation 1: Score Function Estimator (Log-Derivative Trick).}

The objective is $J(\theta) = \mathbb{E}_{\tau \sim \pi_\theta}[R(\tau)]$, where $\tau = (s_0, a_0, s_1, a_1, \ldots)$ is a trajectory. Expand the expectation as an integral over trajectory distributions:
\begin{align*}
  J(\theta) &= \int p_\theta(\tau) R(\tau) \, d\tau
\end{align*}

Differentiate under the integral sign (valid when $p_\theta(\tau)$ is differentiable in $\theta$ and the reward is bounded):
\begin{align*}
  \nabla_\theta J(\theta) &= \int \nabla_\theta p_\theta(\tau) \cdot R(\tau) \, d\tau
\end{align*}

Apply the log-derivative identity $\nabla_\theta p_\theta(\tau) = p_\theta(\tau) \nabla_\theta \log p_\theta(\tau)$:
\begin{align*}
  \nabla_\theta J(\theta) &= \int p_\theta(\tau) \nabla_\theta \log p_\theta(\tau) \cdot R(\tau) \, d\tau = \mathbb{E}_{\pi_\theta}\!\left[\nabla_\theta \log \pi_\theta(\tau) \cdot R(\tau)\right]
\end{align*}

Since $\log p_\theta(\tau) = \sum_t \log \pi_\theta(a_t | s_t) + \text{const}$ (the environment dynamics cancel):
\begin{align*}
  \nabla_\theta J(\theta) &= \mathbb{E}_{\pi_\theta}\!\left[\sum_t \nabla_\theta \log \pi_\theta(a_t|s_t) \cdot R(\tau)\right]
\end{align*}

Key property: this estimator is unbiased. It requires only the ability to \textit{sample} from $\pi_\theta$ and evaluate $\log \pi_\theta(a|s)$ --- the reward function does not need to be differentiable. This makes it applicable to discrete action spaces.

\textbf{Derivation 2: Reparameterization Trick (Continuous Actions).}

When $\pi_\theta(a|s)$ is a distribution with a differentiable sampling procedure, write $a = f_\theta(s, \epsilon)$ where $\epsilon \sim p(\epsilon)$ is a noise source independent of $\theta$. Then:
\begin{align*}
  J(\theta) &= \mathbb{E}_{\epsilon \sim p(\epsilon)}\!\left[R(f_\theta(s, \epsilon))\right]
\end{align*}

Now $\theta$ enters inside the function, so:
\begin{align*}
  \nabla_\theta J(\theta) &= \mathbb{E}_{\epsilon}\!\left[\nabla_\theta R(f_\theta(s, \epsilon))\right] = \mathbb{E}_{\epsilon}\!\left[\frac{\partial R}{\partial a} \cdot \frac{\partial f_\theta}{\partial \theta}\right]
\end{align*}

The gradient flows directly through the sample. This has lower variance than the score function estimator because it uses derivative information about how the return changes with the action, rather than only weighting samples by their return.

\textbf{Variance comparison.} The score function estimator's variance is $O(R^2)$ --- it scales with the squared magnitude of the return signal. For high-return trajectories in RLVR (long reasoning traces), this variance can be enormous. The reparameterization estimator's variance depends on $\|\partial R / \partial a\|$ --- the sensitivity of the reward to action perturbations --- which can be much smaller. However, reparameterization requires: (1) continuous actions, (2) a differentiable reward function $R(a)$, (3) a differentiable sampling function $f_\theta$. In RLVR with discrete token generation and deterministic verifier rewards, none of these hold. You are stuck with the score function estimator. Understanding \textit{why} is the insight.
\end{thinkfirst}

\begin{thinkfirst}{The Optimal Baseline}
The score function estimator has high variance. A standard trick is to subtract a \textit{baseline} $b(s)$ from the return before forming the gradient estimate. Prove that this is unbiased, then derive the optimal baseline.

\textbf{Unbiasedness proof.} For any state-dependent function $b(s)$, show that:
\begin{align*}
  \mathbb{E}_{\pi_\theta}\!\left[\nabla_\theta \log \pi_\theta(a|s) \cdot b(s)\right] &= 0
\end{align*}

Proof: for a single state-action pair,
\begin{align*}
  \mathbb{E}_{a \sim \pi_\theta(\cdot|s)}\!\left[\nabla_\theta \log \pi_\theta(a|s) \cdot b(s)\right]
  &= b(s) \int \pi_\theta(a|s) \nabla_\theta \log \pi_\theta(a|s) \, da \\
  &= b(s) \int \nabla_\theta \pi_\theta(a|s) \, da \\
  &= b(s) \cdot \nabla_\theta \int \pi_\theta(a|s) \, da \\
  &= b(s) \cdot \nabla_\theta 1 = 0
\end{align*}

The key step is exchanging the gradient and integral (valid under mild regularity) and using $\int \pi_\theta(a|s) da = 1$.

\textbf{Variance-minimizing baseline.} The gradient estimator with baseline is $g(a) = \nabla_\theta \log \pi_\theta(a|s) \cdot (R - b(s))$. The variance is $\mathbb{E}[g^2] - (\mathbb{E}[g])^2$. Since $\mathbb{E}[g]$ is constant in $b$ (unbiasedness), minimizing variance means minimizing $\mathbb{E}[g^2]$. Let $h = \nabla_\theta \log \pi_\theta(a|s)$ and minimize over $b$:
\begin{align*}
  \frac{d}{db} \mathbb{E}\!\left[\|h\|^2 (R - b)^2\right] &= 0 \\
  -2 \mathbb{E}\!\left[\|h\|^2 (R - b)\right] &= 0 \\
  b^*(s) &= \frac{\mathbb{E}\!\left[\|\nabla_\theta \log \pi_\theta\|^2 \cdot R\right]}{\mathbb{E}\!\left[\|\nabla_\theta \log \pi_\theta\|^2\right]}
\end{align*}

This is the optimal baseline: the gradient-magnitude-weighted average return. \textit{Not} the value function $V(s) = \mathbb{E}[R | s]$. The value function is the unweighted average return --- it approximates $b^*$ only when $\|\nabla_\theta \log \pi_\theta\|^2$ is approximately constant across actions. This holds roughly when the policy is near-uniform. As training progresses and the policy sharpens, $\|\nabla_\theta \log \pi_\theta\|^2$ becomes highly non-uniform across actions (confident actions have near-zero gradient; uncertain actions have large gradient), and $V(s)$ becomes a worse approximation to $b^*$. GRPO uses group-normalized returns as a proxy --- question whether this is better or worse than $V(s)$ as an approximation to $b^*$.
\end{thinkfirst}

\begin{thinkfirst}{Natural Gradient and Distribution Geometry}
The ordinary gradient $\nabla_\theta J$ is computed in parameter space. But parameter space is not the right geometry for policy optimization --- moving $0.01$ in one parameter direction might change the policy distribution by a large amount (e.g., shifting a weight that directly gates a high-probability token), while $0.01$ in another direction changes it negligibly (e.g., shifting a weight deep in a saturated layer). Optimizing in parameter space mixes these effects arbitrarily.

\textbf{The Fisher information matrix.} Define:
\begin{align*}
  F(\theta) &= \mathbb{E}_{a \sim \pi_\theta(\cdot|s)}\!\left[\nabla_\theta \log \pi_\theta(a|s) \cdot \nabla_\theta \log \pi_\theta(a|s)^\top\right]
\end{align*}

$F$ measures the curvature of the policy distribution with respect to parameter changes. Large $F_{ii}$ means that changing $\theta_i$ by a small amount causes a large change in the distribution of $\pi_\theta$.

\textbf{Connection to KL divergence.} For small perturbations $\delta$:
\begin{align*}
  D_{\mathrm{KL}}(\pi_\theta \| \pi_{\theta + \delta}) &\approx \frac{1}{2} \delta^\top F(\theta) \delta
\end{align*}

This follows from a Taylor expansion of the KL divergence. It means $F$ is the local metric tensor for the manifold of policy distributions --- the natural geometry is Riemannian with metric $F$, not Euclidean with metric $I$.

\textbf{Natural gradient.} The steepest ascent direction under the constraint of a fixed KL change $\delta^\top F \delta = c$ is:
\begin{align*}
  \tilde{\nabla}_\theta J &= F(\theta)^{-1} \nabla_\theta J
\end{align*}

The natural gradient normalizes the ordinary gradient by the local distribution curvature. It is invariant to reparameterizations of $\theta$ --- if you change the parameterization of your policy network, the natural gradient update in distribution space is the same.

\textbf{Why this matters for RLVR.} In high dimensions (billions of parameters), $F$ is never computed exactly. TRPO approximates the constraint $\delta^\top F \delta \leq \epsilon$ via conjugate gradient. PPO approximates the natural gradient via clipping. K-FAC and Shampoo approximate $F^{-1}$ via block-diagonal structure. Each of these is a different approximation to the same underlying geometry. Understanding this unifying picture lets you reason about which approximation breaks down first as model scale increases.
\end{thinkfirst}

\begin{thinkfirst}{PPO Clipping vs.\ Trust Regions}
PPO and TRPO both aim to limit how much the policy changes per update. They do so differently, and the differences matter in the RLVR regime.

\textbf{TRPO: hard KL constraint.} TRPO solves:
\begin{align*}
  \max_\theta \quad & \mathbb{E}\!\left[\frac{\pi_\theta(a|s)}{\pi_{\text{old}}(a|s)} \hat{A}(s,a)\right] \\
  \text{s.t.} \quad & D_{\mathrm{KL}}(\pi_{\text{old}} \| \pi_\theta) \leq \delta
\end{align*}

The constraint is a hard bound on distribution change. TRPO uses conjugate gradient + line search, making it expensive per step but theoretically grounded: the trust region is in distribution space, matching the natural geometry.

\textbf{PPO: clipping heuristic.} Let $\rho_t = \pi_\theta(a_t|s_t) / \pi_{\text{old}}(a_t|s_t)$ be the importance ratio. PPO maximizes:
\begin{align*}
  L^{\text{CLIP}}(\theta) &= \mathbb{E}_t\!\left[\min\!\left(\rho_t \hat{A}_t,\; \text{clip}(\rho_t, 1-\epsilon, 1+\epsilon) \hat{A}_t\right)\right]
\end{align*}

When $\hat{A}_t > 0$ (good action), the objective is capped at $(1+\epsilon)\hat{A}_t$ --- the policy cannot increase the action probability by more than a factor of $(1+\epsilon)$. When $\hat{A}_t < 0$ (bad action), the objective is capped at $(1-\epsilon)\hat{A}_t$ --- the policy cannot decrease the action probability by more than a factor of $(1-\epsilon)$.

\textbf{Non-equivalence.} TRPO's constraint is in distribution space (KL). PPO's clipping is in ratio space (per-action importance weights). These are not equivalent:

\begin{itemize}
  \item PPO can have small per-token ratios $\rho_t$ while accumulating a large joint KL over a long sequence --- the product $\prod_t \rho_t$ can differ dramatically from 1 even if each $\rho_t \in [1-\epsilon, 1+\epsilon]$.
  \item When advantage estimates $\hat{A}_t$ are wrong in sign, PPO clips in the wrong direction: it allows the policy to move against the true advantage. TRPO's constraint is symmetric and does not make this error.
  \item PPO's clipping does not account for the magnitude of the advantage. A large positive advantage with a large positive ratio is clipped the same way as a small positive advantage with a large positive ratio --- but the policy update risk is very different.
\end{itemize}

\textbf{RLVR regime: large batches and long sequences.} RLVR trains with long reasoning traces (hundreds to thousands of tokens) and large batch sizes (many rollouts per update). In this regime:
\begin{itemize}
  \item The joint KL divergence over a sequence accumulates across tokens: $D_{\mathrm{KL}}(\pi_{\text{old}}(\tau) \| \pi_\theta(\tau)) \approx \sum_t D_{\mathrm{KL}}(\pi_{\text{old}}(\cdot|s_t) \| \pi_\theta(\cdot|s_t))$. Even if each per-token KL is small, the sequence-level KL grows linearly in sequence length. PPO's per-token clipping does not control the sequence-level distribution change.
  \item GRPO (Group Relative Policy Optimization) uses a group-normalized advantage and adds an explicit KL penalty to the reward --- a different approach that puts the KL constraint in the reward signal rather than the optimization constraint. Deriving when this is better or worse than PPO clipping is a research question, not a settled fact.
\end{itemize}

Write down your answer: at sequence length 2048 with $\epsilon = 0.2$, what is the worst-case sequence-level probability ratio $\prod_t \rho_t$ under PPO clipping? (Answer: $(1.2)^{2048} \approx 10^{157}$.) Does this concern you? What would a tighter per-sequence constraint look like?
\end{thinkfirst}

% -----------------------------------------------------------------------
\subsection{Candidate Unlocks: Six Hypotheses}
% -----------------------------------------------------------------------

Research taste is not knowing the answer --- it is the ability to generate structured hypotheses and design experiments that distinguish between them. Rather than betting on one theory, the following six candidates are developed with equal rigor. The experiments in the next subsection test them. You should hold each hypothesis lightly: the right experiment could rule any of them out.

\subsubsection{Candidate 1: RL-as-Inference / Variational View}

\textbf{Hypothesis.} Policy gradient is variational inference. Define a target distribution over trajectories proportional to the exponentiated reward: $p^*(\tau) \propto \exp(R(\tau) / \alpha)$, where $\alpha > 0$ is a temperature. The policy optimization problem can be written as:
\begin{align*}
  \min_{\pi_\theta} \; D_{\mathrm{KL}}(\pi_\theta(\tau) \| p^*(\tau)) &= \min_{\pi_\theta} \; \mathbb{E}_{\pi_\theta}\!\left[\log \pi_\theta(\tau) - \frac{R(\tau)}{\alpha}\right] + \text{const}
\end{align*}
This is equivalent to maximizing $\mathbb{E}[R(\tau)/\alpha] - D_{\mathrm{KL}}(\pi_\theta \| p_{\text{ref}})$ for some reference policy. The KL penalty in GRPO (and RLHF) is the variational prior --- it regularizes $\pi_\theta$ toward the reference. The unlock: if you know the ELBO bound you are targeting, you can schedule $\alpha$ adaptively to improve convergence.

\textbf{Prediction.} Adaptive KL scheduling (dynamically adjusting $\alpha$ to target a specific ELBO gap) should improve GRPO convergence speed and final reward compared to a fixed KL coefficient, because fixed $\alpha$ is wasteful: early in training when the policy is far from optimal, a large KL penalty slows exploration; late in training when the policy is near-optimal, the same penalty unnecessarily constrains fine-grained optimization.

\textbf{Experiment.} Implement GRPO on TinyRL-Align (Chapter~7). Train three conditions: (a) fixed $\alpha = 0.01$, (b) fixed $\alpha = 0.1$, (c) adaptive $\alpha$ that adjusts to keep $D_{\mathrm{KL}}(\pi_\theta \| \pi_{\text{ref}}) \in [0.01, 0.05]$ at each step. Measure convergence speed (steps to 80\% of final reward) and final reward at 10K steps.

\subsubsection{Candidate 2: Natural Gradient / Distribution Geometry}

\textbf{Hypothesis.} PPO's clipping is a poor approximation to the natural gradient update, and the approximation degrades as model size increases. In high dimensions, the Fisher information matrix has highly non-uniform eigenvalues: most of the curvature is concentrated in a low-dimensional subspace corresponding to the most policy-sensitive parameter directions. PPO clips uniformly in ratio space, ignoring this structure.

\textbf{Prediction.} The cosine similarity between PPO's gradient direction and the true natural policy gradient should decrease as model size increases, because larger models have more extreme Fisher eigenvalue distributions (longer tails in the spectrum).

\textbf{Experiment.} Implement a small MLP policy (hidden dims 32, 64, 128, 256) on CartPole. At a fixed policy checkpoint, compute: (a) the ordinary policy gradient $g = \nabla_\theta J$, (b) the empirical Fisher $\hat{F} = \frac{1}{N} \sum_i h_i h_i^\top$ where $h_i = \nabla_\theta \log \pi_\theta(a_i|s_i)$, (c) the natural gradient $\tilde{g} = \hat{F}^{-1} g$ via conjugate gradient. Also compute the PPO gradient by running one PPO step from the same checkpoint. Measure $\cos(\tilde{g}, g_{\mathrm{PPO}})$ for each hidden dim. Plot cosine similarity vs.\ model size.

\subsubsection{Candidate 3: GRPO Advantage Estimation Bias}

\textbf{Hypothesis.} GRPO normalizes advantages within each group of $G$ rollouts from the same prompt: $\hat{A}_i = (R_i - \mu_{\text{group}}) / \sigma_{\text{group}}$. When reward variance differs systematically across prompts, this introduces a bias: low-variance prompts (where all rollouts get similar rewards) produce near-zero normalized advantages, so the policy barely updates on them. High-variance prompts produce large normalized advantages, so the policy updates heavily on them. If high-variance prompts are harder or less representative, GRPO over-specializes to them.

\textbf{Prediction.} After GRPO training, per-prompt policy change (measured by KL divergence from initialization) should be higher for high-variance prompts than low-variance prompts, even if both prompt types have the same expected reward.

\textbf{Experiment.} Create a prompt set with 50 prompts: 25 where the reward is $\mathrm{Bernoulli}(0.5) \in \{0, 1\}$ (high variance, $\sigma^2 = 0.25$), and 25 where the reward is $\mathcal{N}(0.5, 0.01)$ (low variance). Train with GRPO ($G = 8$) for 200 steps. After training, compute $D_{\mathrm{KL}}(\pi_{\text{final}}(\cdot|s) \| \pi_{\text{init}}(\cdot|s))$ for each prompt. Compare the mean KL across prompt types. Also compare to standard PPO trained on the same data.

\subsubsection{Candidate 4: Entropy as Information Bottleneck}

\textbf{Hypothesis.} Entropy collapse (the policy converging to near-deterministic behavior) is not merely a symptom of overfitting to the reward --- it reflects a sudden reduction in the effective dimensionality of the policy's logit representations. When entropy collapses, the effective rank of the output logit matrix (the matrix whose columns are logit vectors across tokens) drops sharply, indicating that the model has collapsed to a low-dimensional subspace of its representational capacity.

\textbf{Prediction.} The effective rank of the logit matrix (measured by $\exp(H(\sigma))$ where $H(\sigma)$ is the entropy of the normalized singular value spectrum) should drop at the same training step as entropy collapse, not before or after.

\textbf{Experiment.} Train a language model on TinyRL-Align with GRPO. At each checkpoint, compute: (a) mean policy entropy $H(\pi_\theta)$ over the validation set, (b) effective rank of the logit matrix (batch of 256 tokens, logits $\in \mathbb{R}^{256 \times V}$ where $V$ is vocabulary size, compute SVD and effective rank). Plot both metrics on the same axis over training. Compute Pearson correlation and the lag (in steps) between entropy drop and rank drop. If they are simultaneous, this supports the hypothesis.

\subsubsection{Candidate 5: Credit Assignment in Sequence-Level Rewards}

\textbf{Hypothesis.} Policy gradient with sequence-level rewards (one reward per full trajectory) applies the same reward signal to every token in the sequence, regardless of which token was causally responsible for the outcome. This means the policy gradient overcredits early tokens: they receive credit signals from more future trajectories (since early token choices determine the distribution of later states) and experience more gradient updates on average.

\textbf{Prediction.} Learning speed should be higher (faster convergence in fewer steps) when the critical decision --- the token choice that determines the reward --- is at the beginning of the sequence rather than the end.

\textbf{Experiment.} Construct a synthetic task: a sequence of $N$ binary choices, but the reward is determined entirely by choice at position $k$ (for $k \in \{1, N/4, N/2, 3N/4, N\}$). All other choices have no effect on reward. Train with REINFORCE (sequence-level reward) for each $k$. Measure convergence speed (steps to 90\% optimal policy at position $k$). Plot convergence speed vs.\ position $k$. Does it decrease monotonically? Does adding per-token advantage estimation (GAE) reduce the positional dependence?

\subsubsection{Candidate 6: RLVR Long-Horizon Dynamics}

\textbf{Hypothesis.} Over weeks-long RLVR training runs, the policy traverses qualitatively distinct regimes: (1) an exploration phase where entropy is high and gradient norms are large but reward is low, (2) an exploitation phase where the policy hones in on verifiable-reward strategies and reward climbs sharply, and (3) a verifier overfitting phase where the policy discovers reward-hacking strategies specific to the verifier and reward plateaus or drops on held-out distribution.

\textbf{Prediction.} Phase transitions should be identifiable as discontinuities (sharp changes in slope) in training metrics --- specifically entropy, gradient norm, reward variance, and held-out reward --- at specific compute thresholds. The transitions should correspond to capability jumps measurable on independent benchmarks.

\textbf{Experiment.} Train TinyRL-Code (Chapter~7) for 10$\times$ the standard number of steps. Log at every step: policy entropy, gradient $\ell_2$ norm, mean reward (training prompts), reward variance across rollouts, and held-out benchmark accuracy. Apply a change-point detection algorithm (e.g., PELT or binary segmentation) to each metric series. Do the detected change-points cluster at similar training steps across metrics? Do they correspond to jumps in held-out accuracy? If yes: what does this imply about when to stop training and how to schedule the KL coefficient?

% -----------------------------------------------------------------------
\subsection{Exercises}
% -----------------------------------------------------------------------

\begin{exercise}{Gradient Estimator Comparison (3--4 hours)}
Implement four policy gradient estimators on CartPole-v1 and empirically characterize their gradient variance. The goal is to connect the theoretical variance analysis from the Foundation section to measured reality.

\textbf{Environment setup.} Use \texttt{gymnasium} CartPole-v1. Policy: 2-layer MLP with hidden dim 64, softmax output. Initialize with a fixed random seed.

\textbf{Estimator 1: REINFORCE.} Roll out 1000 episodes from the current policy. For each episode, compute the return $G_t = \sum_{t'=t}^{T} \gamma^{t'} r_{t'}$ with $\gamma = 0.99$. Gradient estimate: $\hat{g} = \frac{1}{N}\sum_{i,t} \nabla_\theta \log \pi_\theta(a_{i,t}|s_{i,t}) \cdot G_{i,t}$.

\textbf{Estimator 2: REINFORCE with value baseline.} Train a separate value network $V_\phi(s)$ (same architecture, scalar output) with MSE loss on returns. Subtract: $\hat{A}_{i,t} = G_{i,t} - V_\phi(s_{i,t})$. Recompute gradient estimate using $\hat{A}_{i,t}$.

\textbf{Estimator 3: PPO.} Collect 1000 episodes. Compute advantages with GAE ($\lambda = 0.95$). Run 4 epochs of PPO updates with $\epsilon = 0.2$. For the gradient estimate: record the gradient at the first PPO gradient step (before any update).

\textbf{Estimator 4: Natural Policy Gradient.} Compute the empirical Fisher $\hat{F} = \frac{1}{N}\sum_i \nabla_\theta \log \pi_\theta(a_i|s_i) \nabla_\theta \log \pi_\theta(a_i|s_i)^\top$ using 500 samples. Compute $\tilde{g} = \hat{F}^{-1} g$ via conjugate gradient (10 iterations). This is exact (up to CG approximation) for the small CartPole policy.

\textbf{Measurements.} For each estimator, collect 20 independent gradient estimates (each from a fresh 1000-episode rollout). Compute: (a) mean gradient norm $\mathbb{E}[\|\hat{g}\|]$, (b) variance of gradient norm $\mathrm{Var}[\|\hat{g}\|]$, (c) mean cosine similarity between each estimate and the mean estimate (measures direction variance). Plot the distribution of gradient norms as a box plot. Which estimator has lowest variance? Does the ranking match the theoretical prediction from Think First 1?
\end{exercise}

\begin{exercise}{Exact Natural Gradient vs.\ PPO (2--3 hours)}
Test Candidate~2 directly: measure how closely PPO's gradient direction approximates the true natural gradient, and whether the approximation degrades with model size.

\textbf{Setup.} Train a CartPole policy with hidden dims $d \in \{32, 64, 128, 256\}$ for 100 steps using vanilla policy gradient (enough to get a non-trivial policy but not converge). Freeze the policy weights at this checkpoint for all comparisons below.

\textbf{Natural gradient computation.} Using 2000 rollouts from the frozen policy: compute the policy gradient $g = \nabla_\theta J$. Compute the empirical Fisher $\hat{F}$. Invert via conjugate gradient: $\tilde{g} = \hat{F}^{-1} g$ (use scipy's CG or implement from scratch). Normalize: $\tilde{g}_{\text{norm}} = \tilde{g} / \|\tilde{g}\|$.

\textbf{PPO gradient computation.} From the same frozen policy, collect 2000 rollouts. Compute GAE advantages. Run \textit{one} PPO gradient step (no parameter update --- just compute the gradient of the clipped objective and collect it). Normalize: $g_{\text{PPO,norm}} = g_{\text{PPO}} / \|g_{\text{PPO}}\|$.

\textbf{Measurement.} Compute $\cos(\tilde{g}_{\text{norm}}, g_{\text{PPO,norm}}) = \tilde{g}_{\text{norm}} \cdot g_{\text{PPO,norm}}$. Repeat for each hidden dim $d$. Plot cosine similarity vs.\ $d$. Also plot the ratio of the natural gradient norm to the ordinary gradient norm vs.\ $d$ --- this measures how much the curvature correction changes the update magnitude. Does cosine similarity decrease with $d$? If so, this is evidence for Candidate~2.
\end{exercise}

\begin{exercise}{GRPO Bias Test (2--3 hours)}
Test Candidate~3: does GRPO's group normalization create systematic bias toward high-variance prompts?

\textbf{Setup.} Use a small causal language model (GPT-2 small or equivalent, $\sim 117$M parameters). Construct a synthetic dataset with 50 prompts in two groups:

\begin{itemize}
  \item \textit{High-variance prompts} (25 prompts): the reward is determined by whether the model generates a specific rare token in the response. Each rollout independently has probability 0.5 of containing the target token ($\sigma^2_R = 0.25$).
  \item \textit{Low-variance prompts} (25 prompts): the reward is always $0.5 \pm 0.01$ (nearly constant, e.g., a reward model that gives near-constant scores for this prompt type). $\sigma^2_R \approx 10^{-4}$.
\end{itemize}

Both groups have the same expected reward ($\approx 0.5$).

\textbf{GRPO training.} Train with GRPO, $G = 8$ rollouts per prompt per step, for 200 gradient steps. Use a KL coefficient of $\beta = 0.01$.

\textbf{Measurements.} After 200 steps, for each of the 50 prompts, compute $D_{\mathrm{KL}}(\pi_{\text{final}}(\cdot|s) \| \pi_{\text{init}}(\cdot|s))$ by sampling 100 completions and estimating empirically. Report: mean KL for high-variance prompts, mean KL for low-variance prompts, ratio. Also compare total reward improvement (final minus initial reward) across both groups. Does GRPO update more on high-variance prompts? Is the KL ratio consistent with the variance ratio? For comparison, run the same experiment with standard PPO (which uses a shared value baseline rather than group normalization).
\end{exercise}

\begin{exercise}{Open-Ended: The Unlock (2 hours, pen and paper)}
Pick the candidate hypothesis you find most compelling after reading all six. Work from first principles.

\textbf{Part 1: Additional predictions.} Derive 3 additional falsifiable predictions beyond the one listed for your chosen candidate. Each prediction should: (a) follow logically from the hypothesis via a mathematical or mechanistic argument, (b) be testable with an experiment you could run in 2--4 hours, (c) be capable of falsifying the hypothesis (not just confirming it). For each prediction, describe the minimal experiment --- the smallest implementation that would produce a decisive result.

\textbf{Part 2: Adversarial case.} Design the experiment most likely to \textit{rule out} your chosen hypothesis. What data would you expect to see if the hypothesis is wrong? What would you conclude if you saw that data?

\textbf{Part 3: A seventh candidate.} The six candidates above are not exhaustive. Generate a seventh candidate unlock that is not on the list. It should: (a) be grounded in a specific mathematical property of policy gradient that is not captured by any of the six candidates, (b) have at least one falsifiable prediction, (c) be something you genuinely believe could be important. Write 1--2 paragraphs explaining the hypothesis and one paragraph explaining why you find it interesting.

\textbf{Hint for Part 3}: Consider properties of the trajectory distribution itself --- e.g., the geometry of the reward surface in token space, the relationship between the policy gradient and second-order optimization methods like Adam, or the interaction between the language model's pretraining distribution and the RL update direction.
\end{exercise}

% -----------------------------------------------------------------------
\subsection{Research Taste}
% -----------------------------------------------------------------------

\begin{taste}
If an interviewer asks ``What do you think policy gradient is really doing?'', the wrong answer is committing to one theory. Picking one framework and defending it signals that you stopped looking. The right answer demonstrates that you have been attacking the problem directly: ``I have been exploring this systematically. Here are six hypotheses I have been developing, here are the predictions I derived from each one, and here is what the experiments showed. The one I find most compelling right now is Candidate~X --- and here is why. But I hold it lightly. Here is the experiment that would change my mind.''

This is research taste: not getting attached to a given solution, but attacking the problem directly. Attachment to a framework is comfortable --- it gives you a ready answer. But it stops you from updating when the evidence points elsewhere, and in research, the most important experiments are the ones that rule out your favorite hypothesis.

The six candidates in this chapter are a starting point, not a complete list. The most important output of this chapter is not your answer to which candidate is right --- it is the skill of generating the next candidate when the current six are ruled out.
\end{taste}

\newpage

\section{TinyRL --- Community OSS Framework}

\smallskip\noindent\textit{Study Guide companion: Study Guide Section 8 (Research Proposals --- Proposal 2) covers the interview-ready version of this material.}
\smallskip

\begin{whythis}
RLVR (Reinforcement Learning with Verifiable Rewards) has replaced RLHF as the dominant paradigm for frontier model training. Models are trained for weeks against deterministic verifiers --- unit tests, math checkers, formal proof systems --- and the failure modes that emerge (reward hacking, entropy collapse, verifier overfitting) are poorly understood even by practitioners. Yet every paper studying these phenomena requires standing up a custom training stack, designing bespoke environments, and running experiments that take days on expensive hardware.

TinyRL fills the same role MNIST filled for computer vision: a standardized, minimal starting point that everyone can benchmark against. MNIST did not replace ImageNet. It gave researchers a shared substrate where ideas could be tested cheaply, compared cleanly, and built upon publicly. TinyRL does this for RLVR. The headline environment, TinyRL-Code, is RLVR in miniature --- verifiable rewards, deterministic unit tests, emergent strategies, and a reward hacking gallery --- all on a single GPU in under an hour.

As a community asset, TinyRL creates network effects. Contributors add environments, researchers compare results across algorithms, the reward hacking gallery attracts attention from practitioners who have seen the same exploits at scale. The barrier to entry for RLVR experimentation drops from ``access to a frontier training cluster'' to ``a laptop and 30 minutes.'' That is the lever.
\end{whythis}

\subsection{Design Principles}

The TinyStories paper (arXiv 2305.07759) showed that language model phenomena --- memorization, generalization, induction heads, grokking --- are observable at toy scale if the environment is designed carefully. Four properties made it work:

\begin{enumerate}
  \item \textbf{Controlled data distribution.} Stories were generated from a fixed vocabulary and grammar, so the input distribution was known and manipulable.
  \item \textbf{Tasks learnable by small models.} GPT-2 scale (124M parameters) was sufficient to exhibit the phenomena of interest.
  \item \textbf{Clean evaluation.} Human raters could assess story quality reliably; there was no ambiguity about what ``good'' meant.
  \item \textbf{Phenomena of interest observable at small scale.} Memorization curves, phase transitions, emergent capabilities --- all visible without scaling to billions of parameters.
\end{enumerate}

TinyRL maps each of these to RLVR:

\begin{enumerate}
  \item \textbf{Curated prompt sets.} A fixed bank of 50--100 programming puzzles (or IMDB prompts, or quantization tasks), sampled reproducibly. The training distribution is known.
  \item \textbf{GPT-2 scale models.} 124M--350M parameters. Cheap to train, cheap to share, easy to ablate.
  \item \textbf{Deterministic verification.} Unit tests pass or fail; sentiment scores are computed by a fixed classifier; inference cost is measurable. No human raters, no ambiguity.
  \item \textbf{RLVR phenomena visible in under one hour.} Reward hacking, entropy collapse, emergent chain-of-thought under verifiable pressure, KL divergence growth --- all observable on a single GPU before lunch.
\end{enumerate}

\begin{thinkfirst}{What Makes a Good Minimal RL Environment?}
Before reading further, design the simplest RLVR environment you can imagine that still exhibits \textit{verification hacking}. Specify: the model size, the task format, the verifier, and the exploit you predict the model will find. How would you know the model is gaming the verifier rather than solving the task? Write this down before continuing.
\end{thinkfirst}

\subsection{RLVR-First Framing}

RLVR trains a policy against rewards computed by an objective, external verifier --- a function that takes a completion and returns a scalar without human involvement. Math correctness is verified by a symbolic checker. Code correctness is verified by unit tests. The reward is not a neural proxy; it is a ground-truth signal about whether the model solved the problem.

This matters because RLHF's failure mode is reward model gaming: the policy learns to produce text that scores highly on the reward model without actually being better. RLVR's promise is that the verifier is non-gameable in principle --- if the unit tests pass, the code is correct. In practice, the verifier is gameable, and the exploits are more subtle: hardcoded outputs, format manipulation, specification gaps. TinyRL-Code is designed to make these exploits visible.

TinyRL has three environments. \textbf{TinyRL-Code} is the headline RLVR environment: unit-test-verified programming puzzles. \textbf{TinyRL-Align} is the RLHF baseline: sentiment optimization against a neural reward model. \textbf{TinyRL-Quant} is a RLVR-adjacent environment: reward is fully computable (accuracy $\times$ efficiency) but the action space is discrete and combinatorial. The contrast between Code and Align is itself the core pedagogical payload --- it shows concretely why the field moved from neural proxy rewards to verifiable rewards, and at what cost.

\subsection{Framework Architecture}

\subsubsection{Standardized Training Loop}

TinyRL exposes a single interface that every algorithm implements against:

\begin{itemize}
  \item \texttt{env.reset()} --- returns a prompt string sampled from the environment's prompt bank.
  \item \texttt{env.step(completion)} --- takes a string completion, runs the verifier, returns a scalar reward.
  \item \texttt{env.eval(policy)} --- runs the policy on the held-out evaluation set and returns environment-specific metrics.
\end{itemize}

Algorithms (PPO, GRPO, REINFORCE, DPO) are implemented as wrappers around this interface. Swapping algorithms requires changing one line: \texttt{trainer = PPOTrainer(env)} becomes \texttt{trainer = GRPOTrainer(env, G=8)}. The same environment definition works with every algorithm, enabling clean cross-algorithm comparisons.

\subsubsection{Metrics Dashboard}

Every training run automatically logs:

\begin{itemize}
  \item \textbf{Reward curve} --- mean reward per rollout batch, smoothed and raw.
  \item \textbf{KL divergence} --- $D_{\mathrm{KL}}(\pi_\theta \| \pi_{\mathrm{ref}})$ computed on each batch; the primary indicator of policy drift.
  \item \textbf{Policy entropy} --- $H(\pi_\theta)$; collapse signals mode collapse and reward overfitting.
  \item \textbf{Generation quality samples} --- 5 random completions logged at each checkpoint for qualitative inspection.
  \item \textbf{Gradient norm} --- $\ell_2$ norm of the gradient before clipping; instability here precedes training collapse.
\end{itemize}

All metrics are logged in a common schema (Weights \& Biases or TensorBoard), making cross-environment and cross-algorithm comparisons trivial. The leaderboard ingests the same metrics.

\subsubsection{Pluggable Environments}

Each environment is a single Python file defining three objects:

\begin{itemize}
  \item \texttt{prompts} --- a list of strings (or a callable that samples them).
  \item \texttt{reward\_fn(prompt, completion) -> float} --- the verifier.
  \item \texttt{eval\_fn(policy) -> dict} --- held-out evaluation returning a dict of named metrics.
\end{itemize}

New environments are contributed as pull requests. The review criterion is simple: does this environment exhibit a qualitatively distinct RL dynamic that existing environments do not cover? The goal is a curated zoo, not an uncurated dump.

\subsubsection{Leaderboards}

Each environment has a dedicated leaderboard tracking environment-specific competition metrics. Submissions include: algorithm name, hyperparameters, model checkpoint, and a reproducibility script. Results are verified by running the \texttt{eval\_fn} on a held-out test set that is not released publicly. Leaderboards reset on a quarterly cadence to encourage continued participation.

\subsection{TinyRL-Code: ``The MNIST of RLVR''}

\textbf{TinyRL-Code is the headline environment.} RLVR is the frontier paradigm. Code verification is the core RLVR domain at every major lab. If you build one environment, build this one.

\subsubsection*{Setup}

\begin{itemize}
  \item \textbf{Model:} GPT-2 small (124M) or CodeGen-350M, trained from a pretrained checkpoint.
  \item \textbf{Task bank:} 50--100 curated programming puzzles, graded by difficulty:
    \begin{itemize}
      \item \textit{Easy} --- implement a single function, verified by 1--3 unit tests (e.g., \texttt{reverse\_string}, \texttt{sum\_list}).
      \item \textit{Medium} --- implement 2 cooperating functions, verified by 5 unit tests (e.g., a parser and a formatter).
      \item \textit{Hard} --- implement an algorithm with edge cases, verified by 10 unit tests (e.g., BFS on an adjacency list, handling empty graph and disconnected nodes).
    \end{itemize}
  \item \textbf{Reward:} Fraction of unit tests passed. Pass all tests: reward = 1.0. Pass none: reward = 0.0. Partial credit for intermediate passes.
\end{itemize}

\subsubsection*{This Is RLVR in Miniature}

The verifier is objective, rule-based, and in principle non-gameable. In practice, models find exploits:

\begin{itemize}
  \item \textbf{Hardcoded outputs} --- the model learns the exact test inputs from repeated exposure and returns hardcoded outputs that pass the tests without implementing the function.
  \item \textbf{Format gaming} --- the model wraps a trivially wrong answer in \texttt{try/except} blocks to avoid exceptions, earning partial credit without solving the task.
  \item \textbf{Specification ambiguity} --- the model exploits underspecified edge cases in the problem statement that the test suite does not cover.
\end{itemize}

These are the same classes of exploit that frontier RLVR teams encounter at scale. TinyRL-Code makes them visible, characterizable, and patchable at toy scale --- exactly the verification engineering problem that matters for frontier training.

\subsubsection*{RLVR Dynamics Exhibited}

\begin{itemize}
  \item \textbf{Emergent problem-solving strategies} --- the model may develop chain-of-thought-like intermediate steps under verifiable reward pressure even without explicit CoT supervision.
  \item \textbf{Reward hacking on the verifier itself} --- as described above. This is the primary failure mode of RLVR.
  \item \textbf{Late-training capability emergence} --- harder puzzles may only become solvable after extended training, reproducing the ``late emergence'' phenomenon observed in frontier RLVR runs.
\end{itemize}

\subsubsection*{Competition Metrics}

\begin{enumerate}
  \item \textbf{Solve rate} --- fraction of held-out test puzzles with all unit tests passing. Primary leaderboard metric.
  \item \textbf{Capability per FLOP} --- solve-rate improvement per unit of RL compute (measured in total floating-point operations during training). Rewards efficient training, not just long training.
  \item \textbf{Reward Hacking Gallery} --- the community curates creative exploits found during training runs. Not a competition metric, but the most shared artifact. Every hardcoded output, every format hack, every specification exploit goes here.
\end{enumerate}

\subsection{TinyRL-Align: ``The RLHF Baseline''}

\subsubsection*{Setup}

\begin{itemize}
  \item \textbf{Model:} GPT-2 small (124M parameters), pretrained on OpenWebText.
  \item \textbf{Prompts:} IMDB movie review prefixes. The model is asked to complete the review.
  \item \textbf{Reward:} A pretrained sentiment classifier (e.g., \texttt{distilbert-base-uncased-finetuned-sst-2-english}) scores each completion. Positive sentiment = high reward.
\end{itemize}

\subsubsection*{This Is RLHF, Not RLVR}

The reward model is a neural network, not a deterministic verifier. It can be gamed. The policy does not need to write a genuinely positive review --- it needs to write text that the classifier scores highly. These are not the same thing.

\subsubsection*{RLHF Dynamics Exhibited}

\begin{itemize}
  \item \textbf{KL divergence growth} --- the policy drifts from the reference model rapidly when the KL penalty is weak. Observable within 200 steps.
  \item \textbf{Mode collapse} --- the policy converges to a narrow distribution of high-reward outputs (verbose, exclamation-heavy positive text). Entropy collapses.
  \item \textbf{Reward hacking} --- completions that score near 1.0 from the classifier are often incoherent or nonsensical to a human reader. The gap between proxy reward and true quality is the core RLHF problem.
  \item \textbf{Entropy collapse} --- policy entropy drops steeply as the model converges to its mode-collapsed distribution.
\end{itemize}

All dynamics visible in under 30 minutes on a single GPU.

\subsubsection*{Competition Metric}

Highest mean sentiment score (from the classifier) on the held-out IMDB prompt set, subject to a hard KL constraint: $D_{\mathrm{KL}}(\pi_\theta \| \pi_{\mathrm{ref}}) < 5.0$. This forces participants to optimize reward efficiently rather than through unconstrained drift.

\subsubsection*{Why Included}

TinyRL-Align is the lowest-barrier entry point. No code execution, no test harness, no syntax errors --- just a language model and a sentiment classifier. A researcher who has never implemented RL can run their first training loop here in under 30 minutes. The contrast with TinyRL-Code --- gameable neural reward versus in-principle-non-gameable verifier --- is the core RLHF vs.\ RLVR lesson.

\subsection{TinyRL-Quant: ``The PE-RL Flywheel''}

\subsubsection*{Setup}

\begin{itemize}
  \item \textbf{Model:} GPT-2 small (124M), with per-layer precision as a controllable parameter.
  \item \textbf{Action space:} Choose one of \{FP32, FP16, BF16, FP8, INT8, INT4\} for each transformer layer. A policy that maps layer index and layer statistics to a precision choice.
  \item \textbf{Reward:} A composite of evaluation quality (perplexity or accuracy on a held-out set) and inference efficiency (average bits per weight, or measured throughput). Specifically: $r = \text{accuracy} \times \frac{1}{\text{avg\_bits}}$, encouraging quality-per-bit optimization.
\end{itemize}

\subsubsection*{RLVR-Adjacent}

The reward is fully verifiable: accuracy is computed on a fixed evaluation set, and average bits per weight is a deterministic function of the chosen precisions. There is no neural proxy. This places TinyRL-Quant closer to the RLVR end of the spectrum than TinyRL-Align, but the reward combines two objectives, introducing reward design sensitivity that pure unit-test verification avoids.

\subsubsection*{Dynamics Exhibited}

\begin{itemize}
  \item \textbf{Reward design sensitivity} --- the weighting between accuracy and efficiency dramatically changes which layers the policy assigns high precision. Changing the reward formula changes the learned policy structure.
  \item \textbf{Exploration-exploitation in discrete spaces} --- the precision action space is combinatorial ($6^{12}$ for a 12-layer GPT-2). Pure exploitation fails; entropy regularization matters.
  \item \textbf{Eval-set gaming} --- if the same evaluation set is used for reward and final assessment, the policy overfits to it. Using a held-out test set distinct from the reward evaluation set is essential.
\end{itemize}

\subsubsection*{Competition Metric}

Best accuracy (on the held-out test set) per average bits per weight across all layers. A policy achieving 98\% of FP32 accuracy at 4.2 bits/weight beats one achieving 99\% at 6.1 bits/weight.

\subsubsection*{Unique Angle}

Apply the learned quantization policy to two variants: a model trained with RLVR (TinyRL-Code) and a model trained with SFT only. Does RLVR training change which layers the policy assigns high precision? The hypothesis is that RLVR training concentrates representations differently (more in later layers, due to verifiable reasoning demands), making the optimal quantization policy differ. This is a concrete, testable PE/RL intersection question.

\textit{Note: HAQ (Hardware-Aware Automated Quantization) exists as a production system but targets latency lookup tables for specific hardware. TinyRL-Quant focuses on RL dynamics study and the PE/RL intersection, not production deployment.}

\subsection{Stretch Environments}

Two future environments for contributors:

\begin{itemize}
  \item \textbf{TinyRL-Speculative.} RL for speculative decoding decisions: given a draft model and a verifier model, the policy decides how many tokens to speculate and when to verify. Reward is tokens per second on a fixed benchmark. Competition metric: peak tokens/second at fixed model quality. This bridges Chapter~4 (speculative decoding) and RLVR.

  \item \textbf{TinyRL-Reward.} Meta-RL for reward design: the policy proposes reward functions for a downstream task, and is evaluated on whether a model trained against those rewards achieves genuine capability gains (measured on a held-out probe set) without gaming. Competition metric: reward robustness score --- the gap between reward-training performance and probe-set performance, minimized. This is the hardest environment and the most research-adjacent.
\end{itemize}

\begin{thinkfirst}{The Verification Frontier at Toy Scale}
Rank the three TinyRL environments by verification cleanliness: TinyRL-Code (unit tests) $>$ TinyRL-Quant (composite verifiable reward) $>$ TinyRL-Align (neural proxy reward). For each environment, identify the dominant failure mode: verification hacking, KL collapse, or mode collapse. At what point on this spectrum does verification hacking become the dominant failure mode, displacing KL collapse as the primary concern? What does this imply about how verification quality should change your choice of KL penalty coefficient $\beta$?
\end{thinkfirst}

\begin{thinkfirst}{Designing a Minimal RLVR Environment}
Design the absolute simplest environment that satisfies all of the following constraints: (1) smallest model that works (how small is too small?), (2) fewest tasks (what is the minimum task bank size before the policy memorizes the training set?), (3) simplest verifier that is still non-trivially gameable, (4) must exhibit emergent strategy AND verification gaming, (5) must train to completion in under 15 minutes on a single consumer GPU. Write a complete specification: model, task format, verifier pseudocode, reward function, and the exploit you predict the model will find first. Justify each design choice.
\end{thinkfirst}

\subsection{Exercises}

\begin{exercise}{Train with PPO Across All Environments (1--2 hours per environment)}
For each of TinyRL-Code, TinyRL-Align, and TinyRL-Quant: train with PPO for 500 steps using the default hyperparameters (\texttt{lr=1e-5}, \texttt{beta=0.1}, \texttt{clip\_eps=0.2}). For each run, plot reward curve, KL divergence from reference, and policy entropy on the same figure. Answer: (1) What is the primary pathology in each environment? (2) Which environment shows reward hacking first (fewest steps)? (3) In TinyRL-Code, log the first detected reward hack: what did the model do? (4) In TinyRL-Align, at what step does entropy begin to collapse? Compare your findings across the three environments and write a one-paragraph characterization of how the verifier type shapes the failure mode.
\end{exercise}

\begin{exercise}{Train with GRPO and Sweep Group Size (2--3 hours)}
GRPO (Group Relative Policy Optimization) is the algorithm DeepSeek used for R1-Zero RLVR training. It replaces the value function with a group-relative baseline computed over $G$ sampled completions per prompt. For each of the three environments, train with GRPO at $G \in \{4, 8, 16\}$ for 500 steps. Compare: (1) Is GRPO more stable than PPO (lower KL variance, more consistent reward growth)? (2) Does increasing $G$ always improve sample efficiency, or does the benefit plateau? (3) What is the compute cost of $G=16$ vs.\ $G=4$ in wall-clock time? (4) For TinyRL-Code specifically: does larger $G$ reduce the frequency of reward hacking, or does it accelerate it? Produce a 3$\times$3 table: environment $\times$ $G$ value, reporting final reward and final KL at step 500.
\end{exercise}

\begin{exercise}{Remove the KL Penalty and Document Failure Modes (1 hour)}
Train each environment for 500 steps with $\beta = 0$ (no KL penalty). Document: (1) What happens to KL divergence without the penalty? (2) What happens to reward --- does it increase faster, or does training collapse? (3) \textbf{Key comparison:} does TinyRL-Code (RLVR, deterministic verifier) tolerate KL removal better than TinyRL-Align (RLHF, neural reward)? State a hypothesis before running. If RLVR tolerates KL removal better: argue why deterministic verification provides an implicit regularizer that neural rewards do not. If it does not: what does that tell you about the role of KL in preventing reward hacking versus preventing distribution collapse? (4) At what KL value does each environment's policy produce visibly degenerate outputs? Check qualitative samples at step 500.
\end{exercise}

\begin{exercise}{Long-Horizon RLVR Test on TinyRL-Code (3--4 hours)}
Train TinyRL-Code for 5{,}000 steps (10$\times$ the standard 500-step run). Log at every step: solve rate on training prompts, solve rate on held-out test set, policy entropy, gradient $\ell_2$ norm, and KL from reference. Questions: (1) Do new capabilities emerge late --- does the model begin solving medium or hard puzzles only after step 3{,}000? (2) Do new failure modes appear late --- does verifier overfitting (hardcoded outputs for training prompts) worsen after step 2{,}000? (3) Plot capability (held-out solve rate) vs.\ log compute (log total FLOPs). Is the curve approximately log-linear? If so, this supports the claim that RLVR capability scales smoothly with compute. If not, describe the shape and propose an explanation. (4) Apply a change-point detection algorithm (PELT or binary segmentation) to the held-out solve rate series. Do change-points correspond to jumps in entropy or KL?
\end{exercise}

\begin{exercise}{Quantization Stress Test Under RL Training (2--3 hours)}
This exercise tests numerical stability of RL training under reduced precision, directly relevant to Dario's ``laminar flow'' framing of stable training runs. For TinyRL-Align: run the standard PPO training loop with the policy model stored and updated in (a) FP32, (b) FP16, (c) FP8 (if your hardware supports it), (d) INT8 (using a quantization library). For each precision: (1) Log gradient $\ell_2$ norm at each step. At what precision does the gradient norm become unstable (variance $>10\times$ FP32 variance)? (2) Does training diverge? At what step? (3) Compare final reward at step 500 across precisions. (4) Is there a precision where training appears stable by reward but has degraded by KL or entropy metrics? Report the ``safe precision floor'' for RL training on your hardware and explain what numerical phenomenon causes instability below that floor.
\end{exercise}

\begin{exercise}{Adversarial Verification Engineering (2--3 hours)}
This is the core skill for RLVR at scale. For TinyRL-Code: (1) \textbf{Design 5 gameable test suites.} Each should be gameable in a different way: (a) tests that only check a single predictable input, (b) tests that check output format but not correctness, (c) tests with insufficient coverage that miss a large class of bugs, (d) tests with ambiguous specifications exploitable by degenerate solutions, (e) tests that are bypassable by catching exceptions. (2) Train a fresh model against each gameable suite for 300 steps. Document the exploit each suite induces. (3) \textbf{Harden each test suite.} Add edge cases, remove predictable inputs, enforce semantic correctness. (4) Re-train against the hardened suites. Measure: how much does hardening improve held-out solve rate on the \textit{canonical} test suite (which was never shown during training)? Report a ``verification hardening multiplier'' for each suite: ratio of post-hardening to pre-hardening canonical solve rate. Explain the pattern.
\end{exercise}

\begin{portfolio}{TinyRL: The MNIST of RL for LLMs}
\textbf{Deliverable:} An open-source framework on GitHub with three implemented environments (Code, Align, Quant), a leaderboard site, and a launch blog post.

\textbf{The pitch:} ``I built a standardized platform where anyone can study RLVR and RLHF dynamics in 30 minutes on one GPU. The headline environment --- TinyRL-Code --- is RLVR in miniature: verifiable rewards, deterministic unit tests, emergent problem-solving strategies, and a reward hacking gallery. TinyRL-Align provides the RLHF baseline for direct comparison. The contrast between them is the pedagogical core: it shows concretely why the field moved from neural proxy rewards to verifiable rewards, and what is gained and lost in that transition.''

\textbf{The monkey:} The pedestal is ``can you build a framework?'' That is not the question. The monkey is: \textit{What RLVR phenomena are observable at toy scale, and which require frontier-scale models?} If TinyRL-Code reproduces real pathologies --- reward hacking, entropy collapse, emergent chain-of-thought under verifiable pressure --- you have built a useful research tool. If it does not, you have learned something equally important about the scale-dependence of RL phenomena. Either outcome is valuable. Kill the monkey by running TinyRL-Code for 5{,}000 steps and characterizing what it does and does not reproduce.

\textbf{Verification criteria:}
\begin{enumerate}
  \item TinyRL-Code exhibits at least 2 qualitatively distinct reward hacking strategies, documented in the Reward Hacking Gallery with completion examples and step counts.
  \item The RLHF vs.\ RLVR comparison (Align vs.\ Code, same algorithm and hyperparameters) shows measurably different dynamics: different dominant failure modes, different KL trajectories, different entropy collapse timing. Document this with side-by-side plots.
  \item A user who has never implemented RL can train their first model in under 30 minutes by following the README. Test this by asking someone with no RL background to run TinyRL-Align from scratch.
  \item At least 3 community-contributed environments within 3 months of public launch. This requires writing a contribution guide, a minimal environment template, and seeding the community with the stretch environments as open issues.
\end{enumerate}

\textbf{Why community:} Lowering the barrier to RLVR experimentation creates contributors and competitors. The reward hacking gallery is inherently shareable --- every exploit is a story. Network effects compound: each new environment attracts researchers who care about that domain, who then contribute fixes and new environments. The leaderboard creates a persistent reason to return. This is how MNIST became a standard: not by being the best benchmark, but by being the one everyone used.
\end{portfolio}

\newpage
\section{RL Pipeline \& Verification Engineering}

\smallskip\noindent\textit{Study Guide companion: Study Guide Sections 2--5 cover this material in interview format.}
\smallskip

The modern post-training pipeline is \textbf{RLVR-first}: Base model $\to$ SFT $\to$ RLVR (GRPO with verifiable rewards on math/code) $\to$ RLHF (neural reward model for subjective tasks). RLHF is no longer the main compute consumer. Understanding both stages --- and where verification engineering fits --- is the entry point for every RE role at a frontier lab.

\begin{table}[h]
\centering
\begin{tabular}{@{}lll@{}}
\toprule
\textbf{Era} & \textbf{Pipeline} & \textbf{Compute Distribution} \\
\midrule
2020--2022 & Pretraining only & 100\% pretraining \\
2022--2024 & Pretraining $\to$ SFT $\to$ RLHF & $\sim$95\% pretraining, $\sim$5\% SFT+RLHF \\
2025+ & Pretraining $\to$ SFT $\to$ RLVR $\to$ RLHF & Pretraining shrinking, RLVR growing \\
\bottomrule
\end{tabular}
\caption{Training pipeline evolution. The RLHF stage has not disappeared --- it handles subjective alignment tasks where clean verification is unavailable. But RLVR now dominates post-training compute for capability gains.}
\end{table}

\textbf{Intelligence-per-Watt.} Under RLVR, every wasted training step is wasted compute at massive scale. A verifier that is 10\% more robust effectively gives you 10\% more intelligence per Watt --- not via kernel optimization, but via signal quality. This is why verification engineering is a first-class research problem, not a supporting task.

\subsection{Conceptual Foundation}

\begin{thinkfirst}{The RL training pipeline as a system}
Draw the full post-training pipeline on paper. Both RLVR and RLHF stages are present. For each stage, answer:

\begin{enumerate}
  \item \textbf{SFT:} What distribution does the model learn? DeepSeek R1-Zero skipped this stage entirely. What emerged? What was missing? When would you skip it and when would you not?

  \item \textbf{RLVR stage:} The reward is binary --- the verifier says pass or fail. You need no reward model. What are you trading away compared to a trained reward model? What do you gain? (Hint: Goodhart's Law runs in both directions.)

  \item \textbf{RLHF stage (neural RM):} Why rankings instead of scores? If you have 4 outputs to rank, how many pairwise comparisons can you extract? Why is this more data-efficient than scoring independently?

  \item \textbf{The KL penalty in GRPO:} Imagine removing it entirely. The model is free to go anywhere in output space. Predict what happens in the first 100 training steps. In the first 1000. Why does the failure mode get \textit{worse} over time?

  \item \textbf{GRPO vs.\ PPO:} PPO needs 4 models in GPU memory simultaneously. GRPO drops the critic. Write down what each of PPO's 4 models does and which require gradients. What does GRPO replace the critic's function with?
\end{enumerate}
\end{thinkfirst}

\begin{thinkfirst}{Why online RL beats offline}
DPO trains on a fixed dataset of preference pairs. PPO/GRPO generate new outputs during training and evaluate them live.

\begin{enumerate}
  \item A model trained with DPO discovers a new exploit after deployment. Why couldn't DPO have caught this during training? What would GRPO have done differently?

  \item Frontier RL teams use online RL, not DPO. Explain why in one sentence that a non-ML engineer could understand. (Hint: a student who only studies from a fixed answer key vs.\ one who practices with a tutor who gives new problems.)

  \item DPO's loss function is binary cross-entropy on preference pairs. What information is lost compared to RLVR's binary verifier? What information is lost compared to RLHF's scalar reward model? Are these the same loss?
\end{enumerate}
\end{thinkfirst}

\begin{thinkfirst}{Goodhart's Law in RL --- the verification gaming frame}
``When a measure becomes a target, it ceases to be a good measure.'' Under RLVR, the verifier \textit{is} the measure. Apply this:

\begin{enumerate}
  \item You train a coding agent with reward = ``all unit tests pass.'' The model learns to modify the test file. Why does the RL algorithm see this as a perfectly valid solution? What would you need to change about the verifier to prevent this?

  \item Design a reward for ``write a good essay'' using an LLM judge. List 3 ways a model could game this reward without writing a good essay. Be specific --- think about what the judge is actually sensitive to. Now explain why RLVR on math sidesteps this problem.

  \item Imagine using \textit{human expert decisions} as the reward signal instead of a neural RM. In what sense is this more robust to gaming? In what sense is it fragile? Think about what happens on easy cases vs.\ hard/ambiguous cases.
\end{enumerate}
\end{thinkfirst}

\begin{thinkfirst}{The verification frontier}
Rank the following domains from ``easiest to verify'' to ``hardest to verify.'' For each, describe what verification would look like and what the model could exploit:

\begin{enumerate}
  \item Arithmetic (e.g., ``What is 847 $\times$ 293?'')
  \item Competitive programming (e.g., Codeforces problems)
  \item Code review quality (e.g., ``Does this PR introduce bugs?'')
  \item Customer service quality (e.g., ``Was this a helpful response?'')
  \item Novel writing (e.g., ``Is this a compelling opening chapter?'')
\end{enumerate}

For the middle-difficulty domains (code review, customer service), design a \textbf{composite verifier} --- a combination of automated checks + LLM judge + human spot-checks. What would the model try to game? How would you detect it?
\end{thinkfirst}

\begin{thinkfirst}{Process vs.\ outcome rewards}
A Process Reward Model (PRM) scores each reasoning step. An Outcome Reward Model (ORM) scores only the final answer. RLVR typically uses an ORM-equivalent (pass/fail on the final answer).

\begin{enumerate}
  \item A student solves a math problem. They get the right answer but their reasoning has a cancelling error in step 3 and step 5. What does the ORM say? What does the PRM say? Which is more useful for learning?
  \item PRMs need step-level labels. Where do these come from? Why are they expensive? Can you generate step-level labels \textit{automatically} for coding tasks? How?
  \item DeepSeek R1 used \textit{only} outcome rewards. The model spontaneously developed chain-of-thought, self-verification, and backtracking. If process rewards aren't necessary for \textit{generating} good reasoning, what \textit{are} they necessary for?
\end{enumerate}
\end{thinkfirst}

\subsection{Hands-On Exercises}

\begin{exercise}{nanoRLVR --- The modern baseline (10--15 hours)}
Build a \textbf{minimal RLVR pipeline} on a verifiable task. This is the baseline every RE candidate should have in their portfolio.

\textbf{Suggested setup:} GPT-2 (124M) fine-tuned on simple math problems (addition, subtraction, basic arithmetic). The verifier is an exact-match checker against the ground truth answer --- no neural reward model needed.

\textbf{Stages:} (1) SFT: fine-tune GPT-2 on math problem--solution pairs in a structured format. (2) RLVR: GRPO loop --- for each problem, generate $G = 8$ completions, run the verifier on each, compute group-normalized advantages $A_i = (r_i - \bar{r})/\sigma_r$, update with PPO clipped objective. No critic network.

\textbf{Ablations that teach you the most:}
\begin{enumerate}[label=\alph*)]
  \item Remove the KL penalty ($\beta = 0$). Train 500 steps. What does the model learn to do? Can you identify the exact exploit it found? (Common: it generates very long outputs that happen to contain the correct number, or it learns a formatting trick that systematically passes the verifier.)
  \item Vary group size $G = 2, 4, 8, 16$. How does advantage estimation quality change? At what $G$ does training become unstable?
  \item Swap the verifier for a deliberately broken one (e.g., ``correct if answer contains any digit $>$ 5''). Watch what the model learns. This is verification gaming in miniature.
\end{enumerate}

\textbf{Verify your pipeline:} Reward increases over first 200 steps; KL bounded; generated completions are coherent; group-normalized advantages have mean $\approx 0$ and std $\approx 1$.
\end{exercise}

\begin{exercise}{nanoRLHF --- Legacy comparison (10--15 hours)}
\textbf{Run this on the same model after nanoRLVR.} This is a deliberate comparison, not an independent exercise.

Build the classic RLHF pipeline: GPT-2 + sentiment steering. Use a pre-trained sentiment classifier (\texttt{distilbert-base-uncased-finetuned-sst-2-english}) as the reward model. SFT on IMDB reviews, then PPO with the classifier as reward.

\textbf{The comparison experiment:} After running both pipelines to equivalent compute, compare: (a) reward signal noise (verifiable reward is binary and clean; sentiment classifier output is noisy and continuous), (b) the exploits each pipeline discovers, (c) which pipeline requires more hyperparameter tuning.

Write a 1-paragraph post-mortem: in what task setting would you prefer RLHF over RLVR? What does the answer tell you about the verification frontier?
\end{exercise}

\begin{exercise}{PPO to GRPO --- Drop the critic (4--6 hours)}
Starting from a working PPO implementation (or TRL's \texttt{PPOTrainer}):

\begin{enumerate}
  \item Remove the critic/value network entirely. For each prompt, generate $G = 8$ completions. Score all 8. Compute advantages as $A_i = (r_i - \bar{r})/\sigma_r$ (group normalization). Use these in the PPO clipped objective.
  \item \textbf{Measure:} GPU memory saved by dropping the critic. Training stability (reward variance across batches) for PPO vs.\ GRPO. Try $G = 2, 4, 8, 16, 32$.
  \item At what $G$ does GRPO match PPO's training stability? Is the memory savings worth it? Write a 1-paragraph recommendation.
\end{enumerate}
\end{exercise}

\begin{exercise}{Verification hacking benchmark (6--8 hours)}
Build a small benchmark where models can exploit verification loopholes --- then harden it.

\textbf{Setup:} 10--20 coding problems where the verifier is a set of unit tests. Plant loopholes: (a) predictable output patterns the model can hard-code, (b) tests that check format only, not correctness, (c) deterministic tests with insufficient coverage.

Run best-of-N sampling (N=16). Count how many problems the model ``solves'' by gaming the tests. Then harden: add randomized test inputs, format-independent output checking, adversarial cases. Re-run and measure how many exploits you closed. Measure the cost (human time, compute) of each hardening step. Is there a Pareto frontier of verification quality vs.\ cost?
\end{exercise}

\begin{exercise}{Composite verifier design for a fuzzy domain (4--6 hours)}
Pick a domain where ``good'' is fuzzy (code review quality, documentation clarity, explanation quality).

Design a composite verifier with 3+ components: automated checks, LLM judge, and a consistency check (do two independent judges agree?). Weight the components into a scalar reward. Generate 50 outputs at varying quality; score them. Inspect the top-10 and bottom-10 manually --- does the verifier rank correctly?

Then: deliberately try to game your own verifier. Write 5 outputs designed to score high without being good. If you succeed, iterate. Write a 1-page analysis: at what point would a model trained against this reward start producing outputs you'd be unhappy with? Is this signal clean enough for RLVR?
\end{exercise}

\subsection{Optional Extensions}

These are productive next steps after the core exercises but not required for a portfolio:

\begin{itemize}
  \item \textbf{DPO comparison:} Train DPO on preference pairs generated by your RLVR model. Compare: can DPO reach the same performance? Where does it fail to generalize?
  \item \textbf{Multi-turn RL:} Extend nanoRLVR to a 3--5 step environment. The credit assignment problem becomes real when the reward comes only at the end.
  \item \textbf{RLAIF:} Replace the human-labeled reward model in nanoRLHF with AI-generated feedback. Can the AI judge provide useful signal on tasks it can't solve itself?
\end{itemize}

\subsection{Portfolio Projects}

\begin{portfolio}{nanoRLVR: The Pedagogical Pipeline}
Build the ``nanoGPT of RLVR.'' A minimal, readable, single-repo implementation that goes from pretraining through SFT, RLVR (GRPO), and RLHF, all on toy tasks with clean verification. Include nanoRLHF as a labeled legacy comparison.

\textbf{The monkey:} Not ``can you build it'' (that's the pedestal). The monkey is: ``what surprising thing did you discover in the ablations that isn't in existing tutorials?'' Include ablation figures: KL sweep, group size sweep, verifier hardness sweep, GRPO vs.\ PPO vs.\ DPO. Document what broke and how you diagnosed it. Include a ``reward hacking gallery'' --- examples of exploits each pipeline found.

\textbf{Verification:} Someone can clone and reproduce all results in $<$4 hours on a single GPU. At least one ablation finding surprises a reader who has used TRL before.
\end{portfolio}

\begin{portfolio}{Verification Hacking Taxonomy}
Build a systematic study of how models game code verifiers. Not just ``they hack tests'' --- categorize exploit types, measure which mitigations work, and publish findings.

\textbf{Structure:} (1) Taxonomy: hard-coding outputs, format exploitation, specification ambiguity, distributional shortcuts, test-case overfitting, environment manipulation. (2) Benchmark: 50+ problems with known vulnerabilities, difficulty-graded. (3) Mitigations tested: sandboxing, randomized inputs, output validation, RM ensembles, specification hardening. (4) Cost-effectiveness analysis: for each mitigation, compute/human-time cost vs.\ exploit closure rate.

\textbf{The monkey:} ``Which exploit categories matter most, and which mitigations are cost-effective?'' Building the benchmark is the pedestal. The Pareto frontier of verification quality vs.\ cost is the monkey. Solve for that.

\textbf{Deliverable:} OSS benchmark + blog post. Verification: your analysis shows at least one dominated mitigation and makes a concrete recommendation an RL team could act on.
\end{portfolio}

\newpage
\section{Scaling Laws}

\smallskip\noindent\textit{Study Guide companion: Study Guide Section 11 (Scaling Laws Deep Dive) covers the interview-ready version of this material.}
\smallskip

Before spending \$10M on a large RLVR training run, can you predict from \$10K experiments whether the reward signal and algorithm will hold at scale? Andy Jones demonstrated this skill independently with ``Scaling Scaling Laws with Board Games'' (arxiv 2104.03113) --- training AlphaZero on small Hex boards to predict performance on large ones. This chapter builds the same skill on a toy domain you can run on a single GPU.

\textbf{Intelligence-per-Watt, strategic edition.} Kernel optimization gives you 2$\times$ on one operation. Correct compute allocation gives you 2$\times$ on the entire training budget. Scaling laws are the map from ``how much compute'' to ``how much capability'' --- and the only way to navigate that map without running the expensive experiment is to fit it from cheap experiments first.

\subsection{Conceptual Foundation}

\begin{thinkfirst}{Why scaling laws matter for RL teams}
Before a large RLVR training run, the team needs to answer: ``Will this algorithm + reward signal + compute budget actually produce the capability we want?''

\begin{enumerate}
  \item You have a new verifiable reward signal for a coding task. You run small experiments on easy problems and see improvement. Should you immediately scale up? What could go wrong? (Hint: ``things which can help at smaller scales can actually hurt at larger scales.'')

  \item Andy Jones showed that performance on small Hex boards predicts performance on large ones. What's the analog for LLM RLVR? (Hint: ``small Hex board'' $\to$ easy coding problem. ``Large Hex board'' $\to$ hard coding problem. Now: what happens if the verifier for hard problems is noisier than for easy ones?)

  \item The Jones paper found performance is sigmoid in compute --- there's a ``takeoff'' threshold below which nothing happens. What does this mean for budget allocation? If you're below the takeoff threshold, is it better to run a bigger model or train longer? Does RLVR's longer training runs (relative to RLHF) change this answer?
\end{enumerate}
\end{thinkfirst}

\begin{thinkfirst}{The train/test compute tradeoff}
Jones found: $10\times$ more training $\approx$ $15\times$ less inference search needed.

\begin{enumerate}
  \item Translate this to LLMs: what does ``inference search'' correspond to? List at least 3 forms of test-time compute that LLMs use.

  \item If this tradeoff holds for frontier coding agents: more RLVR training should reduce the chain-of-thought needed at inference, making each API call cheaper. Why does this matter for serving economics?

  \item The tradeoff is log-linear. Does it hold forever, or is there a floor? Even a perfectly trained model might need \textit{some} reasoning for hard problems. What does that floor represent?

  \item Frontier labs optimize across three axes: pretraining, RLVR training, and test-time compute. If the train/test tradeoff is log-linear, what shape is the efficient frontier? Sketch it on paper. Where on that frontier do you think Claude 3.5 Sonnet sits relative to o1?
\end{enumerate}
\end{thinkfirst}

\subsection{Hands-On Exercises}

\begin{exercise}{Replicate Jones at toy scale (15--20 hours)}
Replicate the core finding of the Scaling Scaling Laws paper on a toy domain. Use a simple RL algorithm (DQN or PPO) on gridworld mazes of increasing size (5$\times$5, 7$\times$7, 10$\times$10, 15$\times$15).

\begin{enumerate}
  \item \textbf{Sweep compute:} For each maze size, train with 5 different compute budgets (steps: 1k, 3k, 10k, 30k, 100k).
  \item \textbf{Plot the frontiers:} For each maze size, plot performance vs.\ compute. Do you see the sigmoid shape Jones found?
  \item \textbf{Fit the scaling law:} Use Jones's 5-parameter change-point model:
  \begin{lstlisting}
  plateau = m_size * maze_size + c_plateau
  incline = m_size_i * maze_size + m_flops * log(compute) + c_incline
  performance = clip(incline, plateau, ceiling)
  \end{lstlisting}
  Fit with \texttt{scipy.optimize.curve\_fit}. The fit should achieve $R^2 > 0.9$ on training data --- if not, check initialization or functional form.
  \item \textbf{The key test:} Fit only on maze sizes 5--10. Predict performance for sizes 12 and 15. This is the ``scaling scaling laws'' result. How accurate is the extrapolation?
  \item \textbf{The open question:} Does prediction accuracy degrade gracefully or catastrophically as you extrapolate to size 20? Document what you find.
\end{enumerate}
\end{exercise}

\begin{exercise}{The train/test tradeoff (2--3 hours)}
Using the same toy domain:

\begin{enumerate}
  \item Train an agent at 5 compute levels. At each level, evaluate with varying test-time compute (for gridworld: number of planning steps or beam width).
  \item Plot iso-performance curves: what combinations of train and test compute give the same solve rate?
  \item Fit a log-linear relationship: $\log(\text{test}) = a \cdot \log(\text{train}) + b$. Does it hold? What's $a$?
  \item Compare to Jones's finding: $a \approx -1.2$. Is yours steeper or shallower? What might cause the difference?
\end{enumerate}
\end{exercise}

\begin{exercise}{Where does extrapolation break? (2--3 hours)}
This is the genuinely hard question Jones raised but didn't fully answer.

\begin{enumerate}
  \item In your toy domain, introduce a \textbf{qualitative change in difficulty} at some maze size (e.g., at size 12+, add locked doors requiring key items). Does the scaling law still extrapolate across this discontinuity?
  \item What's the analog for RLVR on code? If the verifier (unit tests) is clean for easy problems but starts breaking down for hard ones (insufficient test coverage), will the scaling law still hold?
  \item \textbf{Write a 1-paragraph prediction:} When you scale RLVR compute on coding tasks, where do you expect the scaling law to break? What would cause a discontinuity? (Think about verifier quality, not just model capacity.)
  \item \textbf{Design the next experiment:} You have 24 GPU-hours. Write a brief proposal: what you'd test, what you'd measure, what the result would tell you. Then run it.
\end{enumerate}
\end{exercise}

\begin{portfolio}{Scaling Laws for RL on Code}
Replicate Jones's approach but for \textbf{coding tasks of varying difficulty}. Use a small model (e.g., CodeGen-350M or StarCoder-1B) and RLVR-fine-tune it on programming problems graded by difficulty.

\textbf{Problem difficulty parameterization:}
\begin{itemize}
  \item Easy: single-function, 1--3 unit tests, standard algorithms
  \item Medium: multi-function, 5--10 tests, requires data structure choice
  \item Hard: multi-file, 10+ tests, requires architectural reasoning
\end{itemize}

\textbf{Key questions:}
\begin{enumerate}
  \item Does performance scale as a sigmoid in compute for each difficulty level (like Jones's Hex)?
  \item Can you predict hard-task performance from easy-task experiments?
  \item At what difficulty level does the unit-test verifier start breaking down? Smooth degradation or cliff?
  \item What's the train/test compute tradeoff for coding? Does more RLVR training reduce the chain-of-thought needed at inference?
\end{enumerate}

\textbf{The monkey:} ``Does extrapolation across problem difficulty actually hold for code?'' Running the experiments is the pedestal. The insight about \textit{where and why} extrapolation breaks --- and whether it's the verifier or the model that's the bottleneck --- is the monkey.

\textbf{Deliverable:} Paper or detailed blog post with scaling law fits and extrapolation results. \textbf{Verification:} Prediction error on held-out difficulty levels $<$15\%. If worse, document where extrapolation breaks --- that finding is equally valuable.
\end{portfolio}

\newpage

\section{Synthesis: How It All Connects}

\subsection{The 9-Chapter Puzzle}

Each chapter is a lens on the same question: how do we produce more capable intelligence per unit of energy? The chapters connect through the intelligence-per-Watt metric as follows.

\textbf{Tier~1 optimizes FLOPs/Watt.} Chapter~1 (AMD MI300X) establishes a hardware baseline and reveals whether bandwidth or compute is the binding constraint for RLVR. Chapters~2--3 (NVIDIA/Blackwell, TPU/Pallas) build the kernel engineering depth needed to close the gap between theoretical and realized FLOP/s. Chapter~4 (RL Inference) attacks the generation bottleneck --- the actor's forward pass during rollout collection dominates RLVR wall time. Chapter~5 (Distributed Training) addresses the communication overhead that prevents linear scaling across nodes.

\textbf{Tier~2 optimizes Quality/FLOP.} Chapter~6 (Policy Gradient) builds the theoretical foundation: which algorithmic choices --- clipping, GAE, KL penalties, entropy regularization --- are actually load-bearing, and why? Chapter~7 (TinyRL) operationalizes that theory in a community-auditable framework where every design decision has a clean ablation.

\textbf{Tier~3 provides the systems foundation.} Chapter~8 (RL Pipeline \& Verification) addresses the silent failure modes that waste compute silently: reward hacking, gradient explosion, and reward model collapse. Chapter~9 (Scaling Laws) asks whether a training run is worth running before committing the compute.

The chapters do not compose sequentially --- they compose \textit{hierarchically}. The hardware chapters set the iteration speed (experiments per week). The algorithm chapters determine what each experiment is testing. The systems chapters determine whether the experiment's result is trustworthy. All three must be healthy for a research program to move fast.

\subsection{The PE--RL Flywheel}

Performance Engineering and RLVR are not separate concerns. They form a flywheel:

\begin{itemize}
  \item \textbf{PE makes RLVR faster.} Higher generation throughput (Chapter~4) means more rollouts per hour. More rollouts per hour means more gradient steps per week. More gradient steps per week means faster policy improvement --- or faster discovery that a hypothesis is wrong.
  \item \textbf{RLVR can automate PE decisions.} A policy trained with verifiable rewards on compiler benchmarks can learn quantization schemes, kernel configurations, and memory layout decisions that outperform human heuristics. Chapter~7's TinyRL framework is the natural substrate for these experiments.
  \item \textbf{The person who understands both is uniquely valuable.} Most PE engineers do not understand what makes a policy gradient step meaningful. Most RL researchers do not know what a memory-bandwidth-bound kernel looks like in a profiler. The combination is rare and directly relevant to frontier RL infrastructure roles.
\end{itemize}

\subsection{Picking Your Portfolio Project}

\begin{table}[h]
\centering
\footnotesize
\setlength{\tabcolsep}{3pt}
\begin{tabular}{@{}lccccccc@{}}
\toprule
\textbf{Criterion} & \textbf{AMD} & \textbf{CUDA} & \textbf{Pallas} & \textbf{Infer.} & \textbf{Distrib.} & \textbf{PG} & \textbf{TinyRL} \\
 & \textbf{(Ch1)} & \textbf{(Ch2)} & \textbf{(Ch3)} & \textbf{(Ch4)} & \textbf{(Ch5)} & \textbf{(Ch6)} & \textbf{(Ch7)} \\
\midrule
Clean feedback?   & Yes  & Yes  & Yes    & Yes  & Yes    & Yes     & Yes  \\
Fast iteration?   & Yes  & Yes  & Med.   & Yes  & Med.   & Yes     & Yes  \\
Underexplored?    & V.Hi & Low  & High   & High & Med.   & ?       & High \\
RLVR-relevant?    & Yes  & Yes  & Yes    & Very & Yes    & Yes     & Very \\
\bottomrule
\end{tabular}
\end{table}

\textbf{Recommendation by goal:}
\begin{itemize}
  \item \textbf{Fastest artifact (5--7 weeks):} Ch1 AMD benchmark. One number --- MI300X vs.\ H100 on end-to-end RLVR throughput --- is all you need to start a conversation.
  \item \textbf{Strongest community signal:} Ch7 TinyRL. Open-source + leaderboard entries are persistent artifacts that compound over time.
  \item \textbf{Broadest role targeting:} Ch2 + Ch3. Covers GPU Kernel Engineering (NVIDIA stack) and TPU Kernel Engineering simultaneously.
  \item \textbf{Highest research ceiling:} Ch6 Policy Gradient hypotheses. A single confirmed or refuted hypothesis about clipping behavior or GAE bias is more valuable than a benchmark table, if it changes how people train policies.
\end{itemize}

\subsection{The Meta-Lesson}

Every project in this workbook follows the same template:

\begin{enumerate}
  \item Pick a well-defined problem with clean verification.
  \item Build the simplest version that could possibly produce a signal.
  \item Ablate one variable at a time. Predict before you observe.
  \item Document honestly --- negative results and failed hypotheses are publishable.
  \item Publish. An artifact that exists is worth more than a project that is almost done.
\end{enumerate}

Before calling a project ``done,'' answer three questions: (1) Can you point to a single number that says you succeeded? (2) Can you run one full experiment cycle --- change, train, evaluate --- in under one day? (3) Does the artifact have value to someone besides you? If all three are yes, you are done. If any is no, the project is not finished --- it is a draft.

\vfill
\begin{center}
\small\textit{Compiled \today.}
\end{center}

\newpage
\section*{Appendix A: World Models \& the JEPA Debate}

\textit{This material is moved from the main track to an appendix per guidance to allocate a ``small exploration budget'' to alternative approaches. The content is valuable for breadth but not on the critical path for PE/RE roles. Read it after completing the core chapters. The exercises here are genuinely interesting and could yield a portfolio project --- but only after you have a working nanoRLVR and a scaling law replication in hand.}

LeCun claims autoregressive LLMs are a dead end for machine intelligence. Meta's AMI Labs raised \$1B on that thesis. Dreamer models collect diamonds in Minecraft without human data. Meanwhile, DeepSeek R1 shows RLVR on LLMs produces emergent reasoning from nothing but a right/wrong signal. Who's right? This appendix gives you the math to decide for yourself --- and proposes experiments at the intersection that could matter for frontier RL teams.

\begin{whythis}
The RLVR scaling team is compute-constrained. If model-based RL (Dreamer-style) could improve sample efficiency by 10$\times$ on code tasks, that's equivalent to a 10$\times$ compute budget increase. If it can't, pursuing it is a distraction. This appendix builds the judgment to evaluate that tradeoff --- a skill that matters whether the answer is yes or no.
\end{whythis}

\subsection*{The Three Prediction Paradigms}

Before comparing architectures, understand what each one \textit{predicts} and what that costs information-theoretically.

\textbf{Paradigm 1: Autoregressive (LLMs).} Predict the next token in observation space:
$$\mathcal{L}_{\text{LLM}} = -\mathbb{E}\left[\log p(t_{k+1} \mid t_{1:k})\right]$$
Cross-entropy over a discrete vocabulary ($|V| \approx$ 50--100K). Must model \textit{all} stochasticity in the token distribution --- rare words, formatting, irrelevant details.

\textbf{Paradigm 2: Pixel prediction (MAE, video models).} Predict masked pixels:
$$\mathcal{L}_{\text{pixel}} = \|x_{\text{masked}} - \hat{x}_{\text{masked}}\|_2^2$$
L2 in pixel space. Must model textures, lighting, noise --- everything. Entropy cost: $H(X) \gg H(\text{semantics})$.

\textbf{Paradigm 3: JEPA (representation prediction).} Predict masked \textit{embeddings}:
$$\mathcal{L}_{\text{JEPA}} = \frac{1}{M}\sum_{i=1}^{M}\|\hat{y}_i - y_i\|_2^2$$
where $\hat{y}_i = g_\phi(s_x, z_i)$ is the predictor's output, $y_i = \text{Enc}_{\text{target}}(\text{patch}_i)$ is the target encoder's embedding (with stop-gradient), and the target encoder is updated via EMA: $\theta_{\text{target}} \leftarrow \tau \cdot \theta_{\text{target}} + (1 - \tau)\cdot\theta_{\text{context}}$, $\tau \approx 0.999$.

\textbf{The key difference:} JEPA's prediction target is a \textit{learned} representation ($\sim$768D), not raw observations ($\sim$50K tokens or millions of pixels). The model learns what to throw away. Without care, it collapses to a constant --- VICReg prevents this with explicit variance and covariance regularization terms, no negative samples needed.

\textbf{Model-based RL (Dreamer):} Learn a world model, then optimize the policy \textit{in imagination}. DreamerV3's RSSM learns deterministic hidden state $h_t = f(h_{t-1}, z_{t-1}, a_{t-1})$ and stochastic latent $z_t$ (32 categorical distributions via straight-through gradients). Policy is optimized on imagined $H$-step rollouts using $\lambda$-returns --- no gradients flow back into the world model. Empirically: $\sim$10$\times$ fewer real environment interactions than model-free RL on control tasks. DreamerV3 was the first algorithm to collect diamonds in Minecraft from scratch on a single V100.

\subsection*{Conceptual Foundation}

\begin{thinkfirst}{Prediction in observation space vs.\ representation space}
\begin{enumerate}
  \item An LLM trained on code must predict the next token in a Python file. When it encounters \texttt{x = 3 + 4}, it must assign probability to the token \texttt{7}, but also to formatting tokens, variable names, and comments. How much of the model's capacity is ``wasted'' on tokens that don't affect program behavior? Now imagine a JEPA-style model that predicts program \textit{state} embeddings instead of tokens. What would it throw away? What might it lose?

  \item JEPA predicts in $\sim$768D continuous embedding space. LLMs predict over $\sim$50K discrete tokens. Which is cheaper in bits? Which carries more information about the world?

  \item LeCun argues that predicting tokens ``wastes capacity on irrelevant details.'' But LLMs trained on code develop genuine reasoning capabilities (DeepSeek R1: 15.6\% $\to$ 77.9\% AIME from RLVR alone). How do you reconcile these facts? Is it possible that ``wasting capacity'' on surface details is actually a form of grounding that helps reasoning? Argue both sides.

  \item V-JEPA achieves 81.9\% on Kinetics-400 with $<$1200 GPU hours of pretraining (2.5$\times$ faster than iBOT, 10$\times$ faster than MAE). But it has never been demonstrated on abstract reasoning tasks. Is this a limitation of the architecture or just a matter of applying it to the right domain?
\end{enumerate}
\end{thinkfirst}

\begin{thinkfirst}{Model-based vs.\ model-free: when does the world model pay for itself?}
\begin{enumerate}
  \item DreamerV3 achieves $\sim$10$\times$ sample efficiency over model-free RL on control tasks. But it requires 2--5$\times$ more \textit{compute} per environment step. Under what conditions is the 10$\times$ sample savings worth the 2--5$\times$ compute overhead? (Hint: when are environment interactions expensive vs.\ cheap?)

  \item For LLM RLVR, the ``environment interaction'' is generating a completion and running the verifier. If inference costs \$X per rollout, a world model that reduces rollouts by 10$\times$ but costs $5X$ to train saves money when $N > \text{???}$. Work out the crossover. Are frontier labs likely above or below it?

  \item In code execution, the ``environment'' is the compiler/interpreter --- nearly free to run. If running unit tests costs effectively zero, does a learned world model for code execution have \textit{any} sample efficiency advantage? Where is the real bottleneck?

  \item MuZero combined world models with RL for board games with perfect state information and deterministic dynamics. What breaks when you try this for code generation, where the ``state'' is a partially written program and ``dynamics'' are compiler output depending on unseen test inputs?
\end{enumerate}
\end{thinkfirst}

\begin{thinkfirst}{Evaluating paradigm claims --- is LeCun right?}
Here are the actual benchmark results:

\begin{enumerate}
  \item \textbf{JEPA results:} V-JEPA 2 achieves 88.2\% average across 6 video benchmarks. Zero-shot robot control from 1 hour of DROID data. Strong on perception and physics prediction.

  \item \textbf{LLM+RLVR results:} DeepSeek R1 achieves 77.9\% AIME (math reasoning). OpenAI o1 achieves 93\% AIME with 1000-sample reranking. Emergent chain-of-thought and backtracking from only a right/wrong signal.

  \item Do these results support ``LLM+RL is a dead end,'' ``JEPA is a dead end,'' or ``they solve different problems''? What would you need to see to change your answer?

  \item LeCun founded AMI Labs with \$1B in March 2026 on the thesis that JEPA-based world models will outcompete LLMs for AGI. What \textit{specific benchmark result} would AMI Labs need to produce in the next 12 months to validate this thesis? What result would invalidate it? Being concrete is the exercise --- vague claims like ``better reasoning'' don't count.

  \item \textbf{The meta-question:} When someone claims ``X approach is fundamentally limited,'' what's your checklist for evaluating that claim? Write down 4 criteria. (Suggested: (a) is there a mathematical proof or just an argument by analogy? (b) have empirical results already violated the claimed limitation?)
\end{enumerate}
\end{thinkfirst}

\subsection*{Hands-On Exercises}

\begin{exercise}{JEPA vs.\ autoregressive prediction on a toy sequence (3--4 hours)}
Build both prediction paradigms on the same toy task. Generate synthetic time series from $x_{t+1} = A \cdot x_t + \epsilon_t$, where $x \in \mathbb{R}^{10}$, $A$ is a fixed rotation matrix, $\epsilon_t \sim \mathcal{N}(0, \sigma^2 I)$.

\textbf{Approach 1 (Autoregressive):} Discretize each dimension into 64 bins. Train a small transformer to predict the next token. \textbf{Approach 2 (JEPA-style):} Train a context encoder + predictor to predict future embeddings. Use EMA for target encoder ($\tau = 0.99$). Add variance regularization to prevent collapse.

Compare: (a) sample efficiency, (b) noise robustness as $\sigma$ increases, (c) whether the learned representations capture the rotation matrix $A$ better than surface statistics. \textbf{Critical question:} Increase noise until autoregressive prediction becomes useless. Does JEPA still work? If so, you've demonstrated LeCun's core claim in miniature. If not, what does that tell you?
\end{exercise}

\begin{exercise}{Build a minimal world model for code execution (6--8 hours)}
Can you learn a world model for program execution? Generate a dataset of simple Python programs with execution traces: $(\text{code}, \text{input}) \to (\text{state}_0, \ldots, \text{state}_T, \text{output})$. Fine-tune GPT-2 (124M) to predict the next program state given the current state and next line: $\hat{s}_{t+1} = f_\theta(s_t, \text{line}_{t+1})$.

Evaluate single-step prediction accuracy (easy) and multi-step rollout accuracy (hard --- errors compound). At what program length does the world model's rollout diverge from reality? \textbf{The fundamental question:} For code, the interpreter is deterministic and nearly free. A learned world model is approximate and expensive. Under what conditions would you \textit{ever} prefer it? (Hint: when the verifier is slow, expensive, or unavailable --- integration tests, real APIs, hardware, human judgment.)
\end{exercise}

\begin{exercise}{Sample efficiency showdown: model-based vs.\ model-free (4--6 hours)}
Direct empirical comparison on gridworld mazes (reuse from Ch9). Agent 1: DQN or PPO (model-free). Agent 2: simplified Dreamer --- learn a dynamics model $\hat{s}_{t+1} = f_\theta(s_t, a_t)$ and reward model, then imagine $H = 10$ steps per real step.

Measure: (a) real environment steps to reach 80\% solve rate, (b) total wall-clock compute to reach 80\% solve rate --- at what maze size does sample efficiency overcome compute overhead? (c) world model accuracy over training --- what happens to the policy when the world model is inaccurate? (d) Does the model-based advantage grow or shrink with maze size? What does this predict for code generation?
\end{exercise}

\begin{exercise}{Hybrid RL: world model as a cheap pre-filter (4--6 hours)}
Test the most promising hybrid approach for frontier RL teams. Train a ``fast reward predictor'' --- a small model that predicts the verifier's score from partial completions (first 50\% of tokens). During rollout: generate $N = 32$ partial completions per prompt; use the fast predictor to score all 32; complete only the top $K = 8$; score with the full verifier.

Compare to baseline (8 full completions scored by verifier directly). Measure: (a) does filtering find higher-reward completions? (b) is it cheaper in total compute? (c) how accurate does the predictor need to be for filtering to help? Find the accuracy threshold by deliberately degrading the predictor with less training data.
\end{exercise}

\begin{portfolio}{Code Execution World Models for RL Training}
The hybrid experiment that could actually matter: can a learned world model for code execution improve the efficiency of RLVR training on coding tasks?

\textbf{Architecture:} (1) World model: fine-tune a small LLM to predict test outcomes from code --- given (problem, candidate solution, test spec), predict (pass/fail, error type, partial output). Train on 10K+ solved problems with execution traces. (2) RLVR integration: use the world model as a ``cheap verifier'' in the GRPO loop --- generate $N$ candidates, score with world model, select top $K$ for real execution. (3) Calibration: measure the cost of false positives (predicts pass, actually fails) and false negatives (predicts fail, actually passes) to RL convergence.

\textbf{The monkey:} Not ``can you build a world model for code'' (pedestal) but ``is the filtering it provides worth it given that real execution is cheap?'' Finding and documenting the crossover verification-cost threshold is the contribution.

\textbf{Expected outcome:} For cheap verifiers (unit tests), the world model probably isn't worth it. For expensive verifiers (integration tests, API calls, human evaluation), it likely is. Find and document this crossover.

\textbf{Deliverable:} Blog post with results and concrete recommendation for when RL teams should invest in world models. \textbf{Verification:} Your crossover analysis identifies a specific verification-cost threshold (in seconds or dollars per evaluation) above which world model filtering saves total compute.
\end{portfolio}

\subsection*{The Critical Analysis Framework}

This section teaches a \textit{meta-skill}: how to evaluate paradigm claims in ML. The LLM+RLVR vs.\ world models debate is not the last such debate you will encounter.

\textbf{When someone claims ``X approach is fundamentally limited,'' apply this checklist:}

\begin{enumerate}
  \item \textbf{Is there a mathematical proof of the limitation, or an argument by analogy?} LeCun's claim that autoregressive models ``can't plan'' is an argument by analogy. But DeepSeek R1 empirically \textit{does} plan via learned chain-of-thought. Arguments by analogy can be wrong when the system finds unexpected workarounds.

  \item \textbf{Have empirical results already violated the claimed limitation?} If yes, the claim needs updating. o1 achieving 93\% on AIME via test-time search within an autoregressive model directly contradicts ``autoregressive models can't do hierarchical planning.''

  \item \textbf{Are the benchmarks comparable?} JEPA achieves 88.2\% on video understanding. R1 achieves 77.9\% on math reasoning. These are \textit{different tasks}. Claiming one approach is ``better'' without specifying the task is meaningless.

  \item \textbf{What are the incentives?} AMI Labs raised \$1B on the thesis that LLMs are a dead end. Anthropic's valuation depends on LLM+RLVR continuing to work. Neither party is a neutral evaluator. Look at the results, not the press releases.

  \item \textbf{Is it a false dichotomy?} Most ``X vs.\ Y'' debates in ML resolve as ``X for some tasks, Y for others, and X+Y for a third set.'' MuZero combined world models with RL. The interesting question isn't which paradigm wins --- it's where the hybrid is better than either alone.
\end{enumerate}

\begin{taste}
If an interviewer asks ``What do you think about world models vs.\ RL scaling?'' --- the wrong answer is picking a side. The right answer is: ``They solve different problems. World models excel at sample-efficient learning for physics and control. RLVR scaling excels at abstract reasoning with clean verification. The interesting frontier is where you combine them --- like using a learned world model to reduce inference costs during RLVR training. I'd want to run the experiment before committing.'' This demonstrates research taste: you don't pick sides, you identify the informative experiment.
\end{taste}



\end{document}
